% \documentclass{article}
  % \usepackage{maa-monthly}
 \documentclass{amsart}
% \documentclass{siamltex}


  \let\oldchapter=\chapter
\def\resetpage{\setcounter{page}{1}}
\def\chapter{\expandafter\resetpage\oldchapter}  
 
\usepackage{mathdots}

\usepackage[all]{xy}

\usepackage{bm}
\usepackage{bigints}
\usepackage{epsfig}
\usepackage{epstopdf}

\usepackage{srcltx,amssymb}

\usepackage{graphicx}

    \usepackage{amsthm}
    \raggedbottom
      \usepackage{amsmath}
    %% IF YOU HAVE FONTS INSTALLED
    %\usepackage{mtpro2}
    \usepackage{amscd}
    %\usepackage{mathtime}\bullet
    \usepackage[title]{appendix}
    \usepackage{slashed}
    \numberwithin{equation}{section} 
       
   %\theoremstyle{theorem}
   \newtheorem{theorem}{Theorem}
    \newtheorem*{thm2}{Proposition A}
    \newtheorem*{thm5}{Proposition A$'$}
    \newtheorem*{thm3}{Proposition B}
    \newtheorem*{thm4}{Proposition B$'$}
  %  \newtheorem{corollary}[theorem]{Corollary}

\newtheorem{remark}[theorem]{Remark}

\newtheorem*{conjecture}{Conjecture}

   \usepackage{appendix}
   
 \newtheorem{lemma}[theorem]{Lemma}
  \newtheorem{corollary}[theorem]{Corollary}
  %  \newtheorem{theorem}{Theorem}
     \newtheorem{proposition}[theorem]{Proposition}
     
    \theoremstyle{definition}
   \newtheorem{definition}[theorem]{Definition}
  
    
    
   \numberwithin{theorem}{section} 
    
  %  \newtheorem{definition}[theorem]{Definition}
  
  \setcounter{tocdepth}{1}
    
    \begin{document}
       
             \title{Sorting algebras using uncurling invariants}
    \author{Fred Greensite}
    
   \begin{abstract} The multiplicative inversion operation of a real finite-dimensional unital associative algebra defines a vector field and (assuming the standard topology on the vector space of elements) an associated infinitesimal rotation tensor (and in this regard the multiplicative inversion operation is representative of the algebra's implied nonlinear unary operations).  The group of metrics (symmetric linear transformations) that uncurl this vector field, and the group of resulting integrals, embody various algebra isomorphism invariants that can be used to sort algebras in lieu of an unfeasible algebra classification.  Such sorting is of interest as it reflects an inherent structure generalizing the usual algebra norm in an entity additionally sensitive to an algebra's implicit rotor aspects, leading to recognition of various subtle structural features that isomorphic algebras must share, and consequential combinatorial questions regarding the distribution of these features among real algebras of a given dimension.  \end{abstract}
        
  \maketitle
 % \markright{Novel isomorphism invariants of real algebras} 
 \markboth{F. Greensite }{Sorting algebras using uncurling invariants}
  
 \tableofcontents
 
 \section{Introduction} The general goal of classifying finite-dimensional associative algebras with respect to isomorphism equivalence is deemed to be hopeless according to its recognition as a so-called wild problem, e.g., \cite{eom:2024,drozd:1972,belitskii:2005}.   This inherent incompleteness of structural analysis by no means precludes, but rather heightens, the importance of identifying hidden structures underlying or implied by an algebra.  We may consider that such structures are significant if their invariants allow a useful ``sorting" of algebras (or otherwise interesting descriptions) in lieu of a full classification.   Along these lines, we will present a structural entity, norm-like in character, providing novel insights into a real algebra's implicit rotor aspects. \enlargethispage{.5cm} In principle, geometry, analysis, groups, and combinatorics, all play a role in this ultimately functorial program.
 
\subsection{A short guide to the topics in this paper} 
Concisely, we investigate ``anti-rotors" -  novel structures underlying an algebra that are associated with a trove of invariants.
Given a real unital associative algebra $A$ of nonzero finite dimension $n$ where the vector space of elements $\mathbb{R}^n$ is assumed to have the standard topology, an \mbox{anti-rotor} is the additive group of symmetric linear transformations that annihilate the infinitesimal rotation tensor associated with a chosen nonlinear unary operation implied by the axioms of the algebra.  

Since we are dealing with Euclidean $\mathbb{R}^n$, for computational purposes we can employ the standard basis in the context of the usual dot product `` $\cdot$ ", and for $s\in A$ we write $s=(x_1,\dots,x_n)$ with that understanding.  For reasons indicated below, we will also write a differential 1-form as a formal dot product of an ordered $n$-tuple of vector field components and the ordered $n$-tuple ${\sf d}s \equiv (dx_1,\dots, dx_n)$. Thus, for choice of the multiplicative inversion operation as defined by $s\mapsto s^{-1}$ in a suitable neighborhood of the algebra's multiplicative identity $\mathbf{1}_A$, the 1-form dual to the implied vector field is written as the formal expression $s^{-1}\cdot {\sf d}s$.  The infinitesimal rotation matrix is then expressed by $d(s^{-1}\cdot {\sf d}s)$, and the rotor (also known as the curl) is its Hodge dual. 
The anti-rotor specific to this choice of unary operation (i.e., $s\mapsto s^{-1}$ as above) is then,
\begin{equation}\label{8/30/24.3} \mathbf{u}_A \equiv \{L\in M_n(\mathbb{R}): L^T=L \text{ and } d(\left[L(s^{-1})\right]\cdot {\sf d}s)  = 0\},\end{equation}
where superscript $^T$ indicates transpose or dual vector as applicable.  The \mbox{1-form} $\left[L(s^{-1})\right]\cdot {\sf d}s$ can be alternatively written as another formal expression, $(s^{-1})^TL{\sf d}s$, with both $s$ and ${\sf d}s$ understood to be column arrays. 
We have introduced the \mbox{$n$-tuple} of differentials ${\sf d}s$ (with \mbox{sans serif {\sf d}}) because $(s^{-1})^TL{\sf d}s$ is a very convenient format that streamlines \mbox{various proofs}.  

There are good reasons to suspect that the multiplicative inversion operation sufficiently represents the above set of an algebra's implied nonlinear unary operations vis-a-vis invariants derivable from anti-rotors, though presently this must be formulated as a conjecture.  If the conjecture is false, the purview of the analysis expands accordingly.  For now, we take ``anti-rotor" to be defined by (\ref{8/30/24.3}).  Also, although use of the term ``algebra" in the following means a member of the class of algebras described at the beginning of the second sentence in the first paragraph of this subsection, some generalization of our results is possible.


As we have noted, an anti-rotor is an abelian group under matrix addition.  More specifically, it is a vector space having (it turns out) nonzero finite dimension.  However, that description does not adequately capture the structure of this entity, because the matrix product is intrinsic to the important mapping of anti-rotors to groups of ``energy density" integrals, \mbox{$\mathbf{u}_A \mapsto \int_{{\bf 1}_A}^s [\mathbf{u}_At^{-1}]\cdot {\sf d}t$} (where the notation $\mathbf{u}_A t^{-1}$ means $\{L(t^{-1}): L\in \mathbf{u}_A\}$), and the matrix product is also an intrinsic part of a functor relevant to analysis of \mbox{algebra} endomorphisms via anti-rotors.  Accordingly, anti-rotors are better thought of as objects in a particular category comprised of symmetric matrices with entries that are linear combinations of the components of a $k$-tuple of real parameters, where the morphisms are defined utilizing the matrix product.  This reflects the easily demonstrated feature that $\mathbf{u}_A$ can alternatively be cast as a symmetric \mbox{$(n\times n)$-matrix} $M_{\mathbf{u}_A}$ whose entries are linear combinations of the components of an $m$-tuple of real parameters, $(\alpha_1,\dots,\alpha_m)$, where $m$ is the dimension of $\mathbf{u}_A$ as a vector space - so that the members of $\mathbf{u}_A$ are given by $M_{\mathbf{u}_A}$ when its $m$ parameters are replaced by any set of $m$ real values.

The foundation of our program now consists of two results easy enough to be almost called observations. \enlargethispage{1cm} The surprise is that these lead to recognition of \mbox{subtle structural} features that isomorphic algebras must share, and (in both senses) consequential combinatorial questions. 

\begin{enumerate} 
\item   The first foundational result concerns the elegant manner in which \mbox{anti-rotors} transform.   To wit, if $K:A_1\rightarrow A_2$ is an algebra isomorphism, then \begin{equation}\label{9/6/24.1} \mathbf{u}_{A_1} = K^T \mathbf{u}_{A_2} K. \end{equation} 
Among other things, this allows us to extract numerous algebra isomorphism invariants from the anti-rotor.    Some invariants that can be recognized immediately are, \begin{enumerate} \item The dimension $m$ of $\mathbf{u}_A$.   \item The largest rank  of a member of $\mathbf{u}_A$.  \item The smallest rank of a nonzero member of $\mathbf{u}_A$. \item The dimension of the subspace of the algebra to which ${\rm det}(M_{\mathbf{u}_A})$ is sensitive. \item Invariants of the real algebraic set, \[V_\mathbb{R}({\rm det}(M_{\mathbf{u}_A})) = \{(\alpha_1,\dots,\alpha_m)\in\mathbb{R}^m: {\rm det}(M_{\mathbf{u}_A}) = 0 \}.\] \end{enumerate}
  \item The second foundational result (a consequence of (\ref{9/6/24.1})) leads to a complementary set of algebra isomorphism invariants.  Again let $K:A_1\rightarrow A_2$ be an algebra isomorphism, where for $s_1\in A_1$ we define $s_2\equiv Ks_1$.  Then for $s_1$ in an open ball centered at $\mathbf{1}_{A_1}$ consisting only of units, we have \mbox{the gratifying}, \begin{equation}\label{8/30/24.1} \int_{{\bf 1}_{A_1}}^{s_1} \left[\mathbf{u}_{A_1}t_1^{-1}\right]\cdot {\sf d}t_1 
= \int_{{\bf 1}_{A_2}}^{s_2} \left[\mathbf{u}_{A_2}t_2^{-1}\right]\cdot {\sf d}t_2 . \end{equation} There are now invariants arising from enumeration of the results of integration of the above rational function integrands, i.e., concerning particular sums of rational functions, logarithms, and arctangents.  The three isomorphism classes of two-dimensional algebras each distinctly embody in pure elemental form the three invariants stemming from the above integrals.
\item The above invariants ``sort" the algebras in potentially interesting ways, and there are many combinatorial questions related to the invariants.  The following are some examples relevant to the set of $n$-dimensional algebras as regards the above invariants immediately arising from (\ref{9/6/24.1}). \begin{enumerate} \item How are the algebras distributed between the classes defined by the anti-rotor's dimension (e.g., what are the relative popularities of these classes)?  \item For algebras having the same anti-rotor dimension, how are they further distributed into classes defined by different numbers of irreducible components of $V_\mathbb{R}({\rm det}(M_{\mathbf{u}_A}))$?  \item For which algebras is ${\rm det}(M_{\mathbf{u}_A})$ identically $0$?   \item For those algebras having the same anti-rotor dimension $m$, and for each $j\le m$, how likely is $j$ to be the maximal rank of a member of $\mathbf{u}$?   \end{enumerate}  
Additional combinatorial questions arise from the invariants associated \mbox{with (\ref{8/30/24.1})}. 
\item  The ``usual" algebra norm (the determinant of an element's image under the left regular representation \cite{bourbaki:1989}) is found to be closely associated with a member of the group of ``unital norms" derived from the collection of integrals in (\ref{8/30/24.1}).  This leads to a broadened understanding of the notion of an algebra norm.
\item As indicated above, we are led to a conjecture that, invariant-wise, it is sufficient to consider only the anti-rotors arising from the multiplicative inversion operation (i.e., (\ref{8/30/24.3})).  The truth or falsity of this conjecture greatly impacts the complexity and scope of this program. 
\item Through the functor alluded to above, there is introduction of a category whose objects are symmetric matrices composed of entries that are linear combinations of the components of a $k$-tuple of real parameters, and whose morphisms are defined utilizing the matrix product.  An interesting question that arises is:  which objects in this category are anti-rotors?
\item There may be ample opportunity for generalization of this program to other classes of algebras.  Some nonunital and some not associative algebras are treated in the present paper, but complex and infinite-dimensional algebras are not considered here.
\end{enumerate}

 
\subsection{Conventions}  Unless stated otherwise, ``algebra" means a real unital associative algebra having nonzero finite dimension, i.e., a unital associative algebra whose vector space of elements is $\mathbb{R}^n$. The standard topology on $\mathbb{R}^n$ is assumed, so that for analytic purposes an algebra's vector space of elements is assumed to incorporate the Euclidean norm/inner product.   For computational purposes, the standard basis on $\mathbb{R}^n$ is always assumed, i.e., the inner product as the usual dot product ``\,$\cdot$\," pertains to the elements written as $s=(x_1,\dots,x_n)\in\mathbb{R}^n$.  We will have reason to introduce the column array of differentials ${\sf d}s \equiv \left(\begin{array}{ccc}
dx_1   \\ 
\vdots \\
dx_n   
\end{array}\right)$.  For $v\in\mathbb{R}^n$, and $L\in M_n(\mathbb{R})$, formal expressions like $v\cdot {\sf d}s$ and $L{\sf d}s$ are interpreted as the expected symbol manipulations.   The juxtaposition of two algebra elements indicates their algebra product.  An algebra will typically be denoted by $A$, sometimes subscripted, etc. The multiplicative identity of algebra $A$ is denoted as $\mathbf{1}_A$ or just $\mathbf{1}$ if the algebra in question is understood.   As a vector, an element $s\in A$ has its dual vector denoted by $s^T$. The transpose of a linear operator $L$ is denoted by $L^T$.  To reduce clutter we utilize expressions like $Ls^{-1}$ which are always interpreted to mean $L(s^{-1})$. The $(n\times n)$-exchange matrix (the matrix with entries having value 1 on the anti-diagonal and 0 everywhere else) is denoted by  $\slashed{\bf 1}_n$.  
 

 \section{Anti-rotors: sorting via an algebra's nonlinear unary operations}
 
 \subsection{The particular interest of the multiplicative inversion operation} An algebra is a vector space with an additional product operation satisfying certain axioms.  Based on intractability of the classification problem, we might conclude that this is too much structure to take on all at once.  So we are going to reduce the bilinear binary product operation to its implied nonlinear unary operations and see what happens. 
 
The implied nonlinear unary operations mapping into the algebra are each given by some linear combination of integer powers of an element, the inverse of some linear combination of integer powers of the element, or the product of a collection of such entities.  Among these, the multiplicative inverse operation stands out as one that is assuredly invertible on the space of units for any algebra.   Thus, focusing on the multiplicative inversion operation seems like a good place to start.  In fact, limiting our attention to {\it only} the multiplicative inversion operation as a replacement for the full algebra binary product operation is not quite as reductive as might be at first imagined, since that restricted focus does not entirely eliminate the influence of arbitrary products of algebra elements.  For example, consider elements $a,b$ where $a$ is a unit.  Using only vector space operations and the multiplicative inverse operation, we can construct the element $(a^{-1}+\epsilon b)^{-1}$, where $\epsilon$ is an appropriate choice of real number.  Let us suppose we are working in a representation, so that $a,b$ are matrices.  Then for small $\epsilon$, use of the \mbox{Neumann} series expansion implies $(a^{-1}-\epsilon b)^{-1} = a+\epsilon(ab)a+c$, where the spectral norm of $c$ is $O(\epsilon^2)$.  Thus, the product $ab$ clearly influences $(a^{-1}+\epsilon b)^{-1}$, despite the fact that the latter is constructed from $a$ and $b$ with only vector space operations and the multiplicative inversion operation.  

Of more immediate significance for our ``bilinear product to nonlinear unary operator replacement" program is the fact that the multiplicative inversion operation can be used to approximate the nonlinear operations defined by taking positive or negative integer powers of an element.  Thus, in the context of a matrix representation it is seen that, for an element $s$, the Neumann series expansion leads to \mbox{$(\mathbf{1}-\epsilon s^{- 1})^{-1} - \mathbf{1}- \epsilon s^{-1} = \epsilon^2 s^{-2} + c$} where the spectral norm of $c$ is $O(\epsilon^3)$, etc.  At least in this sense, it seems that the multiplicative inversion operation, combined only with the vector space operations, subsumes the other nonlinear unary operations implied by the algebra's product operation (we will expand on this in \mbox{Section 5.2}).  
 
      \subsection{The uncurling spaces of an algebra} As indicated in Section 1.2, an anti-rotor can be formulated as a basis-independent entity, but it is quite convenient to present things in the language of matrices.  Consequently, we will proceed under assumption of the standard basis.     
       
         \begin{definition}\label{7/21/23.1} With respect to elements $s$ of a given $n$-dimensional algebra $A$ written in the standard basis as $s= \left(\begin{array}{ccc}
x_1   \\ 
\vdots \\
x_n   
\end{array}\right)$, where these elements are in an open ball in $\mathbb{R}^n$ consisting only of units and whose central point is the multiplicative identity element $\mathbf{1}_A$, an {\it uncurling metric $L$} is a symmetric member of $M_n(\mathbb{R})$  satisfying, \begin{equation}\label{4/7/23.1} d([Ls^{-1}]\cdot {\sf d}s) = d\left((s^{-1})^T L {\sf d}s\right) = 0, \end{equation} where ${\sf d}s \equiv \left(\begin{array}{ccc}
dx_1   \\ 
\vdots \\
dx_n   
\end{array}\right)$.  The real vector space of uncurling metrics is the algebra's {\it uncurling space}, $\mathbf{u}_A$, also known as the {\it group of uncurling metrics}, also known as the {\it anti-rotor}.  It may also be denoted by $\mathbf{u}$ if no confusion arises.    \end{definition}

Equation (\ref{4/7/23.1}) is interpreted according to the conventions stated in Section 1.2.  Its format allows us to streamline some proofs.

 A curl (also known as a rotor) is the Hodge dual of the exterior derivative of the differential \mbox{1-form} dual to a vector field (i.e., the curl is the dual of the infinitesimal rotation tensor).  Thus,  (\ref{4/7/23.1}) states that the curl of $Ls^{-1}$ vanishes (\mbox{$L$ ``uncurls" $s^{-1}$}).    In other words, if $L$ is an uncurling metric then $Ls^{-1}$ is a gradient vector field in some neighborhood of $\mathbf{1}$ (to see that an open ball centered at $\mathbf{1}$ and devoid of non-units exists for any algebra, it is sufficient to observe that such a neighborhood clearly exists for the left regular representation of the algebra - where one works within the matrix algebra $ M_n(\mathbb{R})$, and the inverse of an element is simply the reciprocal of the determinant multiplied by the adjugate matrix).
 
 We could similarly define other ``types" of uncurling spaces as associated with any of the algebra's implied nonlinear unary operations.  However, as suggested in Section 2.1, this might not lead to any further insights (also see the Conjecture stated in \mbox{Section 5.2}).  
 
While only matrix addition and scalar multiplication figure in an uncurling space in isolation, the fact that the members of the space share a particular analytical property leads to another layer of structure dependent on the matrix product - which comes to the fore when distinct spaces of uncurling metrics are being compared (which ultimately is the whole point).  This is forcefully illustrated in Theorem \ref{6/6/23.1}, whose implications are expressed in categorical language in Remark \ref{3/31/24.1}. Given its additional layer of structure, an uncurling space is considerably more than a vector space.  Not surprisingly, an uncurling space is often more appropriately viewed as an object in a particular category where the morphisms make intrinsic use of the matrix product (as in Section 2.9).    
We can think of  ``anti-rotor" as emphasizing its meaning as an object subject to such a morphism, ``uncurling space" as recalling its status as a vector space, and ``group of uncurling metrics" as indicating the role of its members in facilitating integration (as in Section 2.6).
 
 
  \subsection{Where we are headed}
    Though our concerns are much broader, it is nonetheless significant that the uncurling space provides a direct route to a generalization of the usual algebra norm (a determinant), so we will give the gist of that result now to provide a taste of where this program leads.  
 
 By definition, we are able to perform path-independent integration of members of the uncurling space applied to the multiplicative inversion operation as in (\ref{11/7/20.2}), alternatively expressed as (\ref{11/7/20.1}).  
We are then led to ``special" functions $\ell_{\rm sp}(s)$ that result from ``normalized" members $L_{\shortparallel}$ of the uncurling space, and these have the \mbox{``determinant-lite"} features of positive homogeneity along with ``inversion consistency", $\ell_{\rm sp}(s^{-1}) = \frac{1}{\ell_{\rm sp}(s)}$.  If $L_{\shortparallel}$ is nonsingular, the inversion consistency implies that (\ref{11/7/20.1}) becomes (\ref{5/22/24.5}), from which the positive homogeneity feature further implies (\ref{5/22/24.6}).  These last two equations characterize each $\ell_{\rm sp}(s)$, indicating that each has a strong kinship with the norm on an inner product space (whose analogous characterizing equations appear in Appendix A as (\ref{11/29/23.1}),  (\ref{11/29/23.2})). Furthermore, based on the above inversion consistency feature, (\ref{5/22/24.5}) can be alternatively written as,  \begin{equation}\label{6/5/24.1} \nabla\ell_{\rm sp}(s) = \frac{1}{\|\mathbf{1}\|^{2}}\ell_{\rm sp}(s) L_{\shortparallel}s^{-1}.\end{equation} This can be understood as a generalization of a consequence of Jacobi's formula, where in the latter case $s$ is a matrix, the determinant of $s$ plays the role of $\ell_{\rm sp}(s)$, and the matrix transpose operation plays the role of $L_{\shortparallel}$ (this all being in the context of a convention regarding the interchange of matrix and vector formats of $s$  -  see the second paragraph of Section 3, and comparison of (\ref{6/5/24.1}) with (\ref{11/5/21.1})).  In other words, (\ref{6/5/24.1}) indicates that each of the functions that derives from some member of the set of normalized uncurling metrics can be thought of as a generalization of the usual algebra norm.  Finally, as shown in Section 3, there is always a normalized uncurling metric that leads to one of the above special functions being closely allied with the algebra's usual norm, a feature also implying that the uncurling space is always at least one-dimensional.  In fact, apart from the case of simple algebras, an uncurling space is typically multi-dimensional.  
      
The dimension of the uncurling space/anti-rotor is an algebra isomorphism invariant.  Essentially, our program concerns an analysis of real unital associative algebras based on invariants deriving from their anti-rotors.      

\begin{remark}\label{11/21/23.1} {\rm If we were to replace $s^{-1}$ with $s$ in (\ref{4/7/23.1}), the resulting ``trivial uncurling space" would not address the given algebra, since the mapping defined by $s\mapsto s$ does not depend on the algebra product.  This would seem to be pointless given the desire to derive an entity sensitive to the algebra and not just the algebra's underlying vector space.  However, the latter trivial uncurling space does forward equations (i.e., (\ref{11/29/23.1}), (\ref{11/29/23.2}), and ultimately, (\ref{11/29/23.5}), (\ref{11/29/23.6})) that can serve as the basis for a new proof of the Pythagorean Theorem, as presented in Appendix A.  It is also seen in Appendix A that, viewed in isolation, the mechanism underlying the new proof can serve as a motivation for our whole program.  However, the original motivation for this program is presented in Appendix B.  }\end{remark}

\subsection{Anti-rotors and algebra isomorphism invariants} The following result, which describes how an anti-rotor transforms, is fundamental.


\begin{theorem} \label{6/6/23.1} If $K:A_1\rightarrow A_2$ is an algebra isomorphism, then \begin{equation}\label{5/23/23.1} \mathbf{u}_{A_1} = K^T \mathbf{u}_{A_2} K. \end{equation} \end{theorem}

\begin{proof} By hypothesis, a nonsingular linear change of variable expression for the elements of $A_2$ in terms of the elements of $A_1$ can be written as $s_2= Ks_1$. \mbox{For $L_2\in \mathbf{u}_{A_2}$} we obtain, \begin{equation} 0 = d\left((s_2^{-1})^T L_2 {\sf d}s_2\right) =  d\left(((Ks_1)^{-1})^T L_2 {\sf d}(Ks_1)\right) = d\left((s_1^{-1})^T \big[K^T L_2K\big]{\sf d}s_1\right), \label{9/21/23.1}  \end{equation} where we have used Definition \ref{7/21/23.1}, the commutativity of a differential with a linear transformation, and $(Ks_1)^{-1} = K(s_1^{-1})$ (i.e., an algebra isomorphism maps a multiplicative identity to a multiplicative identity, so that ${\bf 1}_{A_2} = K\left(\mathbf{1}_{A_1}\right) = K(s_1s_1^{-1}) = K(s_1)K(s_1^{-1})$).  It follows from (\ref{9/21/23.1}) that every uncurling metric  of algebra $A_2$ is associated with an uncurling metric of $A_1$, i.e., from the last equality in (\ref{9/21/23.1}) we have  $K^TL_2K\in\mathbf{u}_{A_1}$.  An analogous argument employing $K^{-1}$ indicates that for $L_1\in\mathbf{u}_{A_1}$ we have the associated $(K^{-1})^T L_1 K^{-1} \in \mathbf{u}_{A_2}$.    Furthermore, $K^T L_2K\in\mathbf{u}_{A_1}$ is such that $(K^{-1})^T [K^T L_2 K]K^{-1} = (K^T)^{-1} [K^T L_2 K]K^{-1} = L_2$.  Similarly, $(K^{-1})^T L_1K^{-1}\in\mathbf{u}_{A_2}$ is such that $K^T[(K^{-1})^T L_1K^{-1}] K = L_1$.  The theorem follows from the implied bijection. 
\end{proof}

\begin{remark} \label{3/31/24.1} {\rm Theorem \ref{6/6/23.1} indicates that the association of algebras with uncurling spaces is functorial.  That is, let $\mathcal{A}$ be the category each of whose objects is an algebra, where the morphisms are algebra isomorphisms, and let $\mathcal{S}$ be the category each of whose objects {\sf S} is a finite-dimensional vector space of real symmetric matrices, where $\begin{CD}  {\sf S}_1 @>\hat{K}>> {\sf S}_2\end{CD}$ is a morphism if there exists a nonsingular real matrix $\hat{K}$ (given the same symbol as the morphism) such that \mbox{${\sf S}_1  = \hat{K}^T  {\sf S}_2 \hat{K}$.}  It is easily verified that (\ref{5/23/23.1}) implies an extension of the mapping of algebras to uncurling spaces, to a functor from $\mathcal{A}$ to $\mathcal{S}$.  This leads to a class of algebra isomorphism invariants that in essence result via secondary functors from $\mathcal{S}$ to the trivial category $\mathbb{R}$ (whose points are objects, and a morphism points only from a point to itself).  Corollary \ref{8/18/23.2} can be interpreted as listing some examples.  A related but more capable functor will be presented in Section 2.9.}\end{remark}

Computation of $\mathbf{u}$ is straightforward in principle.  The algorithm is contained in the proof of the following. \enlargethispage{1.5cm}

\begin{proposition}\label{9/2/24.1} The members of $\mathbf{u}$ are the realizations of a symmetric $(n\times n)$-matrix $M_{\mathbf{u}}$ whose entries are linear combinations of the components of an $m$-tuple of real parameters, where $m$ is the dimension of $\mathbf{u}$, and where each realization of $M_{\mathbf{u}}$ as a real matrix results from replacing its $m$ parameters with the members of any set of $m$ members of $\mathbb{R}$.   
\end{proposition}

\begin{proof} For an $n$-dimensional algebra with elements $s = (x_1,x_2,\dots,x_n)$ with respect to the standard basis, one initially takes $\mathbf{u}$ to be a symmetric $(n\times n)$-matrix $M$ whose upper triangular entries are each a different member of a set of $\frac{n(n+1)}{2}$ real parameters.  Computation of $s^{-1}$ produces a vector whose components are each a rational function of the $x_1,x_2,\dots,x_n$.  So far, $Ms^{-1}$ is a parametrized \mbox{$n$-dimensional} vector field whose components (which can easily be computed) are dependent on the initial set of parameters in $M$ and the $x_1,x_2,\dots,x_n$.  There are then $\frac{n(n-1)}{2}$ constraint equations that must be satisfied in order that (\ref{4/7/23.1}) hold (for $L=M$).  Each of these equations has the form, \begin{equation}\label{9/2/24.2} \frac{\partial}{\partial x_i}([Ms^{-1}]_j) - \frac{\partial}{\partial x_j}([Ms^{-1}]_i) = 0,\end{equation} where $[Ms^{-1}]_k$ indicates the $k$-th component of $Ms^{-1}$.  Evidently, we can write $s^{-1} = (\hat{x}_1,\dots,\hat{x}_n)$, where each $\hat{x}_k$ is a rational function of the components of $s$.  Consequently, $[Ms^{-1}]_k$ is a linear combination of the components $\hat{x}_1,\dots,\hat{x}_n$, where the coefficients defining this linear combination are various of the above parameters.  It follows that $\frac{\partial}{\partial x_p}([Ms^{-1}]_q)$ is a rational function of the components $x_1,\dots,x_n$, but this expression is linear in the parameters.  Consequently, multiplying both sides of (\ref{9/2/24.2}) by an appropriate polynomial in the components of $s$ that is not dependent on the parameters, the left-hand-side of (\ref{9/2/24.2}) becomes a polynomial in the components of $s$, while remaining an expression that is linear in the parameters - and the right-hand-side of (\ref{9/2/24.2}) remains zero.
Thus, satisfaction of each of the $\frac{n(n-1)}{2}$ constraint equations (\ref{9/2/24.2}) ultimately reduces to the requirement that a particular multivariate polynomial be {\it identically} zero.  After simplification of the polynomial to be a linear combination of independent monomials, each monomial term is then comprised of a coefficient consisting of a linear combination of some of the above $\frac{n(n+1)}{2}$ parameters, multiplying an expression consisting of a product of some of the components of $s$, each raised to some power.  Therefore, each of the above coefficients of these (independent) monomials are required to be zero.  Accordingly, the number of surviving independent parameters (i.e., parameters not required to be equated to some linear combination of other parameters, or set to 0) is in general reduced as the $\frac{n(n-1)}{2}$ constraint equations are successively treated.  $M$ is then modified as necessary, as each successive constraint equation is treated.  After treatment of all the constraint equations via the above steps defined by treatment of the constraint equations in some sequence, the final number $m$ of surviving independent parameters is the dimension of $\mathbf{u}$, and the final entries of $M$ evidently yield the requisite $M_\mathbf{u}$.
\end{proof}

Consequently, $\mathbf{u}$ can be represented as a symmetric $(n\times n)$-matrix $M_{\mathbf{u}}$ whose entries are particular linear combinations of $m$ independent parameters (where $m$ is the dimension of $\mathbf{u}$), because each uncurling metric results from a single set of selected real values for those parameters.  The determinant of $M_{\mathbf{u}}$, written as ${\rm det}(M_{\mathbf{u}})$, is a polynomial in the components of an ordered $m$-tuple of parameters, $(\alpha_1,\dots,\alpha_m)$.  If $m$ is small  we will denote the parameters using different Greek letters like $\alpha$, $\beta$, $\gamma$, etc. 

\bigskip

   With the above in mind, it is seen that (\ref{5/23/23.1}) has immediate consequences. 
                 
  \begin{corollary}\label{8/18/23.2} The following are algebra isomorphism invariants. \begin{itemize} \item The dimension $m$ of $\mathbf{u}$.   \item The largest rank  of a member of $\mathbf{u}$.  \item The smallest rank of a nonzero member of $\mathbf{u}$. \item The dimension of the subspace of the algebra to which ${\rm det}(M_{\mathbf{u}})$ is sensitive. \item Invariants of the real algebraic set, \[V_\mathbb{R}({\rm det}(M_{\mathbf{u}})) = \{(\alpha_1,\dots,\alpha_m)\in\mathbb{R}^m: {\rm det}(M_{\mathbf{u}}) = 0 \}.\] \end{itemize} \end{corollary}
  
   It is shown in Section 3 that $\mathbf{u}$ is always at least one-dimensional.   Concrete examples of computation of $\mathbf{u}$ are found in Section 2.7 and Section 4.  
  
\subsection{Combinatorics} As seen in the proof of Proposition \ref{9/2/24.1}, computation of $\mathbf{u}$ in the form of $M_\mathbf{u}$ is straightforward, as is the dimension of $\mathbf{u}$.   Computation of the largest rank of a member of $\mathbf{u}$ is also easy, e.g., using the method of minors.  Computation of the smallest rank of a nonzero member of $\mathbf{u}$ could be problematic, since it might require evaluation of the dimension of a complicated algebraic set (e.g., potentially involving evaluation of the simultaneous vanishing of many minors).  Evaluation of the number of parameters to which ${\rm det}(M_{\mathbf{u}})$ is sensitive is simply a matter of inspection.  Examples of invariants derived from $V_\mathbb{R}({\rm det}(M_{\mathbf{u}}))$ include the real dimension of that real hypersurface (for which an efficient algorithm is \mbox{available \cite{bannwarth:2015})}, and the number of irreducible components of $V_\mathbb{R}({\rm det}(M_{\mathbf{u}}))$.  
 
Various combinatorial questions now arise based on the invariants identified in Corollary \ref{8/18/23.2}.  Examples relevant to the set of algebras of dimension $n$ include, \begin{enumerate} \item How are the algebras distributed between the classes defined by uncurling space dimension (what are the relative popularities of these classes)?  \item For algebras having the same uncurling space dimension, how are they further distributed into classes defined by different numbers of irreducible components of $V_\mathbb{R}({\rm det}(M_{\mathbf{u}}))$?  \item For which algebras is ${\rm det}(M_{\mathbf{u}})$ identically $0$?   \item For those algebras having the same uncurling space dimension $m$, and for each $k\le m$, how likely is $k$ to be the maximal rank of a member of $\mathbf{u}$?   \end{enumerate}  
More combinatorial questions arise from the additional invariants identified in Corollary \ref{7/25/24.1} below. 

Concerning Question (1) above, relating to the distribution of $n$-dimensional algebras in the bins defined by the dimension $\text{dim}(\mathbf{u})$ of their uncurling spaces, we can observe that $\mathbf{u}$ is derived from an initial set of $\frac{n(n+1)}{2}$ parameters, but these are then subjected to  $\frac{n(n-1)}{2}$ constraint equations.  One might then expect to be finally left with $\frac{n(n+1)}{2} - \frac{n(n-1)}{2} = n$ parameters - i.e., $\text{dim}(\mathbf{u}) = n$.  This reasoning is incorrect.  Consider the analogous problem of finding the matrices $L$ that satisfy (\ref{4/7/23.1}) except that $s^{-1}$ is replaced with $s$ in that equation.  One quickly discovers that every constraint equation reduces to $0=0$, so that {\it every} symmetric $n\times n$-matrix is a solution (the dimension of the solution space is $\frac{n(n+1)}{2}$).  So, if some of the constraint equations for uncurling metrics were subject to  that kind of thing, or if the $\frac{n(n-1)}{2}$ constraints did not constitute a set of $n$ independent constraints, perhaps $\text{dim}(\mathbf{u})$ could be some number greater than $n$.  Conversely, it might occur that a constraint equation eliminates more than one parameter, so that $\text{dim}(\mathbf{u})$ turns out to be less than $n$ (along these lines, see the proof of Theorem \ref{4/8/22.1}).  On the other hand, we still might guess that $\text{dim}(\mathbf{u}) = n$ would be the most popular bin.  In fact, for all three isomorphism classes for the two-dimensional algebras, Table I indicates that $\text{dim}(\mathbf{u}) = 2$.  For the three-dimensional algebras, Table II indicates that five of the six algebra classes satisfy $\text{dim}(\mathbf{u}) = 3$, although the sixth satisfies $\text{dim}(\mathbf{u}) = 2$.  Moreover, as regards the four-dimensional algebras, it turns out that $\text{dim}(\mathbf{u}_{M_2(\mathbb{R})}) = 1$ (as follows from Theorem \ref{4/8/22.1}). So a precise understanding of the distribution of $n$-dimensional algebras into the bins defined by $\text{dim}(\mathbf{u})$ remains a goal.  

However, Corollary \ref{8/17/22.1} is a first general result relevant to a different aspect of Question (1), as it indicates that if the uncurling metric group is one-dimensional and has a nonsingular member then the algebra is simple.  



\subsection{The unital norm group of an algebra} By definition, for \mbox{$L\in \mathbf{u}_A$}, the vector field $Ls^{-1}$ has a vanishing curl in some neighborhood of $\mathbf{1}$.  It is thus associated with an ``energy density function" given by its integral (borrowing nomenclature for the integral of a conservative vector field arising in the realm of physics).   
 We can notate the abelian group of these energy density functions as $\mathbf{E}_{\mathbf{u}_A}$.  In fact, using the $m$ parameters of $M_{\mathbf{u}_A}$, the group $\mathbf{E}_{\mathbf{u}_A}$ can be thought of as a kind of single comprehensive energy function associated with $A$.  But since it ultimately has the interesting algebraic interpretation alluded to in Section 2.3 and demonstrated in Section 3, we will instead utilize the following alternative notation and nomenclature.
  
Our later Definition \ref{9/20/23.1} introduces the notion of a ``normalized" uncurling metric.  Its specification will employ the computation of a positive integer $\|\mathbf{1}\|^2 \equiv \mathbf{1} \cdot \mathbf{1}$. As preparation, this factor must first figure in \mbox{Definition \ref{9/21/23.2}}.   Note that $\|\mathbf{1}\|^2$ is not an algebra isomorphism invariant.  That is, for isomorphic algebras $A_1,A_2$, it may be the case that $\|\mathbf{1}_{A_1}\|^2 \ne \|\mathbf{1}_{A_2}\|^2$ (in which case the isomorphism is not isometric).  This factor does not enter into the main results of this subsection, \mbox{Theorem \ref{6/6/23.2}} and Corollary \ref{7/25/24.1}.
  
  \begin{definition}\label{9/21/23.2} Let $L$ be an uncurling metric for some algebra $A$.  For $s$ in an open ball centered at $\mathbf{1}$ and devoid of non-units, the associated {\it unital norm} $ \ell(s)$ is the function 
  \begin{equation}\label{11/7/20.2}   \ell(s) \equiv \exp\left(\frac{1}{\|{\bf 1}\|^2}\int_{\bf 1}^s [Lt^{-1}]\cdot {\sf d}t\right).\end{equation} 
   The set of unital norms associated with an algebra, which is an abelian group $\mathbf{U}_A$ with respect to the usual point-wise operation of multiplication of functions, is called the {\it unital norm group}. $\mathbf{U}_A$ may be denoted as $\mathbf{U}$ if no confusion arises. \end{definition} 
   
  
    From (\ref{11/7/20.2}), an uncurling metric and its associated unital norm satisfy  \begin{equation}\label{11/7/20.1} Ls^{-1}  = \nabla\left[\|{\bf 1}\|^2\log  \ell(s)\right] =   \left(\frac{\|{\bf 1}\|^2}{ \ell(s)}\right)\nabla \ell(s).\end{equation} 
Thus, the pair $(L, \ell(s))$ characterizes the algebra's multiplicative inversion operation $s\mapsto s^{-1}$ in the portion of an open ball centered at $\mathbf{1}$ devoid of non-units that is orthogonal to the kernel of $L$.  If $L$ is nonsingular, this pair fully characterizes the inversion operation $s^{-1}$ in an open ball centered at $\mathbf{1}$ devoid of non-units.  For a given algebra, the set of all such pairs (where $L$ may or may not be singular) is an abelian group with respect to the binary operation 
\begin{equation}\label{5/23/23.2} \big(L_1, \ell_1(s)\big)\big(L_2, \ell_2(s)\big) \equiv 
\big(L_1+L_2, \ell_1(s)  \ell_2(s)\big).\end{equation} 
\begin{definition} For a given algebra, the set of pairs $(L, \ell(s))$ satisfying (\ref{4/7/23.1}) and (\ref{11/7/20.2}), in the context of the operation defined by (\ref{5/23/23.2}), is the algebra's {\it inversion characterization group}. \end{definition}

Any member of the inversion characterization group that is associated with a nonsingular uncurling metric implies the full group, since the member implies the operation $s\mapsto s^{-1}$ via (\ref{11/7/20.1}), from which the full inversion characterization group can be computed.  Obviously, this group is {\it not} $\mathbf{u}\oplus\mathbf{U}$. 

Corollary \ref{8/18/23.2} has already identified invariants associated with $\mathbf{u}$ (whose members supply the first component of the inversion characterization group), but another set of useful invariants becomes apparent on examination of $\mathbf{U}$ (whose members supply the second component of the inversion characterization group).  Based on the mapping from $\mathbf{u}$ to $\mathbf{U}$ defined by (\ref{11/7/20.2}),  we can write \begin{equation}  \mathbf{U} = \exp\left(\int_{\bf 1}^s [\mathbf{u}t^{-1}]\cdot {\sf d}t\right). \label{6/7/23.1}\end{equation}  In the above, the notation $\mathbf{u}t^{-1}$ means $\{Lt^{-1}: L\in \mathbf{u}\}$.  Accordingly, the right-hand-side of (\ref{6/7/23.1}) is understood to be a collection of exponentiated integrals.   Note that in (\ref{6/7/23.1}) we do not need the factor $\frac{1}{\|{\bf 1}\|^2}$, since $\frac{1}{\|{\bf 1}\|^2}\mathbf{u}=\mathbf{u}$.  



The invariants stemming from the $\mathbf{U}$ contribution to the inversion characterization group proceed from the following simple corollary of Theorem \ref{6/6/23.1}, which we will give ``co-equal status" as a theorem because it ultimately addresses the co-important  partner of $\mathbf{u}$, i.e., $\mathbf{U}$. 

\begin{theorem}\label{6/6/23.2} Let $K:A_1\rightarrow A_2$ be an algebra isomorphism, where for $s_1\in A_1$ we define $s_2\equiv Ks_1$.  Then for $s_1$ in an open ball centered at $\mathbf{1}_{A_1}$ consisting only of units, \begin{equation}\label{9/6/24.2} \int_{{\bf 1}_{A_1}}^{s_1} \left[\mathbf{u}_{A_1}t_1^{-1}\right]\cdot {\sf d}t_1 
= \int_{{\bf 1}_{A_2}}^{s_2} \left[\mathbf{u}_{A_2}t_2^{-1}\right]\cdot {\sf d}t_2 . \end{equation} \end{theorem}

\begin{proof} 
Equation (\ref{5/23/23.1}) implies,
\begin{eqnarray*} \int_{{\bf 1}_{A_1}}^{s_1} \left[\mathbf{u}_{A_1}t_1^{-1}\right]\cdot {\sf d}t_1 &=& \int_{{\bf 1}_{A_1}}^{s_1} \left[ K^T\mathbf{u}_{A_2}Kt_1^{-1}\right]\cdot {\sf d}t_1 \\ &=& 
\int_{{\bf 1}_{A_2}}^{s_2} \left[ K^T\mathbf{u}_{A_2}K (K^{-1}t_2)^{-1}\right]\cdot dK^{-1}t_2 \\ &=& \int_{{\bf 1}_{A_2}}^{s_2} [K^T\mathbf{u}_{A_2}t_2^{-1}]\cdot \left(K^{-1}dt_2\right) = \int_{{\bf 1}_{A_2}}^{s_2} \left[\mathbf{u}_{A_2}t_2^{-1}\right]\cdot {\sf d}t_2, 
\end{eqnarray*} 
where we have used $(K^{-1}t_2)^{-1} = K^{-1}t_2^{-1}$ (i.e., $\mathbf{1}_{A_1} = K^{-1}\mathbf{1}_{A_2} = K^{-1}(t_2t_2^{-1}) = (K^{-1}t_2)(K^{-1}t_2^{-1})$), the commutativity of a linear transformation and a differential, and exploitation of an adjoint in the context of an inner product.
\end{proof}



Applying (\ref{9/6/24.2}) to (\ref{6/7/23.1}), it is clear that for isomorphic algebras $A_1,A_2$ the groups $\mathbf{U}_{A_1},  \mathbf{U}_{A_2}$ ``look the same" after accounting for the effect of the isomorphism on the vector space of elements of the algebras.

More specifically, Proposition \ref{9/2/24.1} indicates that $m$-dimensional $\mathbf{u}$ can be written as a matrix $M_\mathbf{u}$ whose entries are each a linear combination of the components of an $m$-vector of real parameters $(\alpha_1, \dots, \alpha_m)$.   Thus, $\mathbf{u}$ can be equivalently thought of as a set of real matrices, each of whose members is given by a choice of real values for the $m$ parameters in the above parametrized matrix version of $\mathbf{u}$, with each matrix in the set being an uncurling metric.  Now, for each $q\le m$ consider the real matrix  $\mathbf{u}_q$ obtained from $\mathbf{u}$ by setting $\alpha_q =1$ with all other parameters set to zero.  These are a set of generators for $\mathbf{u}$. Thus, $\mathbf{u}t^{-1} = M_{\mathbf{u}}[t^{-1}] = \sum_{q=1}^m \alpha_q \mathbf{u}_q t^{-1}$.   Using (\ref{6/7/23.1}), we then obtain \begin{equation}\label{9/26/23.1} \log  \mathbf{U} = \int_{\mathbf{1}}^s (M_{\mathbf{u}} [t^{-1}])\cdot {\sf d}t =\sum_{q=1}^m \alpha_q \int_{\mathbf{1}}^s [\mathbf{u}_q t^{-1}]\cdot {\sf d}t.\end{equation}  

We now consider the algebra elements in an open ball centered at $\mathbf{1}$ consisting only of units, with each element written in coordinates as $s = (s_1,\dots,s_n)$, where $({\sf 1}_1,{\sf 1}_2,\dots,{\sf 1}_n) \equiv \mathbf{1}$.  The integral inside the sum on the right-hand-side of (\ref{9/26/23.1}) can be evaluated on the piecewise linear path with successive endpoints $({\sf 1}_1,{\sf 1}_2,\dots,{\sf 1}_n)$, $(s_1,{\sf 1}_2,\dots,{\sf 1}_n)$, $(s_1,s_2,{\sf 1}_3,\dots,{\sf 1}_n)$,  $\dots$, $(s_1,\dots, s_{n-1}, {\sf 1}_n)$, $(s_1,s_2,\dots,s_n)$.    On a linear portion of the path, the integrand is a rational function of a single variable.  Thus, on each linear path component the indefinite integral will be a sum of rational functions, logarithms, and arctangents.  However, coefficients in that sum will be functions of the components of $s$ that do not vary on the linear piece of the path.  For dimension less than five, we can still obtain a closed form expression for the right-hand-side of (\ref{9/26/23.1}) as a function of the components of $s$ (as in Table I and Table II of Section 2.7, for dimensions two and three).  However, this is not possible in general for dimension five and greater, given the need for an expression for the roots of a polynomial appearing in the denominator of an integrand - these roots being functions of the nonvarying components of $s$ (of course, if definite values are given for the non-varying components of $s$, the integration on each linear path component can be accomplished to any degree of desired accuracy via well-known polynomial root-finding algorithms).

But while lack of a closed form solution for the roots of a polynomial in dimensions five and above limits the ability to write down a closed form expression for the full function $\log\mathbf{U}$, we can evaluate (to any degree of accuracy) the set of $n$ functions given by $\log\mathbf{U}$ computed along paths each beginning at $\mathbf{1}$ and proceeding parallel to each of the component axes of $s$ - because each integrand again has only a single component of $s$ that is varying on the path, but now the other components of $s$ only appear as their known component values in $\mathbf{1}$.  
So, if we successively compute the integrals on the linear path from $\mathbf{1}$ to $(s_1,{\sf 1}_2,\dots,{\sf 1}_n)$, then on the linear path from \mbox{$\mathbf{1}$ to} $({\sf 1}_1,s_2,{\sf 1}_3,\dots,{\sf 1}_n)$, etc., we obtain this set of $n$ functions. 
\enlargethispage{1cm} Now, evaluation of $q$-th term on the right-hand-side of (\ref{9/26/23.1}) is \begin{equation}\label{7/26/24.1} \alpha_q\int_{\mathbf{1}}^s [\mathbf{u}_q t^{-1}]\cdot {\sf d}t = \alpha_q\sum_{k=1}^n \int_{\bar{s}_k} [\mathbf{u}_q t^{-1}]\cdot {\sf d}t_k,\end{equation} with $\bar{s}_k \equiv [(s_1,\dots,s_{k-1},{\sf 1}_{k},{\sf 1}_{k+1}\dots,{\sf 1}_n),(s_1,\dots,s_{k-1},s_k,{\sf 1}_{k+1},\dots,{\sf 1}_n)]$ being the directed line segment with the indicated endpoints.   As we have noted, evaluation of the integrals inside the sum on the right-hand-side of (\ref{7/26/24.1}) is problematic.  But instead consider,  \begin{equation}\label{7/26/24.2}\alpha_q\sum_{k=1}^n \int_{\hat{s}_k} [\mathbf{u}_q t^{-1}]\cdot {\sf d}t_k,\end{equation} with $\hat{s}_k \equiv [({\sf 1}_1,\dots,{\sf 1}_{k-1},{\sf 1}_{k},{\sf 1}_{k+1}\dots,{\sf 1}_n),({\sf 1}_1,\dots,{\sf 1}_{k-1},s_k,{\sf 1}_{k+1},\dots,{\sf 1}_n)]$. 
For reasons stated above, approximations to the integrals in (\ref{7/26/24.2}) in terms of rational functions, logarithms, and arctangents, can be obtained with arbitrary accuracy.  Furthermore, for $s$ in a small enough neighborhood of $\mathbf{1}$, the integrands of each integral in (\ref{7/26/24.2}) are close to the integrands appearing in each integral on the right-hand-side of (\ref{7/26/24.1}).  This means that there will be the same number of respective rational function, logarithm, and arctangent terms in the evaluated indefinite integral version of the $k$-th integral in the sum (\ref{7/26/24.2}) as there would be if one evaluated the indefinite integral version of the $k$-th integral in the sum on the right-hand-side of (\ref{7/26/24.1}) (i.e., if any definite values for the non-varying components of $s$ compatible with the neighborhood of $\mathbf{1}$ were inserted into the integrand of (\ref{7/26/24.1})).  Consequently, a parameter $\alpha_q$ multiplies a rational function (resp., logarithm, arctangent) in the $k$-th term of the sum in (\ref{7/26/24.2}) if and only if $\alpha_q$ multiplies a rational function (resp., logarithm, arctangent) in the $k$-th term of the sum on the right-hand-side of (\ref{7/26/24.1}) were the integral actually computed for any specific values of the non-varying components of $s$ in the range of values compatible with the neighborhood of $\mathbf{1}$.  

So, as regards evaluation of the association of parameters of $M_{\mathbf{u}}$ with rational function, logarithm, and arctangent terms arising from $\log\mathbf{U}$, instead of using the right-hand-side of (\ref{9/26/23.1}) it is sufficient to instead deal with the function,
\begin{equation}\label{7/29/24.1} \sum_{q=1}^m \big(\alpha_q\sum_{k=1}^n \int_{\hat{s}_k} [\mathbf{u}_q t^{-1}]\cdot {\sf d}t_k\big).\end{equation} Again, following evaluation of the integrals, this function is given by respective sums of rational function terms, logarithm terms, and arctangent terms.  Note that for this function resulting from algebra $A$, the rational function terms can be collected into a single rational function term involving some number $\tau_{A,{\rm rat}}$ of distinct parameters.  Furthermore, its logarithm terms can be collected into a single (real) logarithm involving some number $\tau_{A,{\rm log}}$ of distinct parameters.  Similarly, its arctangent terms, which can be alternatively expressed utilizing $\text{arctan}(x) = \frac{i}{2}\log\frac{1-ix}{1+ix}$, can be collected into a single such logarithm term utilizing imaginary coefficients, involving some number $\tau_{A,{\rm arc}}$ of distinct parameters. 
Based on dimensional considerations relating to the respective numbers of distinct parameters associated with the rational function term, the real logarithm term, and the logarithm term with imaginary coefficients, Theorem \ref{6/6/23.2} then implies


\begin{corollary}\label{7/25/24.1} With respect to the set of parameters $\{\alpha_q:q=1,\dots,m\}$ arising in $M_{\mathbf{u}_A}$, and the function defined by (\ref{7/29/24.1}) following evaluation of the integrals, let $\tau_{A,{\rm rat}}$ be the number of distinct parameters multiplying at least one rational function term, let $\tau_{A,{\rm log}}$ be the number of distinct parameters multiplying at least one logarithm, and let $\tau_{A,{\rm arc}}$ be the number of distinct parameters multiplying at least one arctangent.  If $A_1$ and $A_2$ are isomorphic algebras, then \[ (\tau_{A_1,{\rm rat}},\tau_{A_1,{\rm log}},\tau_{A_1,{\rm arc}}) = (\tau_{A_2,{\rm rat}},\tau_{A_2,{\rm log}},\tau_{A_2,{\rm arc}}).\]
\end{corollary}

The invariants identified in Corollary \ref{7/25/24.1}, derived from the unital norm group $\mathbf{U}$, supplement the invariants in Corollary \ref{8/18/23.2} derived from the uncurling metric group $\mathbf{u}$.


\medskip

\subsection{Illustrations using the two-dimensional and three-dimensional algebras}
Up to an isomorphism, there are only three two-dimensional algebras.
For the algebra of complex numbers, $\mathbb{C}$, we can write its elements $s$ as $(x,y)$, with the product given by $(x,y)(z,w) = (xz-yw, yz+xw)$, with $\mathbf{1}=(1,0)$, so that $\|\mathbf{1}\|^2=1$.  Also, $(x,y)^{-1} = (\frac{x}{x^2+y^2},\frac{-y}{x^2+y^2})$.  According to the program indicated immediately prior to Corollary \ref{8/18/23.2}, we start out by writing  $\mathbf{u}_\mathbb{C} = \left[\begin{array}{ccc}
 \alpha & \beta   \\ 
 \beta & \gamma    
\end{array}\right]$.  \mbox{We then} obtain  $\mathbf{u}_\mathbb{C}s^{-1} = \left[\begin{array}{ccc}
 \alpha & \beta   \\ 
 \beta & \gamma    
\end{array}\right]\left(\begin{array}{ccc}
x   \\ 
y   
\end{array}\right)^{-1} = \left(\begin{array}{ccc}
\frac{\alpha x - \beta y}{x^2+y^2}  \\ 
\frac{\beta x - \gamma y}{x^2+y^2} 
\end{array}\right)$.  As regards the right-hand-side of this equation, satisfaction of (\ref{4/7/23.1}) requires that the partial derivative with respect to $x$ of the second component must equal the partial derivative with respect to $y$ of the first component.  Consequently, $\mathbf{u}_{\mathbb{C}} = \left[\begin{array}{ccc}
 \alpha & \beta   \\ 
 \beta & -\alpha.    
\end{array}\right]$.  Using Definition \ref{9/21/23.2}, a simple exercise in integration leads to $\log  \mathbf{U}_{\mathbb{C}} = \frac{\alpha}{2}\log(x^2+y^2) + \beta{\rm arctan}(\frac{y}{x})$.

The split-complex numbers $\slashed{\mathbb{C}}$ constitute an algebra for which the elements are  $(x,y)$ with product $(x,y)(z,w) = (xz+yw,xw+yz)$, $ \mathbf{1} = (1,0)$, and $\|\mathbf{1}\|^2 = 1$.  Also, $(x,y)^{-1} = (\frac{x}{x^2-y^2},\frac{-y}{x^2-y^2})$. \hspace{1cm}  We then find $\mathbf{u}_{\slashed{\mathbb{C}}} = \left[\begin{array}{ccc}
 \alpha & \beta   \\ 
 \beta & \alpha.    
\end{array}\right]$, \hspace{.5cm} and \hfill \\ $\log  \mathbf{U}_{\slashed{\mathbb{C}}} = \frac{\alpha}{2}\log[(x+
y)(x-y)] +  \frac{\beta}{2}\log \frac{x+y}{x-y}$.  In fact, this can be simplified to $\log  \mathbf{U}_{\slashed{\mathbb{C}}} = \frac{\alpha+\beta}{2}\log(x+y) + \frac{\alpha-\beta}{2}\log(x-y)$.  Of course, $\frac{\alpha+\beta}{2}$, $\frac{\alpha-\beta}{2}$ are themselves two independent parameters that we could name $\alpha'$, $\beta'$.

The final real two-dimensional unital algebra (up to an isomorphism) is the algebra of dual numbers $\mathbb{D}$.  The elements are $(x,y)$ with product given by $(x,y)(z,w) = (xz,xw+yz)$, with $\mathbf{1}=(1,0)$ and $\|\mathbf{1}\|^2 = 1$. Also, $(x,y)^{-1} = (\frac{1}{x},-\frac{y}{x^2})$.  We then find $\mathbf{u}_\mathbb{D} = \left[\begin{array}{ccc}
 \alpha & \beta   \\ 
 \beta & 0    
\end{array}\right]$ and $\log  \mathbf{U}_{\mathbb{D}} = \alpha\log x + \beta\left(\frac{y}{x}\right)$. 
   
   Now we consider $\mathbb{R}\oplus\mathbb{R}$, which is isomorphic to $\slashed{\mathbb{C}}$.  We have elements as $(x,y)$, product $(x,y)(z,w) = (xz,yw)$, $\mathbf{1}=(1,1)$, and $\|\mathbf{1}\|^2 = 2$.  Also $(x,y)^{-1} = (\frac{1}{x},\frac{1}{y})$.  We then find $\mathbf{u}_{\mathbb{R}\oplus\mathbb{R}} = \left[\begin{array}{ccc}
 \alpha & 0   \\ 
 0 & \gamma    
\end{array}\right]$ and $\log  \mathbf{U}_{\mathbb{R}\oplus\mathbb{R}} = \frac{\alpha}{2}\log x + \frac{\gamma}{2}\log y$. 



With the observations made at the end of Section 2.6 in mind, let us look further at the isomorphic algebras $\slashed{\mathbb{C}}$ and $\mathbb{R}\oplus\mathbb{R}$.  From the above calculations, it is clear that $\log\mathbf{U}_{\slashed{\mathbb{C}}}$ and $\log\mathbf{U}_{\mathbb{R}\oplus\mathbb{R}}$ are consistent based on the algebra isomorphism resulting from $\left[\begin{array}{ccc}
 1 & -1   \\ 
 1 & 1    
\end{array}\right]\left(\begin{array}{ccc}
x   \\ 
y   
\end{array}\right) = \left(\begin{array}{ccc}
x-y   \\ 
x+y   
\end{array}\right)$.  However, if we ignore that transformation, we could alternatively look at the respective sum of of functions given by $\log\mathbf{U}$ evaluated only on the respective coordinate axes of $\slashed{\mathbb{C}}$ and $\mathbb{R}\oplus\mathbb{R}$.  Based on the above calculations, $\slashed{\mathbb{C}}$ results in (\ref{7/29/24.1}) as $\alpha\log x+ \frac{\alpha}{2}\log[(1+
y)(1-y)] +  \frac{\beta}{2}\log \frac{1+y}{1-y}$ and $\mathbb{R}\oplus\mathbb{R}$ results in (\ref{7/29/24.1}) as $\frac{\alpha}{2}\log x +\frac{\beta}{2} \log y$. Consequently,  \[(\tau_{\slashed{\mathbb{C}},{\rm rat}},\tau_{\slashed{\mathbb{C}},{\rm log}},\tau_{\slashed{\mathbb{C}},{\rm arc}}) = (\tau_{\mathbb{R}\oplus\mathbb{R},{\rm rat}},\tau_{\mathbb{R}\oplus\mathbb{R},{\rm log}},\tau_{\mathbb{R}\oplus\mathbb{R},{\rm arc}}) = (0,2,0). \] Thus, the invariants in Corollary \ref{7/25/24.1} are at least consistent with $\slashed{\mathbb{C}}$ and $\mathbb{R}\oplus\mathbb{R}$ being isomorphic.  Similarly applying Corollary \ref{7/25/24.1}, it is easily seen that neither of $\slashed{\mathbb{C}}$ or $\mathbb{R}\oplus\mathbb{R}$ can be isomorphic to either $\mathbb{C}$ or $\mathbb{D}$ (nor can  $\mathbb{C}$ or $\mathbb{D}$ themselves be isomorphic).   So, even though searching for an isomorphism mapping for algebras that might be isomorphic is in general impractical, invariants derived from $\log\mathbf{U}$ are available to coarsely sort the algebras in a meaningful way.

As will be evident from Corollary 3.3, all unital norm groups have a one dimensional subgroup corresponding to multiples of what is closely allied with the usual algebra norm, and the parameter of $M_{\mathbf{u}_A}$ corresponding to this subgroup {\it always} makes a contribution to $\tau_{A,\log}$.  Thus, in distinguishing algebras, it makes sense to consider a ``reduced" triple of the invariants in Corollary \ref{7/25/24.1} as $(\tau_{A,{\rm rat}},\tau_{A,{\rm log}}-1,\tau_{A,{\rm arc}})$.  The three isomorphism classes of the two-dimensional algebras represented by $\mathbb{C}$, $\slashed{\mathbb{C}}$, and $\mathbb{D}$, respectively correspond to the reduced triples $(0,0,1)$, $(0,1,0)$, and $(1,0,0)$ (i.e., pure elemental forms of the reduced triples associated with algebras in general).
   
 The three isomorphism classes of the two-dimensional algebras are listed in \mbox{Table I} along with their respective uncurling metric groups and unital norm groups.  
 
 The six isomorphism classes of three-dimensional algebras are also instructive.  The structure constants for a representative of each of these are presented in \cite{kobayashi:2021}.  Table II lists those six algebras (in the order given in \cite{kobayashi:2021}) in matrix representations along with their respective uncurling metric groups and unital norm groups.
 
 \vspace{1cm}
 

  \hspace{1cm} The Three Isomorphism Classes of the Two-dimensional Algebras \enlargethispage{2cm}

\vspace{.5cm}


\hspace{-0cm}$ \begin{array}{cccccccc} 
\mbox{\underline{algebra}} & \mbox{\underline{uncurling metric group $(M_{\mathbf{u}})$}} & \mbox{\underline{unital norm group ($\mathbf{U}$)}}  \\ 
\\
\mathbb{C} \approx \left[\begin{array}{ccc}
x  & -y    \\
 y & x \\   
\end{array}\right]   & \left[\begin{array}{ccc}
\alpha  & \beta    \\
 \beta & -\alpha \\   
\end{array}\right]  & (x^2+y^2)^{\frac{\alpha}{2}} \,e^{\beta\,{\rm arctan}\left(\frac{y}{x}\right)}  \\ \\
\slashed{\mathbb{C}} \approx \left[\begin{array}{ccc}
x  & y    \\
 y & x \\   
\end{array}\right]   & \left[\begin{array}{ccc}
\alpha  & \beta    \\
 \beta & \alpha \\   
\end{array}\right]  & [(x+y)(x-y)]^{\frac{\alpha}{2}} \left[\frac{x+y}{x-y}\right]^{\frac{\beta}{2}}  \\ \\ 
\mathbb{R}\oplus\mathbb{R} \approx \left[\begin{array}{ccc}
x  & 0    \\
 0 & y \\   
\end{array}\right]   & \left[\begin{array}{ccc}
\alpha  & 0    \\
 0 & \beta \\   
\end{array}\right]  & x^\frac{\alpha}{2} y^\frac{\beta}{2}  \\ \\
\mathbb{D} \approx \left[\begin{array}{ccc}
x  & y    \\
 0 & x \\   
\end{array}\right]   & \left[\begin{array}{ccc}
\alpha  & \beta    \\
 \beta & 0 \\   
\end{array}\right]  & x^\alpha\,e^{\beta \frac{y}{x}}  \\ \\
\end{array} $

\vspace{.75cm}
 

\noindent {\bf Table I.} {\small Each row indicates a representative  of one of the three two-dimensional algebra isomorphism classes along with its left regular representation, followed by the associated uncurling metric group and unital norm group. Note that  $\slashed{\mathbb{C}}$ and $\mathbb{R}\oplus\mathbb{R}$ are isomorphic, their respective uncurling metric groups are compatible (via the algebra isomorphism $\left[\begin{array}{ccc}
 1 & 1   \\ 
 -1 & 1    
\end{array}\right]$), and their unital norm groups are  compatible (i.e., $(x+y)^{\frac{\alpha+\beta}{2}}(x-y)^{\frac{\alpha-\beta}{2}}$ versus $x^\frac{\alpha}{2} y^\frac{\beta}{2}$).  }


\newpage


  \hspace{1.5cm} The Six Isomorphism Classes of the Three-dimensional Algebras 

\vspace{1cm}


\hspace{-2.4cm}$ \begin{array}{cccccccc} 
\mbox{\underline{algebra}} & \mbox{\underline{uncurling metric group ($M_{\mathbf{u}}$) }} & \mbox{\underline{unital norm group ($\mathbf{U}$)}}  \\ 
\\
 \mathbb{R}\oplus \mathbb{R}\oplus \mathbb{R}  \approx \left[\begin{array}{ccc}
x  & 0 & 0   \\
 0 & y & 0 \\
 0 & 0 & z    
\end{array}\right]  \approx \mathbb{R}\oplus\slashed{\mathbb{C}}  &  \mathbb{R}\oplus M_{\mathbf{u}_{\mathbb{R}\oplus\mathbb{R}}}, \begin{array}{ccc}
V_\mathbb{R}(\alpha\beta\gamma)   \\ 
(3,3,1,3,2,3)   
\end{array} & x^{\frac{\alpha}{3}}y^{\frac{\beta}{3}}z^{\frac{\gamma}{3}} \\ \\
   \mathbb{R}\oplus\mathbb{C} \approx \left[\begin{array}{ccc}
x  & 0 & 0   \\
 0 & y & -z \\
 0 & z & y    
\end{array}\right]   \hfill &  \mathbb{R}\oplus M_{\mathbf{u}_{\mathbb{C}}}, \begin{array}{ccc}
V_\mathbb{R}(-\alpha(\beta^2+\gamma^2))   \\ 
(3,3,1,3,2,2)   
\end{array} &   x^\alpha (y^2+z^2)^{\frac{\beta}{2}} \,e^{\gamma\,{\rm arctan}\left(\frac{z}{y}\right)} \\ \\
   \mathbb{R}\oplus\mathbb{D} \approx \left[\begin{array}{ccc}
x  & 0 & 0   \\
 0 & y & z \\
 0 & 0 & y    
\end{array}\right]  \hfill &  \mathbb{R}\oplus M_{\mathbf{u}_{\mathbb{D}}}, \begin{array}{ccc}
V_\mathbb{R}(-\alpha\gamma^2)   \\ 
(3,3,1,2,1,2)   
\end{array}  &   x^\alpha y^{\beta} e^{\gamma\frac{y}{z}} \\ \\
   \mathbb{T}_3  \approx \left[\begin{array}{ccc}
x  & y & z   \\
 0 & x & y \\
 0 & 0 & x    
\end{array}\right] \hfill &  \left[\begin{array}{ccc}
\alpha  & \beta & \gamma   \\
 \beta & \gamma & 0 \\
 \gamma & 0 & 0   
\end{array}\right], \begin{array}{ccc}
V_\mathbb{R}(\gamma^3)   \\ 
(3,3,1,1,0,1)   
\end{array}  &  x^\alpha e^{ \beta\frac{y}{x}  +\gamma \frac{z}{x} -\gamma \frac{ 1}{2}\left(\frac{y}{x}\right)^2 } \\ \\
\left[\begin{array}{ccc}
x  & y & 0   \\
 0 & x & 0 \\
 0 & z & x    
\end{array}\right] \hfill &    \left[\begin{array}{ccc}
\alpha  & \beta & \gamma   \\
 \beta & 0 & 0 \\
 \gamma & 0 & 0   
\end{array}\right], \begin{array}{ccc}
V_\mathbb{R}(0)   \\ 
(3,2,1,0,3,0)   
\end{array}   &   x^\alpha e^{ \beta\frac{y}{x}  +\gamma \frac{z}{x}  } \\ \\
\left[\begin{array}{ccc}
x  & z & y   \\
 0 & x+y & 0 \\
 y & -z & x    
\end{array}\right]  \approx \left[\begin{array}{ccc}
x  & 0 & 0   \\
 0 & y & 0 \\
 0 & z & x    
\end{array}\right]  \hfill & \hspace{.5cm} \left[\begin{array}{ccc}
\alpha  & \beta & 0   \\
 \beta & \alpha & 0 \\
 0 & 0 & 0   
\end{array}\right] = M_{\mathbf{u}_{\slashed{\mathbb{C}}}}\oplus \{0\}, \begin{array}{ccc}
V_\mathbb{R}(0)   \\ 
(2,2,1,0,2,0)
\end{array}  &  (x+y)^\alpha (x-y)^\beta \\ \\
\end{array} $



\vspace{1cm}


\noindent {\bf Table II.} {\small Each row lists a representative  of one of the six three-dimensional algebra isomorphism classes along with its representation as a subalgebra of $M_3(\mathbb{R})$, followed by the associated uncurling metric group and unital norm group (``$\approx$" means \mbox{``isomorphic to"}).  $\mathbb{T}_3$ is the three-dimensional upper triangular Toeplitz matrix algebra. The invariant ordered sextuple in each row lists the value of the dimension of $\mathbf{u}$, the largest rank of a member of $\mathbf{u}$, the smallest rank of a nonzero member of $\mathbf{u}$, the dimension of the space to which det$(M_{\mathbf{u}})$ is sensitive, the real dimension of the real algebraic set defined by det$(M_{\mathbf{u}}) =0$, and the number of irreducible components of $V_\mathbb{R}(\text{det}(M_{\mathbf{u}}))$ (this hypersurface is indicated above the sextuple).}
 
\newpage

\subsection{Special unital norms} \enlargethispage{1cm}
On an open ball centered at $\mathbf{1}$ devoid of non-units but containing the line segment from $\left(\frac{1}{\alpha}\right)\mathbf{1}$ to $\alpha\mathbf{1}$ for some choice of $\alpha>0$ (such a neighborhood always exists if $\alpha$ is close enough to 1), it is easy to show from (\ref{11/7/20.2}) that for unital norms in general,  \[  \ell(\alpha s) = \alpha^{\frac{\mathbf{1}^T L \mathbf{1}}{\|\mathbf{1}\|^2}} \ell(s),\] using an argument similar to that employed in the proof of Theorem \ref{9/24/23.1}. Thus, $ \ell(s)$ is a positive homogeneous function of degree given by the exponent on the right-hand-side of the above equation.  
   {\it But things become somewhat nicer if elements that are inverses with respect to the algebra's product also behave inversely with respect to the scalar product defined by $L$}.  
    This situation results if uncurling metric $L$ satisfies the additional constraint (\ref{7/7/21.1}).

    \begin{definition}\label{9/20/23.1} A {\it normalized uncurling metric} $L_{\shortparallel}$ is an uncurling metric that satisfies, \begin{equation} \label{7/7/21.1} s^TL_{\shortparallel} s^{-1} =\|{\bf 1}\|^2,\end{equation} on a sufficiently small neighborhood of $\mathbf{1}$.   A {\it special unital norm} $\ell_{\rm sp}(s)$ is a unital norm that is derived from a normalized uncurling metric $L_{\shortparallel}$ according to (\ref{11/7/20.2}). 
\end{definition}


\begin{theorem}\label{9/24/23.1} $ \ell_{\rm sp}(s)$ is a degree-1 positive homogeneous function, and
\begin{equation}\label{5/23/23.3}  \ell_{\rm sp}(s^{-1}) = \big(\ell_{\rm sp}(s)\big)^{-1} = \frac{1}{\ell_{\rm sp}(s)}, \end{equation} when $s,s^{-1}$ are in some open ball centered at $\mathbf{1}$ devoid of non-units. \end{theorem}

\begin{proof} We have already observed (shortly following Definition \ref{7/21/23.1}) that an open ball centered at $\mathbf{1}$ devoid of non-units always exists.  Let $\eta$  be such a neighborhood.

Given $\frac{1}{\alpha}s,s,\alpha s\in\eta$ for some choice of $\alpha>0$, we have from (\ref{11/7/20.2}) that
\begin{eqnarray}\nonumber  \log\ell_{\rm sp}(\alpha s) &=&  \frac{1}{\|{\bf 1}\|^2}\int_{\bf 1}^{\alpha s} [L_{\shortparallel}t^{-1}]\cdot {\sf d}t = \frac{1}{\|{\bf 1}\|^2}\int_{\frac{\mathbf{1}}{\alpha}\mathbf{1}}^{s} [L_{\shortparallel}(\alpha w)^{-1}]\cdot d(\alpha w)  \\ \nonumber &=&  \frac{1}{\|{\bf 1}\|^2}\int_{\mathbf{1}}^{s} [L_{\shortparallel}w^{-1}]\cdot dw + \frac{1}{\|{\bf 1}\|^2}\int_{\frac{\mathbf{1}}{\alpha} \mathbf{1}}^{\mathbf{1}} [L_{\shortparallel}w^{-1}]\cdot dw  \\ &=& \log \ell_{\rm sp}(s) + \log\alpha, \label{4/2/23.1}
\end{eqnarray} where the integral from $\frac{\mathbf{1}}{\alpha} \mathbf{1}$ to $\mathbf{1}$ can be easily evaluated along the line segment with those endpoints, and we have used (\ref{7/7/21.1}) in evaluation of that integral to obtain the term $\log \alpha$ on the right-hand-side of the final equation.  Thus, we have $\ell_{\rm sp}(\alpha s) = \alpha\ell_{\rm sp}(s)$, degree-1 positive homogeneity.

We 
now consider a smooth path $\mathbf{P}$ from $\mathbf{1}$ to $s$ which is inside $\eta$  such that $\mathbf{P}^{-1}$ is also inside $\eta$,  where $\mathbf{P}^{-1}$ is the point set consisting of the multiplicative inverses of the members of $\mathbf{P}$.  The  existence of neighborhood $\eta$ with such a property  is \enlargethispage{1.2cm} easily demonstrated using the left regular representation of the algebra.
\mbox{We then have,} \begin{eqnarray}  \log \ell_{\rm sp}(s^{-1})   &=&   \frac{1}{\|{\bf 1}\|^2} \int_{\bf 1}^{s^{-1}} [L_{\shortparallel}t^{-1}] \cdot {\sf d}t \nonumber \\ &=&  \frac{1}{\|{\bf 1}\|^2} \int_{\bf 1}^s [L_{\shortparallel}y]\cdot d\left(y^{-1}\right) \nonumber \\ &=&  \frac{1}{\|{\bf 1}\|^2}\left(\left(y^{-1}\cdot [L_{\shortparallel}y]\right)\bigg|_{\bf 1}^s - \int_{\bf 1}^s y^{-1}\cdot d\left(L_{\shortparallel}y\right)\right)   \nonumber \\ &=& -\frac{1}{\|{\bf 1}\|^2} \int_{\bf 1}^s y^{-1}\cdot (L_{\shortparallel}dy) = -\frac{1}{\|{\bf 1}\|^2} \int_{\bf 1}^s [L_{\shortparallel}y^{-1}] \cdot dy  \nonumber \\ &=& -\log \ell_{\rm sp}(s),\label{4/2/23.5} \end{eqnarray} where we have used a change of variable (contemplating the left regular representation, it is appreciated that $y=t^{-1}$ is a diffeomorphism on $\eta$), commutativity of a linear transformation and a differential, the property that $L_{\shortparallel}$ is a real symmetric matrix, and  (\ref{7/7/21.1}).  Equation (\ref{5/23/23.3}) then follows.
\end{proof}



As demonstrated in Section 3 (Corollary \ref{11/17/22.1}), every algebra is associated with at least one special unital norm, which ultimately follows from the fact that the transpose operation is a normalized uncurling metric for $M_n(\mathbb{R})$.

It is also seen in Section 3 that at least some members of an algebra's group of unital norms are multiplicative, i.e., $\ell(s_1s_2) = \ell(s_1)\ell(s_2)$ (in fact, all of the unital norms associated with the two-dimensional algebras are multiplicative).   But in general, unital norms are not multiplicative.  However, the unital norms originate from the nonlinear unary multiplicative inverse operation $s\mapsto s^{-1}$ implied by the algebra's bilinear binary product operation - so the particularly desirable fealty would be to that unary operation.  It is thus significant that the special unital norms pay homage to it via (\ref{5/23/23.3}) (also, see comments in the second and third paragraphs of Section 4.2).

Two reasons we have employed the term {\it unital} in ``unital norm" are that the latter entities apply to {\it unital} algebras, and they {\it unite} metric and rotor information.  One reason we have employed the term ``norm" for these entities is that the group includes a member closely related to the usual algebra norm (Corollary \ref{11/17/22.1}).  A second reason is as follows.   Using (\ref{11/7/20.1}), (\ref{5/23/23.3}), the degree-1 positive homogeneity property of $\ell_{\rm sp}(s)$, and assuming the associated uncurling metric $L_{\shortparallel}$ is nonsingular, it is appreciated that a special unital norm has the feature that,
\begin{equation}\label{5/22/24.5}  s = \ell_{\rm sp}(s)\big[\|\mathbf{1}\|^2 L_{\shortparallel}^{-1}[\nabla\ell_{\rm sp}(s^{-1})]\big],\end{equation} with
\begin{equation}\label{5/22/24.6}\ell_{\rm sp}\big(\|\mathbf{1}\|^2 L_{\shortparallel}^{-1}[\nabla\ell_{\rm sp}(s^{-1})]\big) = 1,\,\end{equation}
where $\nabla\ell_{\rm sp}(s^{-1})$ means that $\nabla\ell_{\rm sp}$ is evaluated at $s^{-1}$.  What these equations say is that (as a member of the vector space of elements) an element $s$ is the product of its ``norm" $\ell_{\rm sp}(s)$ and a ``unit norm" vector that is determined by $\ell_{\rm sp}$ (\mbox{a ``direction"}).  This format is very reminiscent of equations satisfied by the norm on an inner product space (see the latter parts of Appendix A and Appendix B for more context).

\subsection{Categorics and the epimorphism problem}
 In line with Remark \ref{3/31/24.1}, our overall program up to this point can be cast in categorical language, though that is hardly necessary, apart from the organizational insight it provides. 
On the other hand, this program can also more generally address algebra epimorphisms, and as regards direct use of uncurling spaces for that purpose, this is most conveniently introduced by a particular functor as follows (though we will not explicitly develop it here, the epimorphism problem can alternatively be addressed via the invariants identified in Corollary \ref{7/25/24.1}, and examples of that are found in Section 5).  

Let $\mathcal{E}$ be the category each of whose objects is an algebra (i.e., a real finite-dimensional unital associative algebra), and whose morphisms are now algebra epimorphisms (we will use the same symbol, e.g., $K$, for both the epimorphism and morphism).  We define category $\mathcal{P}$ as follows. An object in this category is a symmetric $(n\times n)$-matrix where $n$ is any particular positive integer, and where each matrix entry is a linear combination of the components of an $m$-vector of real parameters, $m\ge 1$ (i.e., the $m$ components of the later vector are each a different real parameter).  Thus, for any object of $\mathcal{P}$, the implied set of real matrices resulting from all realizations of the associated $m$-vector of parameters is an $m$-dimensional space (a realization of the $m$-vector of parameters is a real $m$-vector resulting from a choice of real values for the $m$ parameters).  An object of $\mathcal{P}$ is considered to be the same as the latter space of real matrices.  Consider two objects ${\sf P}_1,{\sf P}_2$ of $\mathcal{P}$, where ${\sf P}_1$ is a $(j\times j)$-matrix parametrized by the components of a $k$-parameter vector, and ${\sf P}_2$ is a $(q\times q)$-matrix parametrized by the components of a $r$-parameter vector.  There is a morphism $ \begin{CD}  {\sf P}_1 @>p_1>> {\sf P}_2\end{CD}$ if and only if  there exists a real $(q\times j)$-matrix $W$ such that $WW^T$ is nonsingular and $W^T{\sf P}_2W\subset {\sf P}_1$, where the latter subset expression is 
interpreted to mean that the set of real matrices implied by $W^T{\sf P}_2W$ through all realizations of the $r$-parameter vector associated with ${\sf P}_2$, is a subset of  the set of real matrices implied by  ${\sf P}_1$ through all realizations of the associated $k$-parameter vector.  It is evident that there is at most one morphism pointing from one category $\mathcal{P}$ object to another.  It is also easily seen that the morphisms are associative.  

We define a mapping $\mathbf{F}$ from category $\mathcal{E}$ to category $\mathcal{P}$ as follows.  $\mathbf{F}[
A] \equiv \mathbf{u}_A$, and given a morphism $K$ from $A_1$ to $A_2$, we define $\mathbf{F}(K)$ to be the morphism from $\mathbf{F}[A_1]$ to $\mathbf{F}[A_2]$, if the morphism exists.  
 


\begin{theorem}\label{6/13/22.1}  ${\bf F}: \mathcal{E}\rightarrow \mathcal{P}$, is a covariant functor. 
 \end{theorem}
 \begin{proof}
We will first show that the category $\mathcal{E}$ morphism $ \begin{CD}  A_1 @>K>> A_2\end{CD}$ implies the existence of category $\mathcal{P}$ morphism $ \begin{CD}  {\bf F}[A_1] @>{\bf F}(K)>> {\bf F}[A_2]\end{CD}$.

Morphism $K$ results from an algebra epimorphism $K:A_1\rightarrow A_2$.   Taking $s_1$ to be notation for a variable in $A_1$, we have ${\bf 1}_{A_2} = K(s_1s_1^{-1}) = K(s_1)K(s_1^{-1})$, so that $(Ks_1)^{-1} = K(s_1^{-1})$.  By definition, given the uncurling space $\mathbf{F}[A_2]$ associated with $A_2$ and an arbitrary real matrix  $\hat{\mathbf{F}}[A_2]$ in the set of real matrices implied by $\mathbf{F}[A_2]$, we have $d\left((s_2^{-1})^T\hat{\mathbf{F}}[A_2] ({\sf d}s_2)\right) = 0$ in some neighborhood of ${\bf 1}_{A_2}$, where $s_2$ is notation for a variable in $A_2$.   With respect to appropriate neighborhoods of $\mathbf{1}_{A_1}$ and $\mathbf{1}_{A_2}$, it follows that \begin{eqnarray}\nonumber 0 = d\left((s_2^{-1})^T\hat{\mathbf{F}}[A_2] {\sf d}s_2\right) &=&  d\left(((Ks_1)^{-1})^T\hat{\mathbf{F}}[A_2] d(Ks_1)\right) \\  &=& d\left((s_1^{-1})^T \big[K^T \hat{\mathbf{F}}[A_2]K\big]{\sf d}s_1\right), \label{7/4/22.5}  \end{eqnarray}   using the equation at the end of the second sentence of this paragraph, and the fact that a differential commutes with a linear transformation.  Equation (\ref{7/4/22.5}) is simply an instance of (\ref{4/7/23.1}), so it follows that symmetric $K^T \hat{\mathbf{F}}[A_2]K$ is an uncurling metric associated with $A_1$.  Consequently, $K^T\mathbf{F}[A_2]K\subset \mathbf{F}[A_1]$.  Furthermore, $KK^T$ is nonsingular since $K$ is an epimorphism.  Hence, criteria for the existence of morphism $\mathbf{F}(K)$ are fulfilled (with $W=K)$.

Commutativity of the requisite diagram 
  \[  \begin{CD}  
A_1 @>K_1>> A_2 @>K_2>> A_3\\
 @VVV   @VVV  @VVV\\
 {\bf F}[A_1]   @>{\bf F}(K_1)>>  {\bf F}[A_2]  @>{\bf F}(K_2)>>  {\bf F}[A_3] \\ 
 \end{CD} \]
is then evident from existence and uniqueness of the relevant category $\mathcal{P}$ morphisms, and the mechanism of their composition.  Thus, ${\bf F}$ is a covariant functor. 
\end{proof}

We then have the following useful results.  The first is immediate.
\begin{corollary}\label{6/26/22.2} An epimorphism from one algebra to another algebra can exist only if there is a category $\mathcal{P}$ morphism associated with their respective uncurling spaces. \end{corollary}

\begin{corollary} \label{8/17/22.5} An epimorphism from a first algebra to a second algebra can exist only if the dimension of the uncurling space of the first algebra is greater than or equal to the dimension of the uncurling space of the second algebra. \end{corollary}  

\begin{proof} Given algebra epimorphism $K:A_1\rightarrow A_2$, suppose $A_1$ is a $j$-dimensional algebra with an uncurling space utilizing a $k$-parameter vector (and thereby the uncurling space is a $k$-dimensional space), and $A_2$ is a $q$-dimensional algebra with an uncurling space utilizing a $r$-parameter vector (and thereby the uncurling space is a $r$-dimensional space).  Evidently, $K$ is a $(q\times j)$-matrix with $q\le j$ and $KK^T$ is nonsingular (since $K$ is an epimorphism). The proof of Theorem \ref{6/13/22.1} has shown that $K^T \mathbf{F}[A_2] K \subset \mathbf{F}[A_1]$.  Thus, $K^T \mathbf{F}[A_2] K$ cannot have more parameters than $\mathbf{F}[A_1]$  (i.e., the dimension of $K^T \mathbf{F}[A_2] K$ cannot exceed the dimension of $\mathbf{F}[A_1]$).  Furthermore, $K^T \mathbf{F}[A_2] K$ has at least as many parameters as $KK^T \mathbf{F}[A_2] KK^T$, but it has less than or equal to the number of parameters of $\mathbf{F}[A_2]$.  However, $KK^T \mathbf{F}[A_2] KK^T$ has the same number of parameters as $\mathbf{F}[A_2]$ since $KK^T$ is nonsingular.  Thus, $K^T \mathbf{F}[A_2] K$ has the same number of parameters as $\mathbf{F}[A_2]$.  Hence, the dimension of $\mathbf{F}[A_2]$ cannot exceed the dimension of $\mathbf{F}[A_1]$.
\end{proof}

\begin{corollary}\label{8/17/22.1} If the uncurling space of an algebra is one-dimensional and contains a nonsingular member, then the algebra is simple. \end{corollary}

\begin{proof} Suppose algebra $A_1$ is $j$-dimensional and has a non-trivial ideal, and its uncurling space is one-dimensional and has a nonsingular member.  It follows that $\mathbf{F}[A_1]$ is represented by  $\alpha L_1$, $\alpha\in \mathbb{R}$, where $L_1$ is a real nonsingular $(j\times j)$-matrix.  Since $A_1$ has a non-trivial ideal, there exists an algebra epimorphism $K:A_1\rightarrow A_2$ where $A_2$ has dimension $q<j$.  Corollary \ref{8/17/22.5} indicates that the uncurling space of $A_2$ cannot have dimension greater than one, and thus $\mathbf{F}[A_2]$ is represented by $\alpha L_2$ where $L_2$ is a real $(q\times q)$-matrix with $q<j$. From the proof of Theorem \ref{6/13/22.1}, we must then have that the nonzero matrix $ K^T L_2 K$ is a member of the uncurling space $\mathbf{F}[A_1]$.  But the determinant of the nonzero matrix $ K^T L_2 K$ is zero since $K$ has a nontrivial kernel, while the only member of the uncurling space $\mathbf{F}[A_1]$ having zero determinant is the zero matrix (corresponding to $\alpha=0$), since $L_1$ is necessarily nonsingular.   It follows that $K$ cannot exist, and so $A_1$ is simple. \end{proof}


In view of Theorem \ref{4/8/22.1}, an easy consequence of Corollary \ref{8/17/22.1} is an alternative demonstration that the algebra of real $(n\times n)$-matrices $ M_n(\mathbb{R})$ is simple.  

Note that the converse of Corollary \ref{8/17/22.1} is not true.  For example, $\mathbb{C}$ is simple but has a two-dimensional uncurling space (see Section 2.7).


\section{The unital norm group generalizes the usual algebra norm} 

The ``usual" algebra norm is given by the determinant of an element's image under the left regular representation (e.g., \cite{bourbaki:1989}).  The group $\mathbf{U}$ of unital norms subsumes this definition in the sense of the theorems and corollaries of this section.  This also means that the uncurling space of an algebra is at least one-dimensional.

 We begin with the matrix algebra $ M_n(\mathbb{R})$. For an invertible element $s$, with det$(s)$ being its determinant, we have the well-known (though in non-standard format),
  \begin{equation}\label{11/5/21.1} \nabla \mbox{det}(s) = \big[{\rm Adj}(s)\big]^T = \mbox{det}(s)[s^{-1}]^T,\end{equation} where Adj$(s)$ is the adjugate matrix.  The first of these equalities is special case of Jacobi's formula, and the second follows from the classical matrix inverse expression resulting from Laplace expansion.  The gradient notation on the left-hand-side is employed with the understanding that an element of the algebra (including the full expression on the right-hand-side), as a member of the vector space of elements, is identified as the $n^2$-tuple obtained from sequentially appending the columns of its matrix format (and vice-versa), with a compatible understanding regarding the transpose operation as a particular involutive $(n^2\times n^2)$-matrix, which we will denote as $\mathcal{T}$. 
  
From (\ref{11/5/21.1}), $\mathcal{T}s^{-1} = \frac{\nabla{\rm det}(s)}{{\rm det}(s)} = \nabla[\log{\rm det}(s)]$, so $\mathcal{T}$ is an uncurling metric.  Since ${\rm det}(s)$ is a degree-$n$ positive homogeneous function, the Euler Homogeneous Function Theorem implies $[\nabla{\rm det}(s)]\cdot s = n \,{\rm det}(s)$.  We also have $\|\mathbf{1}\|^2 = n$.  Consequently, (\ref{7/7/21.1}) is satisfied for $L=\mathcal{T}$, so $\mathcal{T}$ is a normalized uncurling metric.  Applying the first equation of this paragraph to (\ref{11/7/20.2}), it is then seen that $(\text{det}(s))^{\frac{1}{n}}$ is a special unital norm.
  Moreover,
 
 
 \begin{theorem} \label{4/8/22.1} $\mathbf{u}_{M_n(\mathbb{R})}$ is one-dimensional and there is only a single normalized uncurling metric, this being the  transpose operation on $(n\times n)$-matrices expressed as an $(n^2\times n^2)$-matrix. The associated special unital norm is $({\rm det}(s))^{\frac{1}{n}}$, i.e., the usual norm raised to the power given by the reciprocal of the dimension of $M_n(\mathbb{R})$. \end{theorem}
 \begin{proof} The vector space of members of $ M_n(\mathbb{R})$ consists of  vectors with $n^2$ components.   Using the sequentially appended matrix columns convention as introduced in the second paragraph of this section, we can equivalently represent a member $s\in M_n(\mathbb{R})$ in either matrix or column vector format.   An uncurling metric for this algebra will be a matrix $L\in M_{n^2}(\mathbb{R})$ that must satisfy (\ref{4/7/23.1}), i.e., $d\left([Ls^{-1}]\cdot {\sf d}s\right)=0$. This constitutes a condition on every pair of distinct components of $Ls^{-1}$ (the latter again having $n^2$ components).  Thus,  the partial derivative of the  $(i,j)$-entry of matrix $Ls^{-1}$ (say), taken with respect to the $(k,m)$-entry of matrix $s$ (say) (with $(k,m)\ne (i,j)$),  minus  the partial derivative of the $(k,m)$-entry of $Ls^{-1}$  taken with respect to the $(i,j)$-entry of matrix $s$, must be zero.
 
 To examine this constraint, let $s_{pq}$ be the $(p,q)$-entry of matrix $s$.   Differentiating both sides of the expression
$ ss^{-1} = {\bf 1}$ with respect to $s_{pq}$ (as in Matrix Calculus), we obtain $\frac{\partial (s^{-1}s)}{\partial s_{pq}} = \frac{\partial (s^{-1})}{\partial s_{pq}}s + s^{-1}\frac{\partial s}{\partial s_{pq}}=0$ (where $0$ is the matrix of all zero entries).  Thus, we have the well-known result 
 \[ \frac{\partial (s^{-1})}{\partial s_{pq}} = -s^{-1}\frac{\partial s}{\partial s_{pq}}s^{-1} = -s^{-1}e_{pq}s^{-1},\]   where $e_{pq}$ is the matrix whose $(p,q)$ entry is $1$, but with all other entries zero (and $\{ e_{pq}\}$ is the canonical basis for $ M_n(\mathbb{R})$).  
 Hence, the partial derivative of $s^{-1}$ with respect to the  $(k,m)$-entry of matrix $s$, which is itself a matrix, (``the derivative of $s^{-1}$ in the $e_{km}$-direction") is, \begin{equation}\label{8/18/22.5} -s^{-1} e_{km}s^{-1} = -(s^{-1})_{:k} (s^{-1})_{m:},\end{equation} where $(s^{-1})_{:k}$ indicates the $k$-th column of $s^{-1}$, and $(s^{-1})_{m:}$ indicates the $m$-th row of $s^{-1}$.    Similarly, the derivative of $s^{-1}$ with respect to the $(i,j)$-entry of $s$ (``the derivative of $s^{-1}$ in the $e_{ij}$-direction")  is, \begin{equation}\label{8/18/22.6} -s^{-1} e_{ij}s^{-1} = -(s^{-1})_{:i} (s^{-1})_{j:}.\end{equation}  It is thus clear that the sets of entries of the $e_{km}$-direction derivative matrix and the $e_{ij}$-direction derivative  matrix share the member  $-(s^{-1})_{mi} (s^{-1})_{jk}$, and they can share no other member because $\{(s^{-1})_{pq}\} = \{\frac{C_{qp}}{\det(s)}\}$, where the right-hand-side of this equation is clearly a set of linearly independent functions since $\{C_{qp}\}$ is the matrix of cofactors of $s$.  The latter observation also indicates that the union of the entries of the $e_{km}$-direction derivative matrix with the entries of the $e_{ij}$-direction derivative matrix is a set of linearly independent functions.

Therefore, since a derivative and a linear transformation commute, the partial derivative in the $e_{km}$-direction of the $e_{ij}$-component of $Ls^{-1}$ is a linear combination of the $n^2$ entries of the matrix on the right-hand-side of (\ref{8/18/22.5}), and  the partial derivative in the $e_{ij}$-direction of the $e_{km}$-component of $Ls^{-1}$ is a linear combination of the $n^2$ entries on the right-hand-side of (\ref{8/18/22.6}).
So satisfaction of \mbox{$d\left([Ls^{-1}]\cdot {\sf d}s\right) = 0$,} based on the observation made at the end of the prior paragraph, requires that all entries in any two rows of $L$ must be zero, except for the one entry in each row corresponding to the shared member of the two relevant partial derivative matrices (the right-hand-sides of (\ref{8/18/22.5}) and (\ref{8/18/22.6})) noted in the prior paragraph - and these two corresponding nonzero entries of the two cited rows of $L$ must have the same value.  This prescription must be true for any two rows of $L$, so all nonzero entries of $L$ must have the same value, and there can be at most one nonzero value in each row of $L$.  Thus, the uncurling space is at most one-dimensional. Assuming the group is one-dimensional, the uncurling space can be represented as a parameter $\alpha$ ranging over $\mathbb{R}$ multiplying a nonzero constant matrix whose nonzero entries are all 1.  
For the special unital norm, the value of $\alpha$ is determined by the requirement that (\ref{7/7/21.1}) be satisfied. 


Hence, there is at most one possible normalized uncurling metric.  But according to observations made in the third paragraph of this section, $ M_n(\mathbb{R})$ does indeed have a normalized uncurling metric - this being the $(n^2\times n^2)$-matrix corresponding to the matrix transpose operation on $(n\times n)$-matrices.  As was shown immediately prior to the theorem statement, $(\mbox{det}(s))^{\frac{1}{n}}$ is the associated special unital norm.  
 \end{proof}
 
 Expanding on Theorem \ref{4/8/22.1}, the following theorem and corollary indicate the manner in which unital norms generalize the usual algebra norm.  Recall that a subalgebra of $M_n(\mathbb{R})$ is defined to be self-adjoint if the transpose operation maps the subalgebra to itself.
 
  \begin{theorem}\label{5/13/22.1} If the left regular representation of an $n$-dimensional algebra is self-adjoint as a subalgebra of $M_n(\mathbb{R})$, then there is a special unital norm given by the usual algebra norm raised to the power of $\frac{1}{n}$.
\end{theorem}

\begin{proof} 
Let $\bar{A}$ be an $n$-dimensional algebra and let $A$ be its left regular representation.  Given $\bar{s}\in\bar{A}$ we denote its image under the left regular representation as $s$.
 So, $\bar{A}$ is isomorphic to $A$, which is a subalgebra of $ M_n(\mathbb{R})$.   $\mathbf{1}_{M_n(\mathbb{R})}$ is the $(n\times n)$-identity matrix, and observe that as an element of $ M_n(\mathbb{R})$, $\mathbf{1}_A$ is likewise the $(n\times n)$-identity matrix.  However, $ M_n(\mathbb{R})$ is $n^2$-dimensional with $\|\mathbf{1}_{M_n(\mathbb{R})}\|^2 = n$, while $A$ is $n$-dimensional with $\|\mathbf{1}_A\|^2 = \|\mathbf{1}_{\bar{A}}\|^2$.  That is, as {\it vector spaces}, $\bar{A}$ is mapped to  $A$ by the identity operation, and the latter are subject to the Euclidean inner product as per the conventions stated in the Section 1.   

As noted shortly prior to the statement of Theorem \ref{4/8/22.1}, the transpose operation on $ M_n(\mathbb{R})$ is a linear transformation $\mathcal{\mathcal{T}}:\mathbb{R}^{n^2}\rightarrow \mathbb{R}^{n^2}$ and, also notating it as $\mathcal{\mathcal{T}}$ with respect to the standard basis, matrix $\mathcal{\mathcal{T}}$ represents a normalized uncurling metric associated with algebra  $ M_n(\mathbb{R})$.  Being an uncurling metric means that integration of $\mathcal{\mathcal{T}}s^{-1}$ is path-independent in some simply connected neighborhood of the multiplicative identity of $ M_n(\mathbb{R})$.  
 Hence, $\int_{\mathbf{1}_{M_n({\mathbb{R})}}}^s (\mathcal{\mathcal{T}}t^{-1})\cdot {\sf d}t$ is path-independent in the latter neighborhood and, in particular, is path-independent for paths in that neighborhood restricted to be in $A$.  By assumption, $\mathcal{\mathcal{T}}$ maps $A$ to $A$. 
 Thus, $\int_\mathbf{1}^s \mathcal{\mathcal{T}}t^{-1}\cdot {\sf d}t$ is path-independent when the domain of integration is restricted to $A$.  This means that 
  $d((\mathcal{\mathcal{T}}s^{-1})\cdot {\sf d}s) = 0$ in some neighborhood of $\mathbf{1}_A$ in $A$.
  
   Now, we can use $\mathcal{\mathcal{T}}$ to define a scalar product on $A$ (not necessarily positive-definite or nondegenerate) as $\langle s_1,s_2\rangle_A \equiv  \frac{\|\mathbf{1}_{A}\|^2}{n}s_1^T \mathcal{\mathcal{T}} s_2$. 
    Since $A$ and $\bar{A}$ are isomorphic, the above scalar product on $A$ can be employed to produce a scalar product on $\bar{A}$ so that the two resulting quadratic spaces are isomorphic.  Specifically,  for $s_1,s_2\in A$, let $ \mathcal{I}:\bar{A}\rightarrow A$ be the left regular representation isomorphism, with $s_1 = \mathcal{I}\bar{s}_1$ and $s_2= \mathcal{I}\bar{s}_2 $.  Then we define the induced scalar product on $\bar{A}$ via, \begin{equation}\label{7/15/22.1} \bar{\mathcal{\mathcal{T}}} \equiv  \frac{\|\mathbf{1}_{\bar{A}}\|^2}{n}\mathcal{I}^T\mathcal{\mathcal{T}}\mathcal{I},\end{equation} i.e., $\langle \bar{s}_1,\bar{s}_2\rangle_{\bar{A}} \equiv \bar{s}_1^T \bar{\mathcal{\mathcal{T}}} \bar{s}_2$.  Applying (\ref{7/15/22.1}) and the previously noted $\|\mathbf{1}_A\|^2 = \|\mathbf{1}_{\bar{A}}\|^2$, we have \begin{equation}\label{8/25/22.1} \langle \bar{s}_1,\bar{s}_2\rangle_{\bar{A}} = \bar{s}_1^T \bar{\mathcal{T}} \bar{s}_2 = \frac{\|\mathbf{1}_{\bar{A}}\|^2}{n}( s_1^T \mathcal{T} s_2) = \langle s_1,s_2\rangle_{A}.\end{equation}  The fraction $\frac{\|\mathbf{1}_{\bar{A}}\|^2}{n}$ is employed in (\ref{7/15/22.1}) in order that $\langle \mathbf{1}_{\bar{A}},\mathbf{1}_{\bar{A}}\rangle_{\bar{A}} = \|\mathbf{1}_{\bar{A}}\|^2$, given that $\mathbf{1}_{M_n(\mathbb{R})}^T \mathcal{T}\,\mathbf{1}_{M_n(\mathbb{R})} = n = \|\mathbf{1}_{M_n(\mathbb{R})}\|^2$. Due to this (isometric) isomorphism of $A$ and $\bar{A}$ as quadratic spaces,  in small neighborhoods of the respective multiplicative identities $\mathbf{1}_A$ and $\mathbf{1}_{\bar{A}}$, the previously demonstrated $d((\mathcal{T}s^{-1})\cdot {\sf d}s) = 0$ implies 
\mbox{$d((\bar{\mathcal{T}}(\bar{s})^{-1})\cdot d\bar{s}) = 0$.} Thus, $\bar{\mathcal{T}}$ is an uncurling metric for $\bar{A}$.


In the context of $M_n(\mathbb{R})$, we know from Theorem \ref{4/8/22.1} that $L=\mathcal{T}$ satisfies (\ref{7/7/21.1}), where $\|\mathbf{1}\|^2$ in  (\ref{7/7/21.1}) is $\|\mathbf{1}_{M_n({\mathbb{R})}}\|^2 = n$.  Combining these facts with (\ref{8/25/22.1}), we have \begin{equation}\label{9/27/23.1} (\bar{s})^T \bar{\mathcal{T}} (\bar{s})^{-1} = \frac{\|\mathbf{1}_{\bar{A}}\|^2}{n}( s^T \mathcal{T} s^{-1}) = \|\mathbf{1}_{\bar{A}}\|^2 .\end{equation} 
That is, $\mathcal{\bar{T}}$ is a normalized uncurling metric for $\bar{A}$.  Hence, a special unital norm on $\bar{A}$ results from (\ref{11/7/20.2}) (with $\mathbf{1}=\mathbf{1}_{\bar{A}}$)  via the normalized uncurling metric $L_{\shortparallel} = \bar{\mathcal{T}}$.  
  We then obtain, 
\begin{eqnarray}\nonumber \log\ell_{\rm sp}(\bar{s}) = \frac{1}{\|\mathbf{1}_{\bar{A}}\|^2} \int_{\mathbf{1}_{\bar{A}}}^{\bar{s}} (\bar{\mathcal{T}}(\bar{s})^{-1})\cdot d\bar{s}
 &=& \frac{1}{\|\mathbf{1}_{\bar{A}}\|^2} \int_{\mathbf{1}_{A}}^{s} \frac{\|\mathbf{1}_{\bar{A}}\|^2}{n}(\mathcal{T}s^{-1})\cdot {\sf d}s \\ &=& \frac{1}{n} \log\text{det}(s),\label{9/27/23.2}\end{eqnarray} where the first equality is from (\ref{11/7/20.2}), the second equality follows from application of (\ref{7/15/22.1}), and the final equation results from Theorem \ref{4/8/22.1}.
The theorem then follows. 
 \end{proof} 
 

So now we consider the general case.
 
  \begin{corollary}\label{11/17/22.1}
  Let $m$-dimensional algebra $\underline{A}$ be such that there is an epimorphism $\Omega:\bar{A}\rightarrow \underline{A}$, and such that the left regular representation of $\underline{A}$ is self-adjoint as a subalgebra of $M_m(\mathbb{R})$. There is at least one such $\underline{A}$ for which $m\ge 1$, and for each one there is a member of the special unital norm family of  $\bar{A}$ given by $[{\rm det}(\underline{\mathcal{I}}\Omega \bar{s})]^{\frac{1}{m}}$, for $\bar{s}\in\bar{A}$, where $\underline{\mathcal{I}}$ is the isomorphism from $\underline{A}$ to its left regular representation.
  \end{corollary}
 
  \begin{proof}   There exists an epimorphism $\Omega_1$ mapping $\bar{A}$  to its one-dimensional subalgebra $\{\alpha\mathbf{1}_{\bar{A}}:\alpha\in\mathbb{R}\}$, resulting from the improper ideal.   It is evident that the left regular representation of  $\{\alpha\mathbf{1}_{\bar{A}}:\alpha\in\mathbb{R}\}$ is self-adjoint as a subalgebra of $M_1(\mathbb{R})$. Thus, there is at least one subalgebra of $\bar{A}$ satisfying the theorem hypotheses and such that $m\ge 1$.  
  
  By Theorem \ref{5/13/22.1}, there exists a normalized uncurling metric $\underline{L}_{\shortparallel}$ for $\underline{A}$ with the associated special unital norm as $\ell_{\rm sp}(\underline{s}) =  [\text{det}(s)]^\frac{1}{m}$, for $s$ the image of $\underline{s}\in\underline{A}$ under the left regular representation map of $\underline{A}$.  It is easy to see that this induces a normalized uncurling metric for $\bar{A}$ given by $\bar{L}_{\shortparallel} = \frac{\|\mathbf{1}_{\bar{A}}\|^2}{m} \Omega^T \underline{L}_{\shortparallel} \Omega$, since (\ref{7/7/21.1}) is satisfied because $\Omega$ is an epimorphism (i.e., $\Omega\bar{s} =  \underline{s}$ implies $\Omega(\bar{s})^{-1} =  \underline{s}^{-1}$).  According to (\ref{11/7/20.2}), $\bar{L}_{\shortparallel}$ is associated with the special unital norm on $\bar{A}$ given by $\ell_{\rm sp}(\bar{s}) = [\text{det}(\underline{\mathcal{I}}\Omega\bar{s})]^{\frac{1}{m}}$, where $\underline{\mathcal{I}}$ is the isomorphism from $\underline{A}$ to its left regular representation. 
  \end{proof}
  
   

 The following is an alternative statement of part of the above corollary.   
 
 \begin{corollary}\label{1/1/23.1} An algebra's uncurling space has dimension greater than zero. \end{corollary}
  
Half of the examples of associative algebras $\bar{A}$ collected in Table III of \mbox{Section 5} are such that their left regular representations are not self-adjoint as subalgebras of $M_n(\mathbb{R})$, thereby illustrating Corollary \ref{11/17/22.1}.  Calculations related to such algebras are explicitly highlighted in the examples of Section 4.4 and Section 4.5.3.

\section{More examples, and application to some not associative algebras}
\subsection{Upper triangular Toeplitz matrix algebras}

The uncurling space and unital norm group associated with an upper triangular Toeplitz matrix algebra have some interesting features.
\begin{itemize}\item The uncurling space is given by the left regular representation (the algebra whose elements are lower triangular Toeplitz matrices) multiplied on the left by the exchange matrix. \item The usual norm is sensitive to only a single component of an element, but almost all unital norms are sensitive to all components of an element.\end{itemize}


 \begin{definition}\label{5/4/22.1} $\mathbb{T}_n$ is the $n$-dimensional vector space of real upper triangular Toeplitz $(n\times n)$-matrices. For $s\in \mathbb{T}_n$, the notation $s=(x_1,\dots,x_n)$ indicates that $s$ corresponds to the Toeplitz upper triangular $(n\times n)$-matrix whose entries on the successive diagonals, beginning with the main diagonal and proceeding to the right, are given by the successive vector component values $x_1,\dots,x_n$.   
\end{definition}

In the sequel, for $s\in \mathbb{T}_n$, $s$ can refer either to its format as an $n$-vector or its format as an $(n\times n)$-matrix, depending on the context.


The following result is known, but we include the proof for completeness, and to introduce some notation.  



 \begin{proposition}\label{5/10/22.1} Under the usual matrix product, $\mathbb{T}_n$ is a real $n$-dimensional unital associative algebra.\end{proposition} 
 
 \begin{proof}
Let ${\bf z}$ be the $(n \times n)$-matrix with entries having value $1$ on its superdiagonal  (the diagonal immediately to the right of the main diagonal) and value $0$ otherwise.  It is easily shown that for $j<n$ its $j$-th power, ${\bf z}^j$, is the $(n\times n)$-matrix with entries $1$ on the $j$-th diagonal to the right of the the main diagonal, with all other entries being $0$.   Thus, we can express the matrix format of $s\in\mathbb{T}_n$ as a polynomial in ${\bf z}$, such that
\begin{equation}\label{5/11/22.5} s = x_1 {\bf 1} + x_2 {\bf z} + x_3 {\bf z}^2 + \cdots + x_n {\bf z}^{n-1}.\end{equation}  It is similarly easy to show that ${\bf z}^k$ is the matrix of all zero entries if $k\ge n$.  It follows that for $s,t\in \mathbb{T}_n$, we have $st\in \mathbb{T}_n$.   Finally, for $x_1\ne 0$, we have \begin{equation}\label{5/11/22.6} s^{-1} =  (x_1 {\bf 1} + x_2 {\bf z} + x_3 {\bf z}^2 + \cdots + x_n {\bf z}^{n-1})^{-1}.\end{equation}  The matrix adding to $x_1{\bf 1}$ inside the parentheses on the right-hand-side of (\ref{5/11/22.6}) has only zero as an eigenvalue, so the full expression on the right-hand-side of (\ref{5/11/22.6}) has a power series expansion in ${\bf z}$.  Since ${\bf z}^k$ is the zero matrix for $k\ge n$, the power series expansion is  a $(n-1)$-degree polynomial in ${\bf z}$, which indicates that $s^{-1}\in \mathbb{T}_n$.
 \end{proof}
 
We will show that upper triangular Toeplitz matrix algebras have unital norms sensitive to all element components, providing a great contrast with the usual norm for these algebras, which is sensitive to only one element component.  One can easily generate other algebras with unital norms ``similarly dominating" the usual norm, as seen later.

To approach the problem of deriving the uncurling space for $\mathbb{T}_n$, the following will be important.

\begin{lemma} \label{5/7/22.15} For a unit $s\in \mathbb{T}_n$,  $n>1$, let $\left( s^{-1} \right)_p$ be the $p$-th component of $s^{-1}$ in the latter's $n$-vector format.  Given $k\le n$, then for all $j<k$ we have \begin{equation}\label{5/6/22.10} \frac{\partial \left( s^{-1} \right)_k}{\partial x_j} = \frac{\partial \left( s^{-1} \right)_{k-j+1}}{\partial x_1}.\end{equation}
\end{lemma}

\begin{proof} The proof is by induction. Equation (\ref{5/6/22.10}) is a tautology for the $k=2$ case.  We thus assume the inductive hypothesis that (\ref{5/6/22.10}) is true  for all cases up through $k-1$, and then establish its truth for the $k$-th case.

\enlargethispage{3cm}
We must first express $s^{-1}$ in terms of the components of  $s=(x_1,\dots,x_n)\in \mathbb{T}_n$ with $x_1\ne 0$.   From the proof of Proposition \ref{5/10/22.1},  $s^{-1}$ is an upper triangular Toeplitz matrix.  But as an $n$-vector, it is easily seen that we can write it as  \begin{equation}\label{5/9/22.1} s^{-1} = (\frac{1}{x_1}, \frac{y_1}{x_1},\dots, \frac{y_{n-1}}{x_1}),\end{equation} where the $y_i$ are to be determined. 
Based on the notation introduced in the proof of Proposition \ref{5/10/22.1}, we can use (\ref{5/11/22.5}) and (\ref{5/11/22.6}) to write 
\begin{align}\nonumber s &= x_1 \left({\bf 1} + \frac{x_2}{x_1} {\bf z} + \frac{x_3}{x_1} {\bf z}^2 + \cdots + \frac{x_n}{x_1} {\bf z}^{n-1}\right)\\ s^{-1} &= x_1^{-1} \left({\bf 1} + y_1 {\bf z} + y_2 {\bf z}^2 + \cdots + y_{n-1} {\bf z}^{n-1}\right).  \label{5/18/22.10} \end{align}  Substituting the above two equations into $s^{-1} s = {\bf 1}$ leads to following set of relations, which allow the $y_i$ to be computed recursively:
\begin{align} -y_1 &= \frac{x_2}{x_1} \nonumber \\ -y_2 &= \frac{x_2}{x_1}y_1 + \frac{x_3}{x_1} \nonumber \\ &\;\vdots \nonumber \\ -y_{k-1} &= \frac{x_2}{x_1} y_{k-2} + \frac{x_3}{x_1} y_{k-3} + \cdots + \frac{x_k}{x_1} \label{5/7/22.1} \\ &\;\vdots \nonumber \end{align} 
 
 \noindent For $k>1$, (\ref{5/18/22.10}) implies  $\left( s^{-1} \right)_k = \frac{y_{k-1}}{x_1}$.  Applying this to (\ref{5/7/22.1}) with $k>1$, we find
\begin{eqnarray} \nonumber \left( s^{-1} \right)_k &=& -\frac{x_2}{x_1}\left( s^{-1} \right)_{k-1} - \frac{x_3}{x_1}\left( s^{-1} \right)_{k-2} -\dots  - \frac{x_k}{x_1}\left( s^{-1} \right)_1 \\ &=& -\frac{1}{x_1}\sum_{m=1}^{k-1} x_{m+1}\left( s^{-1} \right)_{k-m}. \label{5/7/22.2}\end{eqnarray}
Differentiating with respect to $x_j$, we find \begin{equation}  \frac{\partial \left( s^{-1} \right)_k}{\partial x_j} = -\frac{1}{x_1} \sum_{m=1}^{k-j} x_{m+1} \frac{\partial \left( s^{-1} \right)_{k-m}}{\partial x_j} -  \frac{1}{x_1}\left( s^{-1} \right)_{k-j+1} \label{5/7/22.4}, \end{equation} where the upper limit on the sum on the right-hand-side is $(k-j)$ rather than $(k-1)$ because the partial derivatives in the sum are zero for $m> k-j$ (i.e., $\left( s^{-1} \right)_p$ is not a function of $x_j$ for $p<j$).  

Next, substituting $(k-j+1)$ for $k$ in (\ref{5/7/22.2}) and then differentiating with respect to $x_1$, we have   \begin{eqnarray} \nonumber
 \frac{\partial \left( s^{-1} \right)_{k-j+1}}{\partial x_1} &=& \frac{1}{x_1^2} \sum_{m=1}^{k-j} x_{m+1} \left( s^{-1} \right)_{k-j-m+1}    - \frac{1}{x_1} \sum_{m=1}^{k-j} x_{m+1} \frac{\partial \left( s^{-1} \right)_{k-j-m+1}}{\partial x_1}  \\ &=& - \frac{1}{x_1}\left( s^{-1} \right)_{k-j+1}  -  \frac{1}{x_1} \sum_{m=1}^{k-j} x_{m+1} \frac{\partial \left( s^{-1} \right)_{k-j-m+1}}{\partial x_1},\label{5/7/22.3}\end{eqnarray} where the equality of the first terms on the right-hand-sides of the two above equalities comes from (\ref{5/7/22.2}) itself (with $(k-j+1)$ substituting for $k$). 
 Now, according to the induction hypothesis, $\frac{\partial \left( s^{-1} \right)_{k-m}}{\partial x_j} = \frac{\partial \left( s^{-1} \right)_{k-j-m+1}}{\partial x_1}$.  Inserting this into (\ref{5/7/22.3}), we see that the left-hand-sides of (\ref{5/7/22.4}) and (\ref{5/7/22.3}) are equal, as was required to be shown.
 \end{proof} 

We recall that $\slashed{\bf 1}_n$ is the exchange matrix.

\begin{theorem} \label{5/4/22.2} For any $t\in \mathbb{T}_n$ let $\hat{t}$ be its $(n\times n)$-matrix format.    Then $\hat{t} \slashed{\bf 1}_n$ is an uncurling metric.  Conversely, each member of the uncurling space associated with $\mathbb{T}_n$ can be expressed as $\hat{t}\slashed{\bf 1}_n$ for some $t\in \mathbb{T}_n$.   \end{theorem}

\enlargethispage{1cm}
\begin{proof} The following notation is more convenient for the moment.  Given $(\gamma_n,\dots,\gamma_1)$ $\in \mathbb{T}_n$, let $G_n$ be its upper triangular Toeplitz $(n\times n)$-matrix format (which happens to be the transpose of its left regular representation), and consider the resulting Hankel matrix, 
\begin{equation}\label{5/7/22.20} H_n \equiv G_n \slashed{\bf 1}_n = \begin{bmatrix} \gamma_1 & \gamma_2 & \gamma_3 & \dots & \gamma_n \\ \gamma_2 & \gamma_3 & \dots & \gamma_n & 0 \\ \gamma_3 & \dots & \iddots & \iddots & 0  \\ \vdots & \iddots & \iddots & \iddots & \vdots \\ \gamma_n & 0 & \dots & \ldots & 0 \end{bmatrix}.\end{equation}
For $\left( s^{-1} \right)_p$ as defined in the statement of Lemma \ref{5/7/22.15}, we then have
\begin{multline}  H_n s^{-1}   =  \Big( \gamma_1\left( s^{-1} \right)_1  + \dots + \gamma_n \left( s^{-1} \right)_n, \\ \gamma_2\left( s^{-1} \right)_1  + \dots + \gamma_n \left( s^{-1} \right)_{n-1}, 
  \dots , 
 \gamma_n \left( s^{-1} \right)_1 \Big). \label{5/7/22.40}
\end{multline}
 We will first demonstrate that $H_n$ is an uncurling metric.  According to (\ref{4/7/23.1}) as referred to by Definition \ref{9/20/23.1}, we are required to demonstrate 
\begin{equation}\label{5/7/22.30} d\left([H_n s^{-1}]\cdot {\sf d}s\right) = 0.\end{equation}  

To verify (\ref{5/7/22.30}), we apply an induction argument.  The $n=2$ case is easily demonstrated to satisfy the above equation (e.g., see Section 4.4).  We assume the induction hypothesis that (\ref{5/7/22.30}) holds for all cases less than $n$, and now need to show that it holds for the $n$-th case.  
We begin by showing that the partial derivative of the first component of $H_n s^{-1}$ with respect to any $x_j$, with $j>1$, minus the partial derivative of the $j$-th component of $H_n s^{-1}$ with respect to $x_1$, is zero. We will then show that the analogous difference of partial derivatives of {\it any} two different components of $H_ns^{-1}$ is also zero.

So, from (\ref{5/7/22.40}), 
\begin{equation}\label{5/18/22.5} \frac{\partial \left(H_n s^{-1}\right)_j}{\partial x_1} = 
\sum_{i=j}^n \gamma_i\frac{\partial \left(s^{-1}\right)_{i-j+1}}{\partial x_1}. 
\end{equation}
But, \begin{equation}\label{5/18/22.6}  \frac{\partial \left(H_n s^{-1}\right)_1}{\partial x_j} = \sum_{i=1}^n \gamma_i\frac{\partial \left(s^{-1}\right)_i}{\partial x_j}
= \sum_{i=j}^n \gamma_i\frac{\partial \left(s^{-1}\right)_i}{\partial x_j} = \sum_{i=j}^n \gamma_i\frac{\partial \left(s^{-1}\right)_{i-j+1}}{\partial x_1},
\end{equation} where the second equality follows (\ref{5/9/22.1}) and (\ref{5/7/22.1}), which indicate that $\frac{\partial \left(s^{-1}\right)_i}{\partial x_j} = 0$ for $i<j$, and the third equality follows from Lemma \ref{5/7/22.15}.
Equations (\ref{5/18/22.5}) and (\ref{5/18/22.6})  demonstrate that the partial derivative of the first component of $H_n s^{-1}$ with respect to any $x_j$, with $j>1$, minus the partial derivative of the $j$-th component of $H_n s^{-1}$ with respect to $x_1$, is equal to zero.

Now we will consider the vanishing of the partial derivative of the $k$-th component of $H_n s^{-1}$ with respect to $x_j$ minus the partial derivative of the $j$-th component of $H_n s^{-1}$ with respect to $x_k$, for all possible $k>j>1$ (i.e, all choices relevant to (\ref{5/7/22.30}) other than considered in the prior paragraph).  In fact, based on the telescoping format of (\ref{5/7/22.40}), we already have the result that this difference vanishes, because it merely reduces to an instance of the theorem for the case of a matrix of the form $H_ks^{-1}$ with $k<n$ - and this case is assumed to hold according to the induction hypothesis.  So, (\ref{5/7/22.30}) is established.

Furthermore, a Hankel matrix of the form (\ref{5/7/22.20}) is the only possibility for an uncurling metric $L$ because any other choice resulting in the coefficients of the terms in any two components of $Ls^{-1}$ cannot lead to a null value for the difference of the two requisite partial derivatives (as is necessary for $d\left([L s^{-1}]\cdot {\sf d}s\right) = 0$) due to linear independence of the individual terms summing to the components. That is, replacing  $H_n$ with matrix \enlargethispage{1cm}$L=(a_{p,q})_{p,q=1}^n$ in (\ref{5/18/22.5}) 
and (\ref{5/18/22.6}), and invoking Lemma \ref{5/7/22.15}, we obtain
\begin{eqnarray}\nonumber \frac{\partial \left(L s^{-1}\right)_j}{\partial x_1}  -  \frac{\partial \left(L s^{-1}\right)_1}{\partial x_j} &=&  \sum_{i=1}^n a_{j,i}\frac{\partial \left(s^{-1}\right)_i}{\partial x_1} - \sum_{i=j}^n a_{1,i}\frac{\partial \left(s^{-1}\right)_i}{\partial x_j} \\  &=& \sum_{i=1}^n a_{j,i}\frac{\partial \left(s^{-1}\right)_i}{\partial x_1} - \sum_{i=j}^n a_{1,i}\frac{\partial \left(s^{-1}\right)_{i-j+1}}{\partial x_1} \nonumber \\
&=& \sum_{i=1}^n a_{j,i}\frac{\partial \left(s^{-1}\right)_i}{\partial x_1} - \sum_{i=1}^{n-j+1} a_{1,i+j-1}\frac{\partial \left(s^{-1}\right)_i}{\partial x_1}.
\label{5/9/22.10}
\end{eqnarray}
This indicates that we must have $a_{j,i}=0$ for $i>n-j+1$, and otherwise have $a_{j,i} = a_{1,i+j-1}$.  Since this must hold for any $j>1$, $L$ must be a Hankel matrix of the form (\ref{5/7/22.20}).  Thus, the $H_n$ are the only possibilities for uncurling metrics. 

Since $t=(\gamma_n,\dots,\gamma_1)$ is an arbitrary member of $\mathbb{T}_n$, the theorem follows.
\end{proof}

\begin{corollary}\label{7/1/22.1} For any $t\in \mathbb{T}_n$, denote the value of the $(1,n)$ entry of its matrix format $\hat{t}$ as $\hat{t}_{1,n}$.  Then $\frac{1}{\hat{t}_{1,n}}\hat{t}\slashed{\bf 1}_n$ is a normalized uncurling metric if $\hat{t}_{1,n}\ne 0$, and all normalized uncurling metrics can be written in this form for some $t\in \mathbb{T}_n$. \end{corollary}
\begin{proof} Theorem \ref{5/4/22.2} has shown that all uncurling metrics of $\mathbb{T}_n$ have the form (\ref{5/7/22.20}).
We now consider constraint (\ref{7/7/21.1}) of Definition \ref{9/20/23.1}.  Let $H_n^{(1)}$ be matrix $H_n$ except with the  $(1,1)$-entry having the value $1$. 
Since $\|{\bf 1}\|^2 =1$ for this algebra, we must show that \[ s^T H_n^{(1)} s^{-1} = s\cdot [H_n^{(1)}s^{-1}] = 1.\] From (\ref{5/9/22.1}) and (\ref{5/7/22.20}), we have
\begin{eqnarray}\nonumber s^T H_n s^{-1} &=& (x_1,\dots,x_n)\cdot \left[ H_n \left(\frac{1}{x_1}, \frac{y_1}{x_1},\dots,\frac{y_{n-1}}{x_1}\right)^T\right] \\ 
 &=&  \gamma_1 + \gamma_2y_1 + \gamma_3y_2 \hspace{3mm} + \hspace{3mm}\gamma_4y_3 \hspace{3mm} + \dots   \hspace{6mm}+\gamma_ny_{n-1} \nonumber \\
&+&  \gamma_2 \frac{x_2}{x_1} + \gamma_3\frac{x_2}{x_1}y_1 + \hspace{2mm} \gamma_4\frac{x_2}{x_1}y_2 + \dots +  \gamma_n\frac{x_2}{x_1}y_{n-2} \nonumber \\ &+&\gamma_3 \frac{x_3}{x_1}  + \hspace{3mm}\gamma_4\frac{x_3}{x_1}y_1 +  \dots \hspace{4mm} +  \gamma_n\frac{x_3}{x_1}y_{n-3} \nonumber \\ 
&+&  \hspace{5mm} \vdots \hspace{2.5cm} \iddots \bigskip\bigskip \nonumber \\ &+& \gamma_{n-1} \frac{x_2}{x_1} \hspace{3mm} + \gamma_n \frac{x_2}{x_1}y_1 \nonumber \\ &+& \hspace{2mm} \gamma_n\frac{x_n}{x_1}. \label{6/23/22.1}
 \end{eqnarray}
 For any $i>1$, if we collect all the terms on the right-hand-side above that each have $\gamma_i$ as a coefficient, and then apply (\ref{5/7/22.1}), it is quickly seen that the sum of all terms having $\gamma_i$ as a coefficient vanishes.  It follows that
 $s^T H_n s^{-1} =\gamma_1$.   Thus, the requisite matrix satisfying constraint (\ref{7/7/21.1}) of Definition \ref{9/20/23.1} will be the right-hand-side of (\ref{5/7/22.20}) with $\gamma_1=1$.  Setting $\frac{1}{\hat{t}_{1,n}}t = (\gamma_n,\dots,\gamma_2, 1)$, the corollary follows.
\end{proof}

\begin{corollary}\label{5/17/22.1} For $s= (x_1,\dots,x_n)\in \mathbb{T}_n$ with $x_1\ne 0$, the usual norm is $x_1^n$.  The special unital norm family of $\mathbb{T}_n$ is the $(n-1)$-parameter family,
\begin{equation} \label{5/5/22.6} \ell_{\gamma_2,\dots,\gamma_n}(x_1,\dots,x_n) = x_1 e^{P_{\gamma_2,\dots,\gamma_n}\left(\frac{x_2}{x_1},\dots,\frac{x_n}{x_1}\right)},\end{equation} where $P_{\gamma_2,\dots,\gamma_n}(w_1,\dots,w_{n-1})$ is a real $(n-1)$-degree multivariate polynomial with coefficients dependent on the $(n-1)$ real parameters $\gamma_2,\dots,\gamma_n$.    \end{corollary}

\begin{proof} The usual norm of $s=(x_1,\dots,x_n)\in \mathbb{T}_n$ is immediately seen to be $x_1^n$.

With respect to the notation $G_n,H_n$ introduced at the beginning of the proof of Theorem \ref{5/4/22.2}, each $(\gamma_n,\dots,\gamma_2,1)\in\mathbb{T}_n$ is associated with a Toeplitz matrix $G_n^{(1)}$ having its $(1,n)$-entry equal to $1$.  Corollary \ref{7/1/22.1} indicates that $H^{(1)}_n\equiv G_n^{(1)}\slashed{\bf 1}_n$ is a normalized uncurling metric. From (\ref{11/7/20.2}),  the associated unital norm is $e^{\left[\int_{\bf 1}^s (H_n^{(1)} v^{-1})\cdot dv\right]}$.  Elementary integration and a simple induction argument quickly results in (\ref{5/5/22.6}) and the $(n-1)$-degree of the multivariate polynomial in question.
\end{proof}
 
 \subsection{Extension to certain not associative algebras}
The two-dimensional real unital algebras (there are three, up to an isomorphism) have generalizations as the Cayley-Dixon sequence algebras ($\mathbb{R}$ and $\mathbb{C}$ being its lowest dimensional members),  the Spin Factor Jordan algebras (generalizing $\slashed{\mathbb{C}}$), and real upper triangular matrix algebras with at least some redundant diagonal components (generalizing the lowest dimensional case, $\mathbb{D}$).  Their unital norm groups will be the subject of this subsection. 

Spin Factor Jordan Algebras and Cayley-Dixon Algebras after the quaternions are power-associative but not associative.  While the usual norm is thereby inapplicable, there is an ``Algebraic Norm" replacement \cite{jacobson:1963} (based on the minimal polynomial, and thus analogous to the usual norm as given by the constant term of a characteristic polynomial).  For the Spin Factor Jordan Algebras and Cayley-Dixon Algebras, the Algebraic Norms are essentially the  Minkowski  norm and Euclidean norm, respectively.  Thus, useful norms are not restricted to associative (or alternative) algebras, even though it turns out that the multiplicative property (i.e., $N(s_1s_2)=N(s_1)N(s_2)$) must then be sacrificed  \cite{schafer:1963}. 
So, despite a strong bias against a norm not satisfying the multiplicative property, these examples indicate that ``algebraic" norms that are not multiplicative can nevertheless be very significant \cite{mccrimmon:2004}. 

Of course, apart from the omnipresent unital norm group member comparable to the usual norm as guaranteed by Corollary \ref{11/17/22.1} for real finite-dimensional unital associative algebras (as well as all unital norms resulting from two-dimensional real unital algebras), unital norms are not typically multiplicative.  The multiplicative property is instead replaced by (\ref{5/23/23.3}).  The latter equation also replaces the multiplicative property for the Algebraic Norm in the cases of the not associative Spin Factor Jordan Algebras and Cayley-Dixon Algebras. 
Like the Algebraic Norm, the unital norm/uncurling metric constructions remain valid for these latter two classes of algebras.  The way we deal with these particular not associative algebras is as follows.

\subsubsection{Unital norms for $^*$-algebras whose Hermitian elements are real} 

Consider an algebra associated with an involutive antiautomorphism ``
$^*$ " (known as a $^*$-algebra, star-algebra, or involutive algebra) such that the Hermitian elements are members of $\mathbb{R}$ (thus, $s+s^*$ and $s^*s$ are both real; also, $s^*s = ss^*$ because $(s+s^*)s = s(s+s^*)$ since $(s+s^*)$ is real). 
Then \[ X^2 - (s+s^*)X + s^*s, \] is the minimal polynomial of any not real element $s$, and annihilates all elements, real and not real.  Consistent with comments in the second paragraph of Section 4.2, we define its constant term to be the ``Algebraic Norm" of an element $s$, by analogy with the characteristic polynomial defined for an associative algebra, which supplies the norm from the constant term of that element-associated polynomial annihilating all elements.


On these $^*$-algebras, the multiplicative identity $\mathbf{1}$ is $1$ (where the latter is thought of as the algebra element that is the multiplicative identity on the subalgebra $\mathbb{R}$ rather than a field element).  Accordingly, for these algebras, we define the expression $\|{\bf 1}\|^2$ (e.g., that appears in Definition \ref{9/20/23.1}) to be $1$.  We also have,

\begin{lemma}\label{9/14/22.2} The inverse of a unit $s$ in a finite-dimensional $^*$-algebra whose Hermitian elements are real is unique, and given by $s^{-1}=\frac{s^*}{ss^*}$.\end{lemma}

\begin{proof} 
We can write $s=\frac{s+s^*}{2} +\frac{s-s^*}{2}$, where the first  term on the right-hand-side is real and the second term on the right-hand-side is antisymmetric with respect to application of the $^*$-operation.  We define $\sigma$ to be the first term on the right-hand-side, and $\bar{s}$ to be the second term on the right-hand-side.  Suppose $x = \chi + \bar{x}$ is an inverse of $s$, where $\chi$ is real, and $\bar{x}$ is antisymmetric. Then \begin{equation}\label{9/14/22.1} sx = \sigma\chi + \sigma\bar{x}+\chi\bar{s} + \bar{s}\bar{x} = 1,\end{equation} and $xs = \sigma\chi + \sigma\bar{x}+\chi\bar{s} + \bar{x}\bar{s} = 1$.  Thus, $\bar{s}\bar{x} = \bar{x}\bar{s}$.  Since $\bar{s}^* = (\frac{s-s^*}{2})^* = -(\frac{s-s^*}{2}) = -\bar{s}$, we have $(\bar{s}\bar{x})^* = \bar{x}^*\bar{s}^* = \bar{x}\bar{s} = \bar{s}\bar{x}$.  Hence, $\bar{s}\bar{x}$ is real.  It follows from (\ref{9/14/22.1}) that $\sigma\chi+\bar{s}\bar{x} =1$ and $\sigma\bar{x}+\chi\bar{s} = 0$.  The latter has has a unique solution, yielding $\chi+\bar{x} = \frac{\sigma-\bar{s}}{\sigma^2-\bar{s}\bar{s}} = \frac{s^*}{ss^*}$, where the second equality follows since the right-hand-side is obviously an inverse ($1=s\frac{s^*}{ss^*} = \frac{s^*}{ss^*}s$ because $s^*s=ss^*$ is real and thus self-adjoint), and the inverse has just now been seen to be unique.
\end{proof}

Extending Definition \ref{9/20/23.1} to include this class of $^*$-algebras, we can now provide a special unital norm  by identifying a linear transformation $L$ satisfying (\ref{4/7/23.1}), (\ref{7/7/21.1}) and then evaluating (\ref{11/7/20.2}).  
 
 \begin{proposition}\label{1/11/22.1} Given a $^*$-algebra with vector space of elements $\mathbb{R}^n$ and whose Hermitian elements are members of $\mathbb{R}$, the Algebraic Norm on the $^*$-algebra in a sufficiently small neighborhood of ${\bf 1}$ is the square of a special unital norm. \end{proposition}

\begin{proof} By definition, the Algebraic Norm is \begin{equation}\label{1/15/22.1} N(s)= ss^* = N(s^*)\in\mathbb{R}.\end{equation} 
 We then have $N(\alpha s) = \alpha^2 N(s)$ for $\alpha\in \mathbb{R}$, i.e., degree-2 homogeneity of $N(s)$.  
For $N(s)\ne 0$, Lemma \ref{9/14/22.2} implies 
\begin{equation}\label{1/12/22.1} s^{-1} = \frac{s^*}{N(s)}.\end{equation}
Applying $N(\cdot)$ to both sides of  (\ref{1/12/22.1}), the degree-2 homogeneity and (\ref{1/15/22.1}) imply $N(s)N(s^{-1})=1$, so that \begin{equation}\label{1/15/22.2} \sqrt{N(s)}\sqrt{N(s^{-1})}=1. \end{equation}

We are required to show that $\sqrt{N(s)}$ is a special unital norm.  Thus, we must first verify that $\sqrt{N(s)}$ satisfies the requirements of the extended  version of Definition \ref{9/20/23.1} as indicated prior to the statement of this proposition. 

To begin with, (\ref{1/15/22.2}) already satisfies  (\ref{5/23/23.3}).
 Equation (\ref{1/15/22.1}) indicates that $N(s)$ is a quadratic form, so there is an orthonormal basis with respect to which the members of the algebra can be written as $s= (s_1,\dots, s_i,\dots)$ with $N(s) = \sum_i \delta_i s_i^2$ and $\delta_i = -1$, $0$, or $1$, since $^*$ is an involution.  Let $K_1 \equiv {\rm diag}\{\delta_i\}$.   Then, \begin{equation}\label{2/6/22.1} N(s) = s\cdot(K_1 s).\end{equation}
If $s^{-1}$ exists (i.e., $N(s)\ne 0$), we then have \begin{equation}\label{2/6/22.2} \nabla\left[\sqrt{N(s)}\right] = \frac{K_1s}{\sqrt{N(s)}}.\end{equation}  Define $K_2$ to be the algebra's involutive antiautomorphism operation $^*$.  It follows that $K_1K_2K_2[\frac{s}{\sqrt{N(s)}}] = \nabla\left[\sqrt{N(s)}\right]$. Defining $L\equiv K_1K_2$, dividing both sides of the equation in the last sentence by $\sqrt{N(s)}$, and using (\ref{1/15/22.1}) and (\ref{1/12/22.1}), we then have \begin{equation}\label{1/12/22.2} Ls^{-1}  = \frac{\nabla\left[\sqrt{N(s)}\right]}{\sqrt{N(s)}} .\end{equation} This equation  is valid for $N(s)>0$, and the latter pertains in a sufficiently small neighborhood of ${\bf 1}$.  Thus,  (\ref{11/7/20.1}) is satisfied.  From (\ref{2/6/22.1}), (\ref{2/6/22.2}), and (\ref{1/12/22.2}), we have \[ s\cdot [Ls^{-1}] = s\cdot K_1\left[\frac{s}{N(s)}\right] = 1,\] which is (\ref{7/7/21.1}).  
 
 To establish that $\sqrt{N(s)}$ is a unital norm, we  now only need to verify that (\ref{4/7/23.1}) is satisfied, i.e., $d\left([Ls^{-1}]\cdot {\sf d}s\right)=0$.  According to (\ref{2/6/22.2}), (\ref{1/12/22.2}), and the definition  of $K_1$, we have $d\left([Ls^{-1}]\cdot {\sf d}s\right) = d\left(\sum_i \frac{\delta_is_i ds_i}{N(s)}\right)$.  Performing the requisite differentiation to obtain the component of the exterior derivative associated with a choice of $i$ and $j$, we find \[ \frac{\partial [Ls^{-1}]_j}{\partial s_i} - \frac{\partial [Ls^{-1}]_i}{\partial s_j} = -2\frac{(\delta_j s_j)(\delta_i s_i)}{N^2(s)} +
2\frac{(\delta_i s_i)(\delta_j s_j)}{N^2(s)} = 0,  \] where $\frac{\partial [Ls^{-1}]_j}{\partial s_i}$ is the partial derivative of the $j$-th component of $[Ls^{-1}]$ with respect to the $i$-th component of $s$, etc. Thus,  $d\left([Ls^{-1}]\cdot {\sf d}s\right) = 0$.

Consequently, $\sqrt{N(s)}$ satisfies the special unital norm requirements of the extended version of Definition \ref{9/20/23.1}.  
\end{proof}


\subsubsection{The Cayley-Dixon sequence}  From the proof of Proposition \ref{1/11/22.1}, it is easily shown that a normalized uncurling metric associated with an algebra of the Cayley-Dixon sequence (all of whose members are $^*$-algebras whose Hermitian elements are real) is given by ${\rm diag}\{1,-1,-1,\dots,-1\}$, with the associated special unital norm seen to be the Euclidean norm. Although it is easily shown from results in Section 2.3 that $\mathbb{C}$ is associated with a 1-parameter family of normalized uncurling metrics, a simple argument (using Lemma \ref{9/14/22.2}) indicates that all other algebras of the sequence have only the above normalized uncurling metric, i.e., the special unital norm family has only one member.  
 
\subsubsection{Spin Factor Jordan Algebras}  Consider the bilinear product ``$\,\bullet\,$" such that for
$\alpha,\beta\in \mathbb{R},\,\, {\sf a},{\sf b}\in \mathbb{R}^m$,  \begin{equation}  (\alpha+ {\sf a})\bullet (\beta + {\sf b}) \equiv
    [\alpha \beta + {\sf a} \cdot {\sf b}] + [ \beta {\sf a} + \alpha {\sf b}] \in \mathbb{R}\oplus \mathbb{R}^m. \label{12/24/18.1}\end{equation}  This makes the vector space $\mathbb{R}\oplus \mathbb{R}^m$ a unital algebra with the multiplicative identity ${\bf 1}$ as simply $1$ (where in this instance $1$ is understood to be an element of the algebra rather than a member of the field).  For $m\ge 2$, this is known as a Spin Factor Jordan Algebra \cite{mccrimmon:2004}.

 There are two involutive antiautomorphisms on each of these commutative algebras, one of which is the identity and the other is the conjugate operation, $(\alpha+ {\sf a})^* = (\alpha- {\sf a})$. The latter is the only one of the two involutive antiautomorphisms  for which the Hermitian elements are real.  Applying Proposition \ref{1/11/22.1}, it is easily shown that the Algebraic Norm is the special unital norm given by the Minkowski norm.  Similar to the case of the Cayley-Dixon algebras, it is easy to show that this is the only special unital norm.

 
 \subsection{The algebra with vector space of elements $\mathbb{R}^n$ and component-wise multiplication} A higher dimensional analogue of $\mathbb{R}\oplus \mathbb{R}$, taken as the algebraic direct sum of $n$ copies of $\mathbb{R}$, is associated with an $(n-1)$-parameter family of  special unital norms.

 This algebra is isomorphic to the subalgebra of $ M_n(\mathbb{R})$  whose members are diagonal matrices. In this example, ${\bf 1} = (1,\dots,1)$ and $\|{\bf 1}\|^2 = n$.   It is easily appreciated that (\ref{4/7/23.1}) forwards the uncurling space ${\rm diag} \{ \sigma_1, \sigma_2,\dots,\sigma_n\}$, so the uncurling space is the same as the members of the left regular representation of the algebra.  This means that for $s=(x_1,\dots,x_n)$, we must have $L[(x_1,x_2,\dots,x_n)^{-1}] = (\frac{\sigma_1}{x_1}, \frac{\sigma_2}{x_2},\dots,\frac{\sigma_n}{x_n})$.     Evaluation of (\ref{11/7/20.2}) indicates that the members of the unital norm group are given by  \begin{equation}\label{1/15/22.5} \ell_{\sigma_1,\dots,\sigma_n}(x_1,\dots,x_n) = (x_1^{\sigma_1} x_2^{\sigma_2} \dots x_n^{\sigma_n})^\frac{1}{n}. \end{equation} From (\ref{7/7/21.1}), the special unital norms are given by (\ref{1/15/22.5}) under the constraint $\sum_1^n \sigma_i = n$. Thus, for this algebra we have a $(n-1)$-dimensional special unital norm family.   
 
 The geometric mean is the special unital norm given by (\ref{1/15/22.5}) with $\sigma_1=\sigma_2=\dots =\sigma_n = 1$.  
   
 \subsection{Upper triangular matrix algebras with independent entries on the main diagonal} Consider the real three-dimensional algebra with elements $(x,y,z)$ which multiply according to the matrices 
$\left[
\begin{array}{ccc}
 x & z   \\
 0 & y    
\end{array}
\right]$.  Or in other words, the algebra is $\left\{ \left[
\begin{array}{ccc}
 x & z   \\
 0 & y    
\end{array}
\right],x,y,z\in\mathbb{R}\right\}$.  We have $(x,y,z)^{-1} = \left( \frac{1}{x}, \frac{1}{y}, -\frac{z}{xy} \right)$, and ${\bf 1}= (1,1,0)$, and $\|{\bf 1}\|^2 = 2$.  Starting with a generic symmetric $L = \left[
\begin{array}{ccc}
 \alpha & \beta & \gamma   \\
 \beta & \delta & \epsilon \\ \gamma & \epsilon & \xi    
\end{array}
\right]$, 
it is required that the curl of $L(x,y,z)^{-1}$, which is \[\left( \frac{\epsilon}{xy} - \frac{\epsilon}{y^2} + \frac{\xi z}{xy^2},\,\,\, \frac{\gamma}{x^2} - \frac{\gamma}{xy} - \frac{\xi z}{x^2y}, \,\,\, \frac{\beta}{y^2} - \frac{\beta}{x^2} - \frac{\gamma z}{xy^2} + \frac{\epsilon z}{x^2y} \right),\] be such that each of its components vanishes in a neighborhood of $(1,1,0)$.  This implies that all entries of the matrix $L$ must be zero except for $\alpha$ and $\delta$.   Thus, $L(x,y,z)^{-1} = (\frac{\alpha}{x}, \frac{\delta}{y},0)$.  From (\ref{7/7/21.1}), we have the requirement that a normalized uncurling metric must satisfy \[ (x,y,z)\cdot [L(x,y,z)^{-1}] = \alpha +\delta =  2.\]  Equivalently, we can alternatively specify $\alpha \equiv 1+\sigma$ and $\delta \equiv 1-\sigma$, for arbitrary $\sigma\in\mathbb{R}$.   We thus have the normalized uncurling metric family composed of all  \begin{equation}\label{11/17/22.2} L = \left[
\begin{array}{ccc}
 1+\sigma & 0 & 0   \\
 0 & 1-\sigma & 0 \\ 0 & 0 & 0    
\end{array}
\right],\end{equation} $\sigma\in\mathbb{R}$.  In this case, each member of the family is singular.  Equation (\ref{11/7/20.2}) then leads to the 1-parameter family of special unital norms, \begin{equation}\label{12/21/21.1} \ell_\sigma(x,y,z) = x^\frac{1+\sigma}{2} y^\frac{1-\sigma}{2} = (xy)^{\frac{1}{2}}\left(\frac{x}{y}\right)^{\frac{\sigma}{2}}.\end{equation} 

The left regular representation of an element $(x,y,z)$ is given by $\left[\begin{array}{ccc}
x & 0 & 0   \\
 0 & y & 0 \\ 0 & z & x    
\end{array} \right]$.  Thus, the $\mathbb{R}^{3\times 3}$ transpose operation maps outside the algebra's representation.  However, there is an epimorphism to the two-dimensional subalgebra defined by elements $(x,y,0)$, and the latter's left regular representation is 
 $\left[\begin{array}{ccc}
x & 0  \\
 0 & y  
\end{array} \right]$.  The usual norm of the subalgebra is thus $xy$, and one of the special unital norms on the subalgebra is $(xy)^{\frac{1}{2}}$, and these differ only by an exponent - consistent with Theorem \ref{5/13/22.1}.   Note that the latter unital norm on the subalgebra provides the values for a special unital norm on the full algebra, corresponding to the choice of normalized uncurling metric in (\ref{11/17/22.2}) given by $\sigma = 0$, i.e., $\ell_0(x,y,z) = (xy)^{\frac{1}{2}}$, consistent with \mbox{Corollary \ref{11/17/22.1}}.  The fact that the transpose operation on $\mathbb{R}^{3\times 3}$ does not map the above left regular representation of $\left\{\left[\begin{array}{ccc}
x & z  \\
 0 & y  
\end{array} \right]:(x,y,z)\in \mathbb{R}^3\right\}$ to itself explains why there is no unital norm that is a power of the usual norm $x^2y$.
  
\subsection{Some higher dimensional algebras with unital norms that are sensitive to more element components than is  usual norm}

The usual norm on $ M_n(\mathbb{R})$ is sensitive to all components of an element of the algebra (because the determinant of an element of $ M_n(\mathbb{R})$  is sensitive to the values of all entries of its matrix).    However, this is not true in general for associative algebras.  For example, for any algebra isomorphic to a triangular matrix algebra, the usual norm is only sensitive to the components on the main diagonal of the triangular matrix representation.  This is also true for the unital norm if all entries on the main diagonal of the triangular matrix representation are allowed to be distinct.  But if  at least two entries on the main diagonal always take equivalent values for the particular algebra,  then almost all of its unital norms family members will be sensitive to one or more components not on the main diagonal, in addition to all of the components on the main diagonal (these are matrix algebras for which an element that is not a diagonal matrix is never diagonalizable).  

It is easily seen from the presentation in Section 2.7 that the dual numbers are a two-dimensional example of this, and it
generalizes to higher dimensions as exemplified by Corollary \ref{5/17/22.1}.  We will now highlight an additional feature using the next three examples.  

\subsubsection{Another three-dimensional algebra $(\mathbb{T}_3)$}Consider the real algebra with elements  $(x,y,z)$ that multiply according to the matrices $ \left[
\begin{array}{ccc}
x & y & z   \\
0 & x & y \\
0 & 0 & x  
\end{array}
\right]$.    The constraints (\ref{7/21/23.1}) and (\ref{7/7/21.1}) lead to the normalized uncurling metrics, $ L = \left[
\begin{array}{ccc}
1 & \beta & \gamma   \\
\beta & \gamma & 0 \\
\gamma & 0 & 0  
\end{array}
\right]$, $\beta,\gamma\in\mathbb{R}$, which is just the $n=3$ case of Corollary \ref{7/1/22.1}.    Evaluating (\ref{11/7/20.2}), we now have a 2-parameter family of special unital norms, \begin{equation}\label{3/31/21.3} \ell_{\beta,\gamma}(x,y,z) = x\, e^{ \beta\frac{y}{x}  +\gamma \frac{z}{x} -\gamma \frac{ 1}{2}\left(\frac{y}{x}\right)^2 }.\end{equation}    The choice $\beta = \gamma = 0$ produces the special unital norm, $\ell_{0,0}(x,y,z)= x$, which is comparable to the usual norm, $x^3$. 
 The other special unital norms are distinguished by their dependence on $y$ and/or $z$, i.e., unlike the usual norm.  The non-units are associated with a different type of unital norm singularity (as $x$ approaches 0) than seen in examples considered so far other than $\mathbb{T}_n$.
 
 \subsubsection{A four-dimensional algebra} Consider the real unital algebra with elements 
$(x,y,z,w)$ which multiply according to the matrices 
$\left[
\begin{array}{ccccc}
x & y & w   \\
0 & x & z \\
0 & 0 & x  
\end{array}
\right]$.
Straightforward application of the constraints (\ref{7/21/23.1}) and (\ref{7/7/21.1}) leads to the normalized uncurling metric family 
$\left[
\begin{array}{cccc}
1 & \beta & \gamma & \delta   \\
\beta & 0 & \frac{1}{2}\delta & 0 \\
\gamma & \frac{1}{2}\delta &  0 & 0   \\
\delta & 0 &  0 & 0
\end{array}
\right]$, and the associated special unital norm family,
\[ \ell_{\beta,\gamma,\delta}(x,y,z,w) = x\, e^{\beta\frac{y}{x} + \gamma\frac{z}{x} +\delta\frac{w}{x}- \frac{\delta}{2}\left(\frac{yz}{x^2}\right)}.\]


\subsubsection{A five-dimensional algebra}Consider the real unital algebra with elements 
$(x,y,z,v,w)$ which multiply according to the matrices
$\left[
\begin{array}{ccccc}
x & z & w   \\
0 & x & v \\
0 & 0 & y  
\end{array}
\right]$.
Straightforward application of the constraints (\ref{7/21/23.1}) and (\ref{7/7/21.1}) leads to the normalized uncurling metric family 
$\left[
\begin{array}{ccccc}
1+\sigma & 0 & \beta & 0  & 0   \\
0 & 1-\sigma & 0 & 0 & 0 \\
\beta & 0 &  0 & 0 & 0  \\
 0 & 0 &  0 & 0 & 0  \\
0 & 0 &  0 & 0 & 0
\end{array}
\right]$, and the associated special unital norm family 
\[ \ell_{\sigma,\beta}(x,y,z,v,w) =  (xy)^{\frac{1}{2}}\left(\frac{x}{y}\right)^{\frac{\sigma}{2}} e^{\frac{\beta}{2} \frac{z}{x}}.\]

In this example, almost all unital norms are influenced the element's $z$-component, while this sensitivity is not present with the usual norm (given by $x^4y$).  However, like the usual norm,  the unital norm remains insensitive to the $w$ and $v$ components.  Once again, there is a new unital norm singularity type as non-units are approached.  

Finally, the left regular representation of this algebra with elements $(x,y,z,v,w)$ is given by $ \left[
\begin{array}{ccccc}
x & 0 & 0 & 0 & 0   \\
0 & y & 0 & 0 & 0 \\
z & 0 &  x & 0 & 0  \\
0 & v &  0 & x & 0  \\
0 & w &  0 & z & x
\end{array}
\right].$ 
The transpose operation on $\mathbb{R}^{5\times 5}$ does not map this representation algebra into itself, due to its effect on the last three columns and last three rows - which explains why there is no unital norm that is a power of the usual norm $x^4y$.  However, there is an epimorphism to the two-dimensional subalgebra with elements $(x,y,0,0,0)$.  The subalgebra's left regular representation is $\left[\begin{array}{ccc}
x & 0  \\
 0 & y  
\end{array} \right]$, which leads to the special unital norm $(xy)^\frac{1}{2}$ according to \mbox{Theorem \ref{5/13/22.1}}.  We can now make analogous observations to those made in Section 4.4, which include consistency with \mbox{Corollary \ref{11/17/22.1}}.
 
\section{Applying the functor relevant to epimorphisms, and a conjecture}
\subsection{Epimorphisms} Table III collects previously presented examples of uncurling spaces and associated special unital norms.  Greek letters represent parameters ranging over $\mathbb{R}$.  The Spin Factor Jordan Algebras and Cayley-Dixon sequence algebras that are not associative are treated via their formulation as \mbox{$^*$-algebras} whose Hermitian elements are real, as in Section 4.2. The uncurling metrics and special unital norms are written with respect to the standard basis.  Each of the fourteen rows of Table III refers to the algebra given as the first entry of the row.  $\slashed{\mathbb{C}}$ is the algebra of split complex numbers, $\mathbb{D}$ is the algebra of the dual numbers, and $\mathbb{H}$ is the quaternion algebra.  $\mathbb{T}_n$ is the $n$-dimensional upper triangular Toeplitz matrix algebra.  The notation $A_q$ refers to the algebra in the $q$-th row.  In the third entry of row 14, the exponent is a multivariate polynomial in the indicated variables, with the subscripted parameters multiplying constituent monomial terms.  $\slashed{\mathbf{1}}_n$ is the $(n\times n)$-exchange matrix.  The expressions $\mathbb{C}_\text{left}$, $\slashed{\mathbb{C}}_\text{left}$, $[\mathbb{R}\oplus\mathbb{R}]_\text{left}$, $\mathbb{D}_\text{left}$, and $[\mathbb{T}_n]_\text{left}$, indicate the left regular representations.  


\bigskip

The examples in Table III illustrate many of the points made earlier in this paper, and in particular the role of the functor of Section 2.9 and its resulting corollaries.    
 
\begin{itemize}\item $\mathbf{u}_{M_n(\mathbb{R})}$ is one-dimensional (Theorem 3.1), from which one concludes that $M_n(\mathbb{R})$ is simple (Corollary \ref{8/17/22.1}) - a demonstration distinct from the usual one.
  \item Up to an isomorphism, there are only three two-dimensional real unital algebras, e.g., $\mathbb{C}$, $\slashed{\mathbb{C}}$, and $\mathbb{D}$.  Algebras $\mathbb{R}\oplus\mathbb{R}$ and $\slashed{\mathbb{C}}$ are isomorphic via the linear transformation $\left[\begin{array}{ccc}
 1 & -1   \\ 
 1 & 1    
\end{array}\right]$, but they are not isometrically isomorphic, since the isomorphism does not have unit determinant.  The respective left regular representations of $\mathbb{C}$, $\slashed{\mathbb{C}}$, and $\mathbb{D}$, are $\left[\begin{array}{ccc}
 x & -y   \\ 
 y & x    
\end{array}\right]$, $\left[\begin{array}{ccc}
 x & y   \\ 
 y & x    
\end{array}\right]$, and $\left[\begin{array}{ccc}
 x & 0   \\ 
 y & x    
\end{array}\right]$.  For each of the algebras $\mathbb{C}$,  $\slashed{\mathbb{C}}$, and $\mathbb{D}$, the set of members of the algebra's uncurling space is the same as the set of members of the algebra's left regular representation algebra except with each of the latter's matrix members multiplied on the left by the exchange matrix $\slashed{\mathbf{1}}_2$ (a similar statement holds for $\mathbb{T}_n$ and $\slashed{\mathbf{1}}_n$).  Also, for algebra $\slashed{\mathbb{C}}$, the set of members of the algebra's uncurling space is the same as the set of members of its left regular representation algebra, which is also true for $\mathbb{R}\oplus\mathbb{R}$. 
 \item While $\slashed{\mathbb{C}}$ and $\mathbb{R}\oplus\mathbb{R}$, are isomorphic, an isomorphism between any two of $\mathbb{C}$, $\slashed{\mathbb{C}}$, and $\mathbb{D}$ is excluded by applying the third and fourth bullet points of \mbox{Corollary \ref{8/18/23.2}}. 
\item  There is no morphism pointing from $\mathbf{u}_{\mathbb{C}}$ to $\mathbf{u}_{\mathbb{R}}$. That is, there is clearly no nonzero real matrix $W$ such that $W^T[\sigma]W \subset \left[\begin{array}{ccc}
 \alpha & \beta   \\ 
 \beta & -\alpha    
\end{array}\right]$, since all the real matrix members of $W^T[\sigma]W$  must have determinant zero - while $\left[\begin{array}{ccc}
 \alpha & \beta   \\ 
 \beta & -\alpha    
\end{array}\right]$ 
 has determinant zero if and only if $\alpha,\beta = 0$ and is thereby the zero matrix. This is distinct from the usual demonstration that $\mathbb{C}$ is simple.  On the other hand, there is clearly a morphism from $\mathbf{u}_{\mathbb{R}\oplus\mathbb{R}}$ to $\mathbf{u}_{\mathbb{R}}$.
 \end{itemize}

 
\hspace{-3.5cm} $ 
\begin{array}{cccccccc}
{\rm \underline{algebra}} & \mbox{\underline{usual norm}} & \mbox{\underline{special unital norm family}} & \mbox{\underline{uncurling space}} &  \bigskip \\ 
M_n(\mathbb{R}) & ({\rm det}(s))^n & ({\rm det}(s))^{\frac{1}{n}} & \alpha \,(n^2\times n^2)\mbox{\rm -``transpose\, matrix"}  \\  \\ 
\mathbb{C} & x^2+y^2 = r^2 & \sqrt{x^2+y^2}\,e^{\beta{\rm arctan}\left(\frac{y}{x}\right)} = re^{\beta\theta} & \left[\begin{array}{ccc}
 \alpha & \beta   \\ 
 \beta & -\alpha    
\end{array}\right]  =  \slashed{\mathbf{1}}_2 \, \mathbb{C}_\text{left} \\ \\ \slashed{\mathbb{C}}
 & x^2-y^2 & \begin{array}{cc}
 \sqrt{x^2-y^2}\, e^{\beta{\rm arctanh}\left(\frac{y}{x}\right)}   \\
 = (x^2-y^2)^{\frac{1}{2}}\left( \frac{x+y}{x-y}\right)^\frac{\beta}{2}    
\end{array} & \left[\begin{array}{ccc}
 \alpha & \beta   \\
 \beta & \alpha    
\end{array}\right]  = \slashed{\mathbb{C}}_\text{left} \\ \\
\mathbb{R}\oplus\mathbb{R} & xy & (xy)^{\frac{1}{2}}\left(\frac{x}{y}\right)^{\frac{\sigma}{2}} & \left[\begin{array}{ccc}
\sigma_1 & 0   \\
 0 & \sigma_2    
\end{array}\right] = [\mathbb{R}\oplus\mathbb{R}]_\text{left}  \\  \\
\vspace{-1mm} \mathbb{D}  & x^2 & x\, e^{\beta\frac{y}{x}} & \left[\begin{array}{ccc}
 \alpha & \beta   \\
 \beta & 0    
\end{array}\right] = \slashed{\mathbf{1}}_2\mathbb{D}_\text{left} \\ \\

\vspace{-1mm} \mathbb{H}  & (x^2+y^2+z^2+w^2)^2 & \sqrt{x^2+y^2+z^2+w^2} &  \,{\rm diag}\{\alpha,-\alpha,-\alpha,-\alpha\}   \\ \\

\mbox{\rm Cayley-Dixon seq., } \mbox{after }\mathbb{H} & (\mbox{\rm Euclidean norm})^2 & {\rm Euclidean\, Norm} &  \,{\rm diag}\{\alpha,-\alpha,-\alpha,\dots,-\alpha\} &  \\ \\
\mbox{Spin Factor Jordan} & (\mbox{\rm Minkowski norm})^2 & \mbox{\rm Minkowski norm} &  \,{\rm diag}\{\alpha,\alpha,\dots,\alpha\} & \\ \\ \bigoplus_{i=1}^n \mathbb{R}
& \prod x_i & \left(\prod x_i^{\sigma_i}\right)^\frac{1}{n}, \sum\sigma_i=n & {\rm diag} \{\sigma_1,\dots,\sigma_n\} & \\ \\
\left\{ \left[\begin{array}{ccc}
 x & z   \\
 0 & y    
\end{array}\right]: (x,y,z)\in \mathbb{R}^3 \right\} & x^2y & (xy)^{\frac{1}{2}}\left(\frac{x}{y}\right)^{\frac{\sigma}{2}} & \left[\begin{array}{ccc}
\sigma_1  & 0 & 0   \\
 0 & \sigma_2 & 0 \\
 0 & 0 & 0    
\end{array}\right] & \\ \\
\left\{  \left[\begin{array}{ccc}
x  & y & z   \\
 0 & x & y \\
 0 & 0 & x    
\end{array}\right]: (x,y,z)\in \mathbb{R}^3   \right\}& x^3 & x\, e^{ \beta\frac{y}{x}  +\gamma \frac{z}{x} -\gamma \frac{ 1}{2}\left(\frac{y}{x}\right)^2 } &  \left[\begin{array}{ccc}
\alpha  & \beta & \gamma   \\
 \beta & \gamma & 0 \\
 \gamma & 0 & 0    
\end{array}\right]  \\ \\ \left\{ \left[
\begin{array}{ccccc}
x & y & w   \\
0 & x & z \\
0 & 0 & x  
\end{array}
\right] :  (x,y,z,w)\in\mathbb{R}^4 \right\} & x^4 & x\, e^{\beta\frac{y}{x} + \gamma\frac{z}{x} +\delta\frac{w}{x}- \frac{\delta}{2}\left(\frac{yz}{x^2}\right)}
& \left[
\begin{array}{cccc}
\alpha & \beta & \gamma & \delta   \\
\beta & 0 & \frac{1}{2}\delta & 0 \\
\gamma & \frac{1}{2}\delta &  0 & 0   \\
\delta & 0 &  0 & 0
\end{array}
\right]\\ \\
\left\{  \left[\begin{array}{ccc}
x  & z & w   \\
 0 & x & v \\
 0 & 0 & y    
\end{array}\right]:  (x,y,z,v,w)\in\mathbb{R}^5 \right\}  & x^4y & (xy)^{\frac{1}{2}}\left(\frac{x}{y}\right)^{\frac{\sigma}{2}}  e^{\frac{\beta}{2}  \frac{z}{x}}  & \left[\begin{array}{ccccc}
\sigma_1  & 0 & 0 & \beta & 0   \\
 0  & \sigma_2 & 0 & 0 & 0  \\
0   & 0 & 0 & 0 & 0 \\
 \beta  & 0 & 0 & 0 & 0  \\
 0  & 0 & 0 & 0 & 0    
\end{array}\right]  \\  
\mathbb{T}_n & x_1^n & x_1 e^{P_{\gamma_2,\dots,\gamma_n}\left(\frac{x_2}{x_1},\dots,\frac{x_n}{x_1}\right)} & \slashed{\mathbf{1}}_n[\mathbb{T}_n]_\text{left} & \bigskip \end{array} 
$

\bigskip\bigskip


\noindent {\bf Table III.} {\small Miscellaneous examples.}

\newpage

\begin{itemize}
 \item There is an epimorphism from $A_{10}$ to $\mathbb{R}\oplus\mathbb{R}$ associated with the ideal $\left\{\left[\begin{array}{ccc}
 0 & z   \\ 
 0 & 0    
\end{array}\right]:z\in \mathbb{R}\right\}$.  The functor defined in Section 2.6 implies that there must then exist a morphism from $\mathbf{u}_{A_{10}}$ and $\mathbf{u}_{\mathbb{R}\oplus\mathbb{R}}$.  Indeed, existence of the latter morphism is obvious from the last entries in rows 10 and 4 of \mbox{Table III}.
 \item There is an epimorphism from $A_{13}$ to $\mathbb{R}\oplus\mathbb{R}$ associated with the ideal

 $\left\{\left[\begin{array}{ccc}
0  & z & w   \\
 0 & 0 & v \\
 0 & 0 & 0    
\end{array}\right]:z,w,v\in\mathbb{R}\right\}$. The functor then implies that there must exist a morphism from $\mathbf{u}_{A_{13}}$ and $\mathbf{u}_{\mathbb{R}\oplus\mathbb{R}}$.  Indeed, existence of the latter morphism is obvious from the last entries in row 13 and row 4.  

On the other hand, Corollary \ref{8/17/22.5} indicates that there can be no epimorphism between $A_{13}$ and its subalgebra $A_{12}$ based on the dimensions of $\mathbf{u}_{A_{13}}$ and $\mathbf{u}_{A_{12}}$, indicated in the last entries of row 13 and row 12.  This demonstration does not employ an ideal-based argument.
 \item 
There is an epimorphism from $A_{12}$ to  $\mathbb{D}$ associated with the ideal

 $\left\{\left[\begin{array}{ccc}
0  & 0 & w   \\
 0 & 0 & v \\
 0 & 0 & 0    
\end{array}\right]:w,v\in\mathbb{R}\right\}$.
The functor then implies that there must exist a morphism from $\mathbf{u}_{A_{12}}$ to $\mathbf{u}_{\mathbb{D}}$.  Indeed, existence of the latter morphism is obvious from the last entries in rows 12 and 5.
 \item There is an epimorphism from $A_{11}$ to $\mathbb{D}$ associated with the ideal 
 $$ \hspace{-5cm} \left\{\left[\begin{array}{ccc}
0  & 0 & w   \\
 0 & 0 & 0 \\
 0 & 0 & 0    
\end{array}\right]:w\in\mathbb{R} \right\} .$$ Accordingly, there must exist a morphism between $\mathbf{u}_{A_{11}}$ and $\mathbf{u}_{\mathbb{D}}$  - and it clearly does exist by inspection of the depicted uncurling spaces at the end of rows 11 and 5.  
\item Based on the listed uncurling spaces, one can observe that Corollary \ref{8/17/22.5} excludes an epimorphism from $\mathbb{H}$ to its subalgebra $\mathbb{C}$.  This is an argument that does not directly involve ideals.  

Furthermore, Corollary \ref{8/17/22.1} indicates that $\mathbb{H}$ is simple - an alternative to the ideal-based argument that $\mathbb{H}$ is simple because it can be shown to be a division algebra.  
 \end{itemize}

  The above observations related to Table III all refer to the uncurling metric component, $\mathbf{u}$, of the inversion characterization group.  However, the unital norm component, $\mathbf{U}$, can also be applied via Corollary \ref{7/25/24.1}.  For example, a glance at the uncurling metrics for $\slashed{\mathbb{C}}$ versus those of $\mathbb{R}\oplus\mathbb{R}$ fails to immediately indicate that these algebras are intimately related.   In contrast, a glance at the special unital norm families for these algebras immediately suggests an intimate relationship between the two algebras. Conversely, a glance at the listed special unital norm families for $\mathbb{C}$, $\slashed{\mathbb{C}}$, and $\mathbb{D}$ immediately
indicates that they cannot be isomorphic.  This kind of observation can be expected to be much more useful in higher dimensions.


\subsection{Further classification considerations} As was remarked at the beginning of Section 2, we have resolved to reduce the algebra bilinear binary product operation to the implied nonlinear unary multiplicative inversion operations.  It was suggested in Section 2.1 (based on Neumann series) that nothing would be gained by also considering the other implied nonlinear unary operations (so we have proceeded from Definition \ref{7/21/23.1}). In fact, there is some practical evidence in support of this suggestion (and more).  Specifically, additional uncurling spaces can be defined by replacing $s^{-1}$ with $s^j$, for $j\in\mathbb{Z}$, $j\ne 0,1$ in Definition \ref{7/21/23.1}, each referred to as ``type $j$".  For each of the algebras listed in the first column of Table III, these variant types of uncurling spaces are easily computed by hand when  $|j|$ is ``small" (and $n$ is ``small" if applicable).  
Thereby one quickly discovers that, within each algebra, those uncurling spaces of different type are {\it identical}.  Despite the informal Neumann series argument presented in Section 2.1, this is still surprising.  That is, let $\hat{s}$ be the left regular representation of $s$, so that we can write $s^j = (\hat{s})^{j-1}[s]$, where $s$ ranges over the algebra elements in an appropriate neighborhood of $\mathbf{1}$. With respect to the basis defined by the columns of an uncurling metric $L$, the above behavior implies that the various vector fields $s\rightarrow s^j = (\hat{s})^{j-1}[s]$ are simultaneously irrotational, at least for small $|j|$ and $n$.  Now consider the vector field $F(s)[s]$ where $F$ is an \mbox{$(n\times n)$-matrix} whose entries are smooth functions of $s$, and suppose $F(s)[s]$ is \mbox{irrotational}.  In general, this does {\it not} imply that $(F(s))^j[s]$ is irrotational for $j\ne 0,1$ (e.g., for \mbox{$n=2$}, $s=\left(\begin{array}{ccc}
 x   \\ 
y    
\end{array}\right)$, and $j=2$, let $F= \left[\begin{array}{ccc}
 f(x) & 1   \\ 
 1 & g(y)    
\end{array}\right]$ with nonconstant $f(x)$, $g(y)$).  So there is something special about the map defined by $\hat{s}$ that results in the unexpected behavior cited above. 
In practical terms, that behavior is of additional interest because it impacts the potential complexity of our general approach (i.e., the desirability of computing more than one uncurling space type).  This important consideration is identified as a matter for future investigation. At this point we are simply left with,
\begin{conjecture} For a given algebra, the uncurling spaces of all types $j\ne 0,1$ are identical. 
\end{conjecture}
\noindent The conjecture is a particularly emphatic assertion that unary operations defined by nontrivial powers each carry the same information regarding an algebra's structure.  
 
 \section{The case of nonunital algebras} Nonunital algebras can be addressed by modifying Definition \ref{7/21/23.1} by removing the reference to an algebra's units, and replacing $s^{-1}$ with $s^j$, $j \in \mathbb{N}$.   However, unital norms are no longer germane.  That is, integration of $Ls^j$, $j>0$ does not have the interest that integration of $Ls^{-1}$ has in the case of unital algebras (there is nothing comparable to the unital norm group).  Nevertheless, one would at least be left with uncurling metric groups.

For nilpotent algebras, the uncurling spaces of different types will be different, something that appears not to be the case for unital algebras.  For example, consider the real three-dimensional nonunital algebra composed of nilpotent elements $(x,y,z)$ that multiply according to the representation
  $ \left[
\begin{array}{ccccc}
0 & x & y & z  \\
0 & 0 & x & y \\
0 & 0 & 0 & x \\
0 & 0 & 0 & 0
\end{array}
\right] $.  We then have $(x,y,z)^2 = (0,x^2, 2xy)$.  The uncurling space resulting from replacement of $s^{-1}$ with  $s^2$ in Definition \ref{7/21/23.1} is then $\left[
\begin{array}{ccccc}
\alpha & \beta & \gamma   \\
\beta & \gamma & 0 \\
\gamma & 0 & 0  
\end{array}
\right] $.  On the other hand, substitution of $s^{-1}$ with $s^3$  deals with $(x,y,z)^3 = (0,0, x^3)$, and the associated uncurling space becomes $\left[
\begin{array}{ccccc}
\alpha & \beta & \gamma   \\
\beta & \delta & 0 \\
\gamma & 0 & 0  
\end{array}
\right] $.  Substitution of $s^{-1}$ by $s^j$, $j>3$, results in the additive group of all symmetric $(3\times 3)$-matrices since $s^j$ is zero for all algebra elements.  The uncurling spaces associated with $s^2$, $s^3$, and $s^4$, must differ because the respective dimensions of the ranges of these operations differ - and this type of behavior will be typical with nilpotent algebras.  
 

\vspace{.5cm}

\appendix

 \section{First motivation: a  new proof of the Pythagorean Theorem}    The isotropy of the Newtonian concept of space, asserting that physical laws are independent of direction (which results in ``conservation of angular momentum" as an example of the renowned ``Noether's Theorem" in theoretical physics), has a non-physical analogue with respect to Euclidean space or, for our present purposes, the Euclidean plane.  Specifically, given a line and a point, the ``Parallel Postulate" (or its equivalent) ultimately indicates that not only is there is a unique line through the given point parallel to the given line, but also that there is a unique line though the given point perpendicular to the given line.  Using the formulation of Euclidean geometry due to George Birkhoff \cite{birkhoff:1932}, which fundamentally incorporates $\mathbb{R}$, it is then easy to derive a bijection between the points of the Euclidean plane and those of $\mathbb{R}\times\mathbb{R}$, using a pair of perpendicular lines intersecting at some chosen point $O$, such that this ``origin" $O$ is mapped to $(0,0)$.  That is, the perpendicular lines through $O$ then become directed ``orthogonal coordinate axes": the real-valued coordinates are due to the assertion of one of Birkhoff's Euclidean geometry postulates that there is a bijection between the points of $\mathbb{R}$ and the points of any line, which also endows the line with a direction as determined by the increasing values of $\mathbb{R}$.  Consequently, this formulation allows incorporation of concepts of real vector spaces, limits, continuity, and one-dimensional Differential Calculus.  Thereby, the above bijection (which was specifically demonstrated in \cite{birkhoff:1932}) leads to an isometric isomorphism. 
To see this, we specify that $\mathbb{R}\times\mathbb{R}$  inherits the Euclidean length function $\ell(P)$, where the latter gives the length of the line segment with endpoints $O$ and $P$, where $P$ is  any point in the Euclidean plane.  Accordingly, we specify that $\ell(x,y) \equiv \ell(P)$ for $(x,y)$ the image of $P$ under the bijection.  Furthermore, the Euclidean plane with distinguished origin $O$ is easily made into a real vector space (dilation/reflection of directed line segments based at $O$ result from multiplication by members of $\mathbb{R}$, and the sum of directed line segments based at $O$ follow from parallelogram constructions).   Since $\mathbb{R}\oplus\mathbb{R}$ is already a real vector space, and both it and the latter Euclidean plane vector space are endowed with the above length function (a degree-1 positive homogeneous function, as implied by one of the geometry postulates), the isometric isomorphism follows without difficulty (this being simply a formalization of Descartes' conception).  An important feature is that (due to the Parallel Postulate or equivalent) such an isomorphism can be constructed from {\it any} pair of perpendicular lines through $O$, i.e., regardless of the angle one of the perpendiculars makes with an arbitrarily selected line through $O$.  So this is the isotropic feature of the Euclidean plane referred to at the beginning of this appendix.    
 
 We will now present a new proof of the Pythagorean Theorem based on that isotropy (fully detailed in \cite{fgreensite:2022b}), and it will be seen that features of the derivation motivate Definition \ref{7/21/23.1}.
 
 First,  we have the usual directional derivative of the inherited length function on $\mathbb{R}\oplus\mathbb{R}$, denoted $D\ell(x,y;u,v)$, this being the value of what is essentially a one-dimensional derivative of the length function at $(x,y)$ with respect to the direction vector $(u,v)\ne (0,0)$ (i.e., employing the limit of an expression involving the usual vector space operations on a one-dimensional subset of $\mathbb{R}\oplus\mathbb{R}$).  For $(x,y)\ne (0,0)$, this derivative value can be shown to exist based on the triangle inequality. The analogous manipulations in the above Euclidean plane vector space allow computation of a derivative similarly notated as $D\ell(P;C)$.  If $P$ and $C$ are the respective pre-images of $(x,y)$ and $(u,v)$ under the isomorphism resulting from a set of orthogonal coordinate axes in the Euclidean plane vector space, then $D\ell(P;C) = D\ell(x,y;u,v)$ (because the isomorphism is isometric).  With respect to a given direction vector, it can be shown that $\ell(x,y)$ is continuously differentiable except at the origin (this being the only subtlety of the proof, since the two dimensional continuity of the derivative must be established in the absence of an expression for the length of an arbitrary line segment; the continuity demonstration ultimately results from the degree-1 positive homogeneity of the length function, the latter essentially being an expression of one of the Euclidean geometry postulates).  Using the Mean Value Theorem of one-dimensional Calculus, it is then easy to show \begin{equation}\label{11/30/23.1} D\ell(x,y;u,v) = \frac{\partial \ell(x,y)}{\partial x}u + \frac{\partial \ell(x,y)}{\partial y}v,\end{equation} where $\frac{\partial \ell(x,y)}{\partial x}$ and $\frac{\partial \ell(x,y)}{\partial y}$ are the directional derivatives of $\ell$ at $(x,y)\ne(0,0)$, specified by the direction vectors $(1,0)$ and $(0,1)$ on the respective coordinate axes in $\mathbb{R}\oplus\mathbb{R}$, where the latter direction vectors are the images of unit length direction vectors associated with a chosen pair of orthogonal coordinate axes in the Euclidean plane vector space that result in a first isometric isomorphism as indicated above.  The latter coordinate axes are respectively labeled as the $x$-axis and the $y$-axis. We conclude from (\ref{11/30/23.1}) that, for fixed $(x,y)\ne (0,0)$,  $D\ell(x,y;u,v)$ is a linear functional, and is evidently the dual of $\left(\frac{\partial\ell(x,y)}{\partial x}, \frac{\partial\ell(x,y)}{\partial y}\right)$.  From that first isomorphism, the vector expression at the end of the prior sentence labels a point $G$ of the Euclidean plane vector space.  An alternative orthogonal coordinate axis system through the origin in the Euclidean plane, resulting in a second isometric isomorphism with $\mathbb{R}\oplus\mathbb{R}$, will provide a different label for this {\it same} point $G$ following computation of a linear functional analogous to (\ref{11/30/23.1}) expressed via that second isomorphism, because each orthogonal coordinate system provides an isometric isomorphism with $\mathbb{R}\oplus\mathbb{R}$ thereby allowing computation of the directional derivative $D\ell(P;C)$ for the Euclidean plane vector space.  In fact, we can find a set of alternative orthogonal coordinate axes, denoted as the $x'$-axis and $y'$-axis, such that $P$ is on the positive portion of the $x'$-axis - so that the associated isomorphism now labels $P$ as $(x',0)$ instead of $(x,y)$.  The Euclidean vector space point $G$ is now labeled by  
 \begin{equation}\label{11/24/23.4} \left( \frac{\partial \ell(x',0)}{\partial x'}, \frac{\partial \ell(x',0)}{\partial y'}\right) = (1,0),\end{equation} where the right-hand-side is an easy calculation since determination of the two directional derivatives on the left-hand-side only involves evaluation of $\ell$ on the coordinate axes.  Equation (\ref{11/24/23.4}) states that $ \left( \frac{\partial \ell(x',0)}{\partial x'}, \frac{\partial \ell(x',0)}{\partial y'}\right)$ is the point of the unit circle on the ray from $(0,0)$ that includes the point $(x',0)$.
 Based on this second isomorphism, this means that $G$ is the point of the unit circle on the ray from $O$ that includes the point $P$.  Hence, based on the first isomorphism, this requires $\left( \frac{\partial \ell(x,y)}{\partial x}, \frac{\partial \ell(x,y)}{\partial y}\right)$ to be the point of the unit circle on the ray from $(0,0)$ that goes through the point $(x,y)$. Consequently, \begin{equation}\label{11/25/23.1} \ell\left( \frac{\partial \ell(x,y)}{\partial x}, \frac{\partial \ell(x,y)}{\partial y}\right) = 1.\end{equation} and, taking into account that $\ell$ is a degree-1 positive homogeneous function,  \begin{equation}\label{11/24/23.5} \left( \frac{\partial \ell(x,y)}{\partial x}, \frac{\partial \ell(x,y)}{\partial y}\right) = \left( \frac{x}{\ell(x,y)},  \frac{y}{\ell(x,y)}\right),\end{equation} since the right-hand-side of (\ref{11/24/23.5}) is (uniquely) both a point on the unit circle while also being on the ray from $(0,0)$ that includes $(x,y)$ - as are now seen to be required for a labeling of $G$ via an isometric isomorphism. 
 
  As regards (\ref{11/24/23.5}), one-dimensional indefinite integration (anti-differentiation) can be applied to the equality of the first coordinates, and then one-dimensional indefinite integration can be independently applied to the equality of the second coordinates.  Reconciliation of the  two resulting expressions for $\ell(x,y)$ in the context of (\ref{11/25/23.1}) and the nonnegativity of $\ell$ quickly leads to $\ell(x,y)=\sqrt{x^2+y^2}$, from which the Pythagorean Theorem easily follows via an argument involving parallelograms to measure lengths of arbitrary line segments using $\ell$. 
   The full details of this proof are found in \cite{fgreensite:2022b}.

  
So now we begin a series of easy generalizations.  

Equations (\ref{11/24/23.5}) and (\ref{11/25/23.1}) can be recast as, \begin{equation}\label{11/29/23.5} s = \ell(s)\nabla\ell(s),\end{equation} and\begin{equation}\label{11/29/23.6}\ell\left(\nabla\ell(s)\right) = 1,\end{equation} now generalized to $s\in \mathbb{R}^n$, $n\ge 2$, where $\nabla\ell(s)$ is simply notation for the $n$-tuple of coordinate-wise derivatives with respect to the standard basis, as on the left-hand-sides of  (\ref{11/25/23.1}) and (\ref{11/24/23.5}).  For nonnegative $\ell(s)$, the solution to (\ref{11/29/23.5}) and (\ref{11/29/23.6}), is similarly determined to be the Euclidean norm.  

But why stop there?  Equations (\ref{11/29/23.5}) and (\ref{11/29/23.6}) are amenable to further generalization, from which  we can obtain {\it any} quadratic space, via the norm $\ell(s)$ \mbox{specified by} \begin{equation}\label{11/29/23.1} Ls = \ell(s)\nabla\ell(s),\end{equation} and \begin{equation}\label{11/29/23.2}\ell\left(L^{-1}[\nabla\ell(s)]\right) = 1,\end{equation} for any appropriately-dimensioned nonsingular symmetric $L$.  The matrix $L$ must be symmetric since the right-hand-side of (\ref{11/29/23.1}) is the gradient of $\frac{1}{2}\ell^2(s)$, so $Ls$ is required to be a gradient.  Equation (\ref{11/29/23.2}) follows from dividing both sides of (\ref{11/29/23.1}) by $\ell(s)$, then applying $L^{-1}$ to both sides of the resulting equation, then applying $\ell$ to both sides of that resulting equation, and finally invoking the assumed degree-1 positive homogeneity of $\ell(s)$. Again,  $\nabla\ell(s)$ is simply notation for the ordered $n$-tuple of  coordinate-wise one-dimensional derivatives, with the standard basis of $\mathbb{R}^n$ assumed.
  
Before getting into generalization of the above to algebras, let us examine the things a bit further.  Equations (\ref{11/29/23.1}) and (\ref{11/29/23.2}) are feasible because $Ls$ is a gradient and the equations continue to hold when $s$ is replaced by $\alpha s$, $\alpha >0$.  This is a consequence of the Euler Homogeneous Function Theorem, which states that if $\ell(s)$ is a degree-1 positive homogeneous function then
  $s\cdot\nabla \ell(s) = \ell(s)$.  For $\alpha>0$, replacing $s$ by $\alpha s$ then leads to \[ \alpha s \cdot \nabla \ell(\alpha s) = \ell(\alpha s) = \alpha \ell(s) = \alpha s\cdot \nabla \ell(s),\] where $\nabla \ell(\alpha s)$ means that $\nabla \ell$ is evaluated at $\alpha s$. Equality of the above left-hand-side expression with the above right-hand-side expression implies, \begin{equation}\label{11/30/23.5} \nabla \ell(\alpha s) = \nabla \ell(s).\end{equation}    
  Given that we require a degree-1 positive homogeneous $\ell(s)$, this means that the left-hand-side of (\ref{11/29/23.1}) and the right-hand-side of (\ref{11/29/23.1}) are consistent when $s$ is replaced by $\alpha s$, $\alpha>0$ (each side simply being multiplied by $\alpha$).  
  
  With that in mind, suppose we are now given a real unital associative algebra with $\mathbb{R}^n$ as its vector space of elements  and, inspired by the above norms on vector spaces, we wish to look for what will ultimately be a degree-1 positive homogeneous function $\ell(s)$ that somehow also acknowledges the algebra.  Along those lines, we recall that the ``vector-space-appropriate" (\ref{11/29/23.1}) exhibits the feature that both sides of the equation are multiplied by a common factor when $s$ is replaced by a multiple of $s$.  So, we want an ``algebra-appropriate" equation that looks similar to (\ref{11/29/23.1}), and retains the feature noted in the prior sentence.  An obvious candidate is simply, \begin{equation}\label{12/22/23.1} Ls^{-1}=\frac{1}{\ell(s)}\nabla\ell(s),\end{equation} since the above looks a lot like (\ref{11/29/23.1}), it addresses the algebra through its use of the nonlinear unary operation $s^{-1}$ and, assuming we retain (\ref{11/29/23.2}) and can ultimately identify a degree-1 positive homogeneous solution, replacement of $s$ by $\alpha s$ leads to both sides of (\ref{12/22/23.1}) being multiplied by a common factor (this time the common factor is $\frac{1}{\alpha}$ instead of $\alpha$).   
  
In fact, assuming $L$ to be nonsingular and then applying $L^{-1}$ to both sides of (\ref{12/22/23.1}), then applying $\ell$ to both sides of the resulting equation, and then invoking the assumed degree-1 positive homogeneity of $\ell$ and the resulting \begin{equation}\tag{\ref{11/29/23.2}}\ell\left(L^{-1}[\nabla\ell(s)]\right) = 1,\end{equation} we find that $\ell(s)$ will also respect the algebra's multiplicative inversion operation, \begin{equation}\tag{\ref{5/23/23.3}} \ell(s^{-1}) = (\ell(s))^{-1} = \frac{1}{\ell(s)}.\end{equation}    Thus, (\ref{12/22/23.1}) becomes \begin{equation}\label{11/27/23.7} Ls^{-1} = \ell(s^{-1})\nabla\ell(s).\end{equation}    Observe that this pairing of (\ref{11/29/23.2}) and (\ref{11/27/23.7}) in the algebra context emulates the pairing (\ref{11/29/23.1}) and (\ref{11/29/23.2}) in the inner product space context.
  
  Hence, we are motivated to find $\ell(s)$ satisfying (\ref{11/27/23.7}) and (\ref{11/29/23.2}). 
  However, there is a new constraint necessarily introduced in this extension to unital algebras.  We began with (\ref{12/22/23.1}), whose right-hand-side is $\nabla\log\ell(s)$, i.e., a gradient.  In general, $s^{-1}$ is not a gradient, nor is $Ls^{-1}$ with an arbitrary symmetric $L$. Consequently, to proceed as above, it is now required that $L$ satisfies \[ d([Ls^{-1}]\cdot {\sf d}s) = 0.\]  It might be speculated that, for any unital algebra with vector space space of elements $\mathbb{R}^n$, the set of $L$ satisfying this constraint always has nontrivial members, and the set often distinguishes different algebras from each other.  It might also be speculated that, given a member $L$ of this set,  the coupled equations, (\ref{11/29/23.2}) and (\ref{11/27/23.7}), emulating (\ref{11/25/23.1}) and (\ref{11/24/23.5}), always has a solution if $L$ is nonsingular.
  These speculations motivate Definition \ref{7/21/23.1}, the beginning of our program.  This was one of the points made in  Remark \ref{11/21/23.1}.  Of course, (\ref{11/27/23.7}) and (\ref{11/29/23.2}) are essentially examples of (\ref{5/22/24.5}) and (\ref{5/22/24.6}). 

But this process can be reversed, which was the other point made in Remark \ref{11/21/23.1}.  Consider the trivial case of Definition \ref{4/7/23.1} where $j=1$.  A type 1 uncurling metric $L$ satisfies $d([Ls]\cdot {\sf d}s)=0$, and thus can be any symmetric matrix of appropriate dimension.  Given the equation in the prior sentence, and assuming we are looking for nonnegative $\ell(s)$, we are entitled to write \begin{equation}\tag{\ref{11/29/23.1}} Ls = \ell(s)\nabla\ell(s),\end{equation} since the right-hand-side is a gradient, i.e., $\nabla(\frac{1}{2}\ell^2(s))$.  
Taking $L$ to be nonsingular, and mandating that $\ell(s)$ be degree-1 positive homogeneous, and then employing the same type of manipulations as used at the beginning of the penultimate paragraph, we obtain \begin{equation}\tag{\ref{11/29/23.2}}\ell\left(L^{-1}[\nabla\ell(s)]\right) = 1,\end{equation} along with $ s = \ell(s)\left(L^{-1}[\nabla\ell(s)]\right)$.  
This format is attractive since $s$ is now expressed as the product of its ``length" and its ``direction" (a ``unit length" vector according to ({\ref{11/29/23.2})).  If we mandate that $\ell(s)$ be {\it isotropic} (i.e., the length of a vector is independent of direction), then it is easy to show that $L$ must be the identity matrix, and we are thus presented with (\ref{11/29/23.5}), (\ref{11/29/23.6}), whose solution is the Euclidean norm. So it now occurs to us that if we can derive these latter two equations from the postulates of Euclidean geometry (as formulated by by George Birkhoff \cite{birkhoff:1932}), we would then have a proof of the Pythagorean Theorem.  Based on the above, such a proof would likely highlight the concept of isotropy. This was the motivation for our new proof of the Pythagorean Theorem via the derivation of (\ref{11/24/23.5}), (\ref{11/25/23.1}).
  
  \bigskip\bigskip
  
  \section{Second motivation: an unusual regularization method}
     Appendix A inverts the reality that, rather than motivating our program, the forwarded new proof of the Pythagorean Theorem was instead itself motivated by our program (see Remark \ref{11/21/23.1}). In this Appendix B, we will present what actually led to the introduction of uncurling metrics and unital norms.  Instead of relying on Euclidean geometry and some Calculus as in Appendix A, this motivation relies on the interplay of deterministic and statistical interpretations of an unusual regularization method in the realm of Inverse Problems (essentially, a geometrical interpretation of a statistically-inspired method) \cite{fgreensite:2022d}.
  
    
 Consider the linear problem described by, \begin{equation}\label{7/7/20.9}  y^\delta=Fx + \nu,\end{equation} where one is given a data vector $y^\delta$ that is a version of an unknown source vector $x$ as transformed by the known linear operator $F:\mathbb{R}^m\rightarrow\mathbb{R}^n$, but where the data vector has been corrupted by a ``noise" vector $\nu$ that is only characterized by   \begin{equation}\label{7/12/20.1} \| y^\delta - Fx\|^2 \le \delta^2, \end{equation}  where $\delta$ is given and $\|\cdot \|$ is the Euclidean norm.  
    
If $F$ is an ill-conditioned matrix (e.g., a discretization of a linear compact operator), a solution estimate $F^{-1}y^\delta$ or $F^\dagger y^\delta$ is usually worthless ($F^\dagger$ is the pseudoinverse of $F$).  That is, consider the singular system for $F$ as  $\{u_i ,\sigma_i, v_i\}$, where $\{u_i\}$ is the set of eigenvectors of $FF'$, $\{v_i\}$ is the set of eigenvectors of $F'F$, and $\{\sigma_i^2 \}$ is the set of eigenvalues of $F'F$ (i.e., $\{\sigma_i \}$ is a set of singular values).   The least squares solution estimate of smallest norm is then, \begin{equation}\label{7/7/20.8} x_{\rm poor} = F^\dagger y^\delta = \sum_{ \{i: \sigma_i \ne 0\}} \frac{u_i' y^\delta}{\sigma_i}v_i,\end{equation} where  an expression like $a' b$ indicates the Euclidean inner product of vectors $a$ and $b$.   It is evident that the existence of small singular values  leads to potentially extreme noise amplification, so that the above is a poor estimate of $x$ as regards the at least possibly approachable $x^\dagger \equiv F^\dagger F x$.
 Instead, a common treatment (upon which zero-order Tikhonov regularization with the ``Discrepancy Principle" is based \cite{vogel:2002,engl:2000}) is to supply the solution estimate of minimum norm that is consistent with the modification of (\ref{7/12/20.1}) as \begin{equation}\label{7/7/20.3} \| y^\delta - Fx\|^2 =\delta^2 .\end{equation}   That is, the solution estimate is the point of smallest Euclidean norm on the ``discrepancy ellipsoid",   \begin{equation}\label{7/15/20.1} \partial\Omega_\delta \equiv \{ x\in \mathcal{D}(F): \| y^\delta - Fx\|^2 =\delta^2\},\end{equation} where $\mathcal{D}(F)$ is the domain of $F$.  We use $\mathcal{D}(F)$ here instead of $\mathbb{R}^m$ because we will want (\ref{7/15/20.1}) to also be relevant when $F$ is a compact operator (the general Inverse Problems setting \cite{engl:2000}).
  A Lagrange multiplier argument then gives the solution estimate as, \begin{equation}\label{7/8/20.9} x^\delta_D = (F'F + \gamma I)^{-1}F'y^\delta,\end{equation} where $I$ is the identity operator, and regularization parameter $\gamma$ is the reciprocal of the
   Lagrange multiplier as determined by the requirement that (\ref{7/7/20.3}) be satisfied (the subscript ``$_D$" indicates the estimate follows from a ``deterministic" treatment).  
 
 Alternatively, one may consider both $x$ and $\nu$ to be zero-mean Gaussian random vectors (``signal" and ``noise", respectively), with the covariance matrix of $x$ being $C_x$ (assumed for the moment to be nonsingular), and the covariance matrix of noise vector $\nu$ being $\epsilon^2 I$ ($\delta^2$ in (\ref{7/7/20.3}) might now be viewed as an estimate of the sum of the autocovariances of the components of $\nu$).  The {\it maximum  a posteriori} (Bayesian) estimate  is then, \begin{equation}\label{7/8/20.5} x^\delta_S = (F'F + \epsilon^2 C_x^{-1})^{-1} F'y^\delta,\end{equation}  where the subscript ``$_S$" indicates that the estimate follows from  a ``statistical" treatment \cite{vogel:2002}. 
  
 We can thus associate the deterministic solution estimate (\ref{7/8/20.9}) with a statistical interpretation (\ref{7/8/20.5}), as resulting from assigned signal and noise covariance matrices each proportional to the identity matrix, with $\gamma$ as the inverse of the square of the  ``signal-to-noise" ratio (this inverse square being the ratio of the respective noise and signal component variances).  Thus, the diagonal entries of the noise covariance matrix are assumed to be identical, and these represent the noise variance $\epsilon^2$ in each of the components of the data vector (i.e., the variances of each component of random vector $\nu$).   So the full noise covariance matrix is assumed known {\it a priori}, perhaps based on knowledge of $\delta^2$ under the supposition that $\delta^2 \approx n\epsilon^2$ in (\ref{7/7/20.3}), where $n$ is the dimension of the range of $F$.  The factor multiplying the identity matrix to give $C_x$, so that this factor is the (assumed uniform) variance of each component of the signal random vector $x$, then follows once $\gamma$ is computed.  Hence, from a statistical standpoint,  (\ref{7/7/20.3}), (\ref{7/8/20.9}) can be re-interpreted as implying a method resulting in {\it simultaneous} estimation of $x$ and the presumably uniform signal component variances $\epsilon^2/\gamma$. 
 
But as long as one is simultaneously estimating the solution and a covariance matrix feature, why stop there?   Instead of assuming that the   signal covariance matrix is proportional to the identity and estimating the proportionality constant simultaneously with the signal, one might only assume that the signal covariance matrix is diagonal in the basis defined by the eigenvectors of $F'F$, and seek to estimate the diagonal elements simultaneously with the signal (i.e., one assumes  that the matrix whose columns are the right singular vectors of  $F$ diagonalizes the signal covariance matrix).
The latter basis, $\{v_i\}$, is the natural one to be used in the context of this type of assumption since, considering the underlying setting where $F$ is a compact operator, one thing we {\it know} is that for the inverse problem to make sense the signal must be in the domain of $F$ (the problem is {\it defined} by $F$ operating on the signal). That is, we already know that the sum of the squares of the signal components in this basis is finite - and (absent other prior information) this is the only basis for which we know that to be true.  So, the assumption is simply that the components of random vector $x$ in this basis are independent, rather than additionally assuming that they all have the same variance (which would correspond to the assumption that the signal covariance matrix is proportional to $I$).

Proceeding in this way, one notes that the value of the $i$-th eigenvalue of $C_x$ is then the autocovariance of the $i$-th component of the signal random vector in the $\{v_i\}$-basis, i.e., the expectation of the square of the $i$-th component of random vector  $x$ in that basis (which is the expectation of $(v_i'x)^2$).   Thus, an alternative approach to the statistical treatment of (\ref{7/7/20.9}), (\ref{7/7/20.3}) (which we do not claim to be optimal in the solitary estimation of $x$)  would be to require that the eigenvalues of $C_x$ be the squares of the components of the solution estimate vector -  inspired by the actual meaning of those covariance matrix values as expectations  of the squares of the signal component random variables in the $\{v_i\}$-basis.  Such a solution estimate would then be  a fixed point of the dynamical system on $\mathbb{R}^m$ defined by iteration of the operator, \begin{equation}\label{10/19/10.3} A_\epsilon[w] \equiv \lim_{\alpha\rightarrow 0}\left(F'F+\epsilon^2\, V\mbox{\rm diag}\{(v_i'w)^2 + \alpha^2\}^{-1}V' \right)^{-1} F'y^\delta, \end{equation} where $V$ is the matrix whose $i$-th column is $v_i$ for each $i$. The role of $\alpha$ and $\lim_{\alpha\rightarrow 0}$ is to account for the possibility that one of the $v_i'w$ is zero.  To get the iterations started,  a solution estimate candidate $w$ (even chosen arbitrarily) is associated with a signal covariance matrix candidate $V \mbox{\rm diag}\{(v _i'w)^2\}V'$.    According to (\ref{7/8/20.5}),  $A_\epsilon[w]$ is then an updated solution estimate using the latter signal covariance matrix.  For the next iteration, the squares of the components of this updated solution estimate  $A_\epsilon[w]$ in the $\{v_i\}$-basis are used to define the entries of an updated version of the diagonal matrix on the right-hand-side of (\ref{10/19/10.3}) (i.e., updating the signal covariance matrix), which then supplies a successive updated solution estimate as $A_\epsilon[A_\epsilon[w]]$, etc. This iteration process can be repeated indefinitely.  A fixed point of the dynamical system evidently supplies a covariance matrix whose eigenvalues are the squares of the solution estimate components in the $\{v_i\}$-basis.  
 The fixed point $x=0$ is associated with a signal covariance matrix as the matrix of zeros. Thus, one would not select $w=0$ as an initial guess for the solution estimate to begin the iteration process (also, $w=0$ is very far from satisfying (\ref{7/7/20.3})).  
 
 
But it is helpful that one can easily identify the fixed points of $A_\epsilon$. Inserting the previously noted
singular system associated with $F$ into the right-hand-side of (\ref{10/19/10.3}),  we obtain \begin{equation}\label{3/31/12.3} A_\epsilon[w] = \sum_i \frac{\sigma_i (u_i'y^\delta)}{\sigma_i^2 + \epsilon^2/(v_i'w)^2}v_i .\end{equation} A fixed point $p$ must satisfy  \begin{equation}\label{4/1/12.1} A_\epsilon[p] = \sum_i  \frac{\sigma_i (u_i'y^\delta)}{\sigma_i^2 + \epsilon^2/(v_i'p)^2}v_i = p = \sum_i  (v_i' p)v_i.\end{equation}  Since the coefficients of $v_i$ in the two sums in (\ref{4/1/12.1}) are necessarily equal, for $v_i'p\ne 0$ we obtain \begin{equation}\label{7/11/10.1} \sigma_iv_i'p = \frac{u_i' y^\delta \pm \sqrt{(u_i'y^\delta)^2 -4\epsilon^2}}{2}.\end{equation}  In the first place, this means that the series for a fixed point will only contain terms $i$ such that $(u_i'y^\delta)^2 >4\epsilon^2$.  Second, using Taylor series, it is easy to show that if the sign in front of the radical on the right-hand-side above is not the same as the sign of $u_i'y^\delta$, then $|\sigma_i v_i'p| =  O(\epsilon^2/|u_i'y^\delta|)$.  Indeed, whatever signal
 was contained in $u_i'y^\delta$ will have then been mostly 
 removed from the right-hand-side of (\ref{7/11/10.1}) (at least until $|u_i'y^\delta|$ is on the order of $\epsilon$, at which point in practical settings the data is mostly noise-dominated).   Thus, any fixed point of interest will be such that the sign in front of the radical is the same as the sign of $u_i'y^\delta$ .   The solution estimate that takes advantage of the greatest number of data subspaces available to this method is then,
\begin{equation}\label{7/19/10.1} p = \frac{1}{2}\sum_{\{i:\,|u_i'y^\delta|>2\epsilon\}} \left(u_i' y^\delta + {\rm sgn}(u_i' y^\delta) \sqrt{(u_i'y^\delta)^2 -4\epsilon^2}\right)\frac{v_i}{\sigma_i}.\end{equation}      
Note that the fixed points of of $A_\epsilon$ are truncations - a feature this solution estimate shares with the (quite different)  commonly used Truncated Singular Value Decomposition (TSVD) solution estimate \cite{vogel:2002}.  

The following result answers an obvious question of interest, though it is tangental to our objective in this appendix.  The proof is found in \cite{fgreensite:2022d}.


\begin{lemma} \label{7/6/21.9} Consider the solution estimate (\ref{7/19/10.1}) but with removal of the terms of the sum on the right-hand-side for which $|u_i'y^\delta|<4\epsilon$.  This is an attracting fixed point of the dynamical system resulting from repeated iterations of $A_\epsilon$. \end{lemma}



For the purposes of this appendix, the point in presenting all this is that the solution estimate $p$ given by (\ref{7/19/10.1}) is a critical point of the {\it geometric mean} of the solution estimate component absolute values on an implied ellipsoid  derived from the discrepancy ellipsoid on which $p$ lies - in contrast to the most commonly used solution estimate, which is a critical point of the Euclidean norm, i.e., a critical point of the {\it quadratic mean}, on the discrepancy ellipsoid on which it lies.  That is, the geometric mean has both a statistical and deterministic interpretation in the context of the regularization of inverse problems, just as the Euclidean norm does. 
 
 
 \begin{theorem} \label{7/6/21.6} Consider \begin{equation} \tag{\ref{7/7/20.9}} y^\delta = F x +\nu, \end{equation} where $y^\delta$ is known, $F:\mathbb{R}^k\rightarrow\mathbb{R}^k$ is a given nonsingular linear transformation associated with singular value decomposition matrix $U\Sigma V'$, and $\{u_i\}$ and $\{v_i\}$ are the sets of columns of $\mathbf{U}$ and $V$, respectively.  Further consider the dynamical system on $\mathbb{R}^k$ defined by iteration of the operator, \begin{equation}\tag{\ref{10/19/10.3}} A_\epsilon[w] \equiv \lim_{\alpha\rightarrow 0}\left(F'F+\epsilon^2\, V\mbox{\rm diag}\{(v_i'w)^2 + \alpha^2\}^{-1}V' \right)^{-1} F'y^\delta. \end{equation} 
As a solution estimate for $x$, the fixed point $p_\epsilon$ of largest Euclidean norm of this dynamical system  has the following interpretations:
\begin{enumerate}\item {\rm Statistical:} If $\epsilon^2$ is the variance of zero-mean white Gaussian noise random vector $\nu$, then 
$p_\epsilon$ is the  {\rm maximum a posteriori} estimate for zero-mean Gaussian signal random vector $x$ of greatest Euclidean norm
 that results from a signal covariance matrix  whose eigenvectors are the members of the basis $\{v_i\}$ and with the additional feature that its eigenvalues are the squares of the components of the solution estimate with respect to those basis vectors.   
 \item  
{\rm Deterministic:} Denote the discrepancy of $p_\epsilon$ as $\delta_\epsilon \equiv \| y^\delta - Fp_\epsilon\|$, which may be adjusted by varying $\epsilon$ so as to allow for the influence of prior knowledge concerning the magnitude of unknown noise vector $\nu$.  Let $\mathbb{V}_\epsilon \subset \mathbb{R}^k$ be the space spanned by the subset of eigenvectors $\{v_i: |u_i'y^\delta| > 2\epsilon \}$.  
 Consider the ellipsoid \[ \partial\Omega \equiv \{ s\in \mathbb{V}_\epsilon: \| y^\delta - Fs\|^2 = \delta_\epsilon^2 \}.\]   Then $p_\epsilon\in \mathbb{V}_\epsilon$ is a critical point of the geometric mean of the absolute values of the coordinates of the points of $\partial\Omega$ with respect to the above eigenvector subset basis.
\end{enumerate}
\end{theorem}

\begin{proof}
The statistical interpretation follows from the prior observations related to the form of the right-hand-side of (\ref{7/8/20.5}) and inspection of the form of (\ref{10/19/10.3}) as regards a fixed point.

We now consider the deterministic interpretation.   
To simplify our notation, without loss we assume an arrangement of the sequence of right singular vectors of $F$ such that $\{v_i\}_{i=1}^n$ is the set $\{v_i: |u_i'y^\delta| > 2\epsilon \}$, and we will write $p_\epsilon$ as simply $p$.  We observe that $p$ is a member of $\mathbb{V}_\epsilon$, since we have previously shown that $p$ is given by (\ref{7/19/10.1}). 
Define $\hat{p} \equiv (v_1'p,\dots,v_n'p)$, i.e., $\hat{p}$ is the coordinate expression of $p$ as a member of $\partial\Omega\subset \mathbb{V}_e$ (thus, $\hat{p}$ is $p$ with suppression of the latter's null-valued coordinates with respect to the full basis of $\mathbb{R}^k$ given by $\{v_i\}$). 
Define $C_p \equiv \mbox{diag}\{(v_i'p)^2\}$ and  $C_{\hat{p}} \equiv  \mbox{diag}\{(v_i'p)^2\}_{i=1}^n$.  Since $p$ is a fixed point of $A_\epsilon$,  (\ref{10/19/10.3}) implies \begin{equation} p = \lim_{\alpha\rightarrow 0} \left(F'F+\epsilon^2\,\left[C_p + \mbox{diag}\{\alpha^2\}\right]^{-1}\right)^{-1} F'y^\delta.\label{12/14/20.1}\end{equation} However, the null coordinates of $p$ correspond to the null diagonal values of $C_p$.  There is no intrinsic loss in suppressing these null entries.  {\it Hence, the right-hand-side of (\ref{12/14/20.1}), can be cast as the solution to the variational problem of finding the point minimizing $s'C^{-1}_{\hat{p}} s$ over $z\in \partial\Omega \subset \mathbb{R}^n$, which is $\hat{p}$}.
Specifically, (\ref{12/14/20.1}) implies, and is replaced by, 
\begin{equation} \hat{p} = \left(\hat{F}'\hat{F}+\epsilon^2\,C_{\hat{p}}^{-1}\right)^{-1} \hat{F}'y^\delta, \label{12/19/20.1}\end{equation} where $\hat{F}$ is the operator resulting from the restriction of the domain of $F$ to $\mathbb{V}_\epsilon$.
  According to the principle of Lagrange multipliers, the point minimizing $s'C^{-1}_{\hat{p}} s$ on \[ \partial\Omega = \{s\in \mathbb{V}_\epsilon: \| \hat{F} s -y^\delta\|^2 = \delta_\epsilon^2\} = \{s\in \mathbb{V}_\epsilon: (s'\hat{F}' - (y^{\delta})')(\hat{F} s - y^\delta) = \delta_\epsilon^2\},\] is the critical point of, \[ \mathcal{L}(s,\lambda) \equiv \lambda\left[ (s'\hat{F}' - (y^{\delta})')(\hat{F} s - y^\delta) -\delta_\epsilon^2\right] + s'C^{-1}_{\hat{p}} s.\] 
Setting to zero all of the partial derivatives of $\mathcal{L}$ with respect to the components of $s$ leads to the solution being given as the right-hand-side of (\ref{12/19/20.1}), where the factor $\epsilon^2$ is the reciprocal of the Lagrange multiplier.  The latter is determined by the final equation $\frac{\partial \mathcal{L}}{\partial \lambda} = 0$, which places the solution on $\partial\Omega$.  This solution is indeed $\hat{p}$ since we are {\it given} (\ref{12/19/20.1}) (it is implied by (\ref{12/14/20.1})) and we know that $\hat{p}$ lies on $\partial\Omega$. 

Now define the function  $h:\mathbb{V}_\epsilon\rightarrow \mathbb{R}$ with $h(s) \equiv s'C^{-1}_{\hat{p}} s$. According to the prior paragraph, $\hat{p}$ minimizes $h(z)$ on  $\partial\Omega$.   Consider the level set of $h$ given by the ellipsoid, \[ \mathcal{H} \equiv \{s\in \mathbb{V}_\epsilon: s'C_{\hat{p}}^{-1} s = n \} = \left\{s\in\mathbb{V}_\epsilon: \sum_{i=1}^n \frac{s_i^2}{p_i^2} = n\right\},\] where $s_j\equiv v_j's$ and  $p_j\equiv v_j'p$ for $v_j\in \{v_i\}_{i=1}^n$.  
The point $\hat{p}$ is evidently a member of $\mathcal{H}$.  The tangent space of ellipsoid $\mathcal{H}$ at $\hat{p}$  has normal vectors proportional to the gradient of $h$ at $\hat{p}$, and with respect to the $\{v_i\}_{i=1}^n$-basis,
\begin{equation}\label{3/18/12.1}  \nabla h(s)\Big|_{s=\hat{p}} = \,\, 2\left(\frac{1}{p_1},\dots, \frac{1}{p_n}\right).\end{equation} But as we have noted above, $\hat{p}$ is the solution estimate minimizing $s'C^{-1}_{\hat{p}} s$ on $\partial\Omega$.  Thus, $\mathcal{H}$ is tangent to $\partial\Omega$ at $\hat{p}$.  
 
We next introduce  $g:\mathbb{V}_\epsilon\rightarrow \mathbb{R}$ with $g(s) = \prod_{i=1}^n |s_i|$ (the $n$-th power of the geometric mean of the absolute values of a point's components with respect to the $\{v_i\}_{i=1}^n$-basis).  The level set of $g$ that contains $\hat{p}$ is,  \[\mathcal{G} \equiv \left\{ s\in \mathbb{V}_\epsilon:\prod_{i=1}^n |s_i| = g(\hat{p})\right\}.\] We observe that at any point of this level set,
   \begin{equation}\label{9/18/22.1} \frac{\partial g}{\partial s_j} = {\rm sgn}(s_j) \prod_{i\ne j} |s_i| = \frac{g(\hat{p})}{s_j}.\end{equation} The tangent space of $\mathcal{G}$ at $\hat{p}$ has normal vectors proportional to the gradient of $g(s)$ at $\hat{p}$ and, with respect to the $\{v_i\}_{i=1}^n$-basis, (\ref{9/18/22.1}) implies \begin{equation}\label{9/15/20.1} \nabla g(s)\Big|_{s=\hat{p}} = \,\, g(\hat{p}) \left(\frac{1}{p_1},\dots, \frac{1}{p_n}\right),\end{equation} which according to (\ref{3/18/12.1}) is proportional to $\nabla h(s)\big|_{s=\hat{p}}$.
Hence, $\mathcal{H}$ and $\mathcal{G}$ both contain $\hat{p}$, and share the same tangent space at $\hat{p}$, i.e., $\mathcal{G}$ and $\mathcal{H}$ are tangent at $\hat{p}$.  But we have seen that $\mathcal{H}$ is tangent to $\partial\Omega$ at $\hat{p}$.  Hence $\mathcal{G}$ is tangent to $\partial\Omega$ at $\hat{p}$, and $\hat{p}$ is thereby a critical point of $g$ on $\partial\Omega$.  The $n$-th root of $g$, which is the geometric mean, then also has a critical point on $\partial\Omega$ at $\hat{p}$, which is equivalent to the final statement in the theorem. 
 \end{proof} 
 
 By the way, the deterministic method is convergent when extended to the infinite dimensional setting, as shown in \cite{fgreensite:2022d} (that reference also provides an illustrative numerical example comparing it with standard zero-order Tikhonov regularization). 
 
 \bigskip
 
Theorem \ref{7/6/21.6} indicates that the regularized solution estimate is a critical point of the geometric mean on a particular discrepancy manifold.  With zero-order Tikhonov regularization, the regularized solution estimate is a critical point of the Euclidean norm on a discrepancy manifold (actually, it is the only relative minimum of the Euclidean norm on the discrepancy manifold) - and in generalized Tikhonov regularization, the minimization is at least performed with seminorms.  So we are already motivated to think of the geometric mean as some kind of a ``norm", albeit one not satisfying the triangle inequality. 

The geometric mean pertains to this regularization method precisely because of (\ref{9/18/22.1}), (\ref{9/15/20.1}), i.e., $\nabla g(s) = s^{-1}g(s)$ where $g(s)$ is the geometric mean raised to the $n$-th power (i.e., $g(s) = \prod_i |s_i|$), , and \begin{equation}\label{12/5/23.1} s^{-1}\equiv\left(\frac{1}{s_1},\dots, \frac{1}{s_n}\right).\end{equation} 
The geometric mean is a degree-1 positive homogeneous function extracted from $g(s) $, but we will modify it by multiplying it by a constant for reasons that will become clear shortly.   Thus, for a point $s = (s_1, \dots, s_n)\in \mathbb{R}^n$, with $s_i\ne 0$ for $i=1,\dots,n$, we define  \begin{equation}\label{12/6/23.3} \ell(s)\equiv\sqrt{n}\left(\prod_i |s_i|\right)^{1/n}.\end{equation}   Combining this with (\ref{12/5/23.1}), we have \begin{equation}\label{12/6/23.1}\ell(s)\ell(s^{-1}) = n = \mathbf{1}\cdot\mathbf{1}\equiv \|\mathbf{1}\|^2,\end{equation} where $\mathbf{1}\equiv (1,1,\dots,1)$.  Thus, for $s$ in some neighborhood of $\mathbf{1}$,
\[ \nabla \ell(s) = \frac{\ell(s)}{n}\left(\frac{1}{s_1},\dots,\frac{1}{s_n}\right) = \frac{s^{-1}}{\ell(s^{-1})},  \] i.e.,  \begin{equation}\label{12/5/23.2} s^{-1} = \ell(s^{-1})\nabla\ell(s),\end{equation}  Applying $\ell$ to both sides of the above equation, and using the degree-1 positive homogeneity of $\ell$,  we obtain \begin{equation}\label{12/6/23.5} \ell(\nabla\ell(s))
= 1.\end{equation}


 
 Equations (\ref{12/5/23.2}), (\ref{12/6/23.5}), describing the geometric mean, look a lot like the equations describing the Euclidean norm,  \enlargethispage{1cm}\begin{equation}\tag{\ref{11/29/23.5}}  s = \ell(s)\nabla\ell(s),\end{equation} and \begin{equation}\tag{\ref{11/29/23.6}} \ell(\nabla\ell(s)) = 1.\end{equation}  That is, the nonnegative solution for $\ell(s)$ from the latter two equations is the Euclidean norm, and the two equations express a point as the product of its Euclidean norm with an associated unit-norm direction.
So from this standpoint also, it seems that the geometric mean {\it is} indeed rather ``norm-like".  In particular, for the case of $\ell$ deriving from the geometric mean as in (\ref{12/6/23.3}) (which, in the above regularization context, we have already come to think of as some kind of a ``norm"), according to (\ref{12/5/23.2}), (\ref{12/6/23.5}) a point in a neighborhood of $\mathbf{1}$ has the decomposition,
\begin{equation}\label{8/4/20.2} s = \ell(s)\nabla \ell(s^{-1}) \mbox{  with   }  \ell(\nabla \ell(s^{-1})) = 1,\end{equation} where the expression $\nabla \ell(s^{-1})$ means that the gradient is obtained following which it is evaluated at $s^{-1}$.  Like the pair of equations (\ref{11/29/23.5}), (\ref{11/29/23.6}) for the Euclidean norm, the equations (\ref{8/4/20.2}) can also be interpreted as expressing a point as the product of its ``norm" with an associated unit-``norm"  direction.


Looking at (\ref{12/6/23.1}), (\ref{12/5/23.2}),  (\ref{12/6/23.5}) in isolation, we have to wonder: what is the significance of the ``inverse of a point", $s^{-1}$?  Well, it suggests that $s$ is an element of a unital algebra.  Which algebra?  In the case at hand, the inverse of a point happens to be given by (\ref{12/5/23.1}).   Equation (\ref{12/5/23.1}) indicates that the multiplicative identity element is ${\bf 1}= (1,1,\dots, 1)$, since this element will be equal to its inverse.  But we must have  $ss^{-1} = (s_1,\dots,s_n)(1/s_1,\dots,1/s_n) = {\bf 1}$, where element juxtaposition denotes the product. The equations of the last two sentences indicate that we must then be dealing with the algebra (whose vector space of elements is $\mathbb{R}^n$) defined by component-wise addition and multiplication, i.e., $\bigoplus_{i=1}^n \mathbb{R}$.  And that gets one thinking: how many other algebras (i.e., real finite-dimensional unital associative algebras) are associated with an entity $\ell(s)$ satisfying something like (\ref{12/6/23.1}), (\ref{12/5/23.2}),  (\ref{12/6/23.5}), i.e., satisfying ``norm-like" equations respecting the multiplicative inversion operation of the algebra?
 

The answer is, all of them.  But to recognize this, we must first deal with the observation that given (\ref{12/6/23.1}) (which mandates respect for the multiplicative inversion operation), the right-hand-side of (\ref{12/5/23.2}) is a gradient, i.e., $\nabla\log \ell(s)$ - but the left-hand-side, $s^{-1}$, is not a gradient for most unital algebras.  
To rectify that, we would have to replace $s^{-1}$ on the left-hand-side of (\ref{12/5/23.2}) with $Ls^{-1}$, where $L$ has the feature that $Ls^{-1}$ {\it is} a gradient (and by the way, the left-hand-side of (\ref{12/6/23.5}) would then become $\ell(L^{-1}\nabla\ell(s))$ leading to an equation comparable to  (\ref{11/29/23.2}), assuming $L$ to be nonsingular).   


The requirement that $Ls^{-1}$ be a gradient means that $L$ must satisfy, \[ d([Ls^{-1}]\cdot {\sf d}s) = 0.\]  We are thus led to Definition \ref{7/21/23.1}, the beginning of our program (and, of course, the equations in (\ref{8/4/20.2}) are special cases of (\ref{5/22/24.5}) and (\ref{5/22/24.6})).

\vspace{.5cm}

\section{Simple observations on sorting versus classifying} 

If one is able to classify, then there is a complete set of invariants, and two objects are isomorphic if and only their invariants are equal.  If classification is impossible, one must settle for sorting into ``bins" rather than isomorphism classes.  Given some incomplete set of invariants with respect to which one sorts, two objects are in the same bin if and only if the latter limited sets of invariants (which we could call ``descriptions" instead of classifications) are the same for the two objects.  

Some sortings/descriptions are better than others, or perhaps more interesting.  One might even assert that a classification is not necessarily ``better" than a sorting, since from one perspective they might have rather distinct objectives.  Consider the set of complex $n\times n$-matrices.  Jordan normal form embodies a complete set of \mbox{invariants} (with respect to similarity), so this canonical form thereby fully classifies.  Instead, suppose we sort via the value of the determinant of the matrix (a single similarity invariant).  One might consider this a very poor system.  And yet, the bin which is filled with the matrices of determinant 1 is quite interesting, starting with its contents being a normal subgroup of the general linear group - and this entity is not especially evident from the invariants comprising the Jordan normal form (similar observations can be made regarding other matrix Lie groups). Furthermore, the fact that homeomorphism-equivalence classification of manifolds above dimension three is impossible by no means detracts from the beauty and usefulness of cohomology.  


Indeed, the role of sorting rather than classifying is implied most anytime a useful functor makes an appearance.

\vspace{.75cm}

\section{Some implied research topics} 
\begin{enumerate} \item Investigate the combinatorial problems related to the invariants of the uncurling metric group (Corollary \ref{8/18/23.2}) as proposed in Section 2.5, and similar combinatorial problems relating to the invariants of the unital norm group (Corollary \ref{7/25/24.1}). \item Regarding the functor in Section 2.9 relating to analysis of algebra epimorphisms, which objects in the particular category of symmetric parametrized matrics are anti-rotors? \item Settle the Conjecture stated in Section 5.2. \item Extend this type of approach to complex algebras. \item Extend this type of approach to infinite-dimensional algebras.  \end{enumerate}
  
 \vspace{.75cm}
  
   \begin{thebibliography}{99}
   
\bibitem{eom:2024} Encyclopedia of Mathematics (accessed June 13, 2024) Representation of an associative algebra. URL:http://encyclopediaofmath.org.
   
    \bibitem{drozd:1972} Drozd YA (1972) Representations of commutative algebras. Functional Analysis and its Applications, 6:286-288.
          
       \bibitem{belitskii:2005} Belitskii G, Lipyanski R, Sergeichuk VV (2005) Problems of classifying associative or Lie algebras and triples of symmetric or skew-symmetric matrices are wild.  Linear Algebra and its Applications, 407:249-262.
       
              \bibitem{bourbaki:1989} Bourbaki N (1989) Algebra I, Springer, Berlin, p. 543.

               
       \bibitem{bannwarth:2015} Bannwarth I, El Din M S (2015) Probabilistic algorithm for computing the dimension of real algebraic sets.  ISSAC'15 - 40th International Symposium on Symbolic and Algebraic Computation, July 2015, Bath UK, pp. 37-44, 10.1145/2755996.2756670, hal-01152751.
       
              \bibitem{kobayashi:2021} Kobayashi Y, Shirayanagi K, Tsukada M, Takahasi S (2021) A complete classification of three-dimensional algebras over $\mathbb{R}$ and $\mathbb{C}$. Asian-European Journal of Mathematics, 14(8) 2150131.                                        
  
              \bibitem{jacobson:1963} Jacobson N  (1963) Generic norm of an algebra. Osaka mathematics Journal, 15(1):25-50.
  
        \bibitem{schafer:1963}  Schafer R (1963) On forms of degree $n$ permitting composition. Journal of Mathematics and Mechanics, 12:777-792. 
         
      
       \bibitem{mccrimmon:2004} McCrimmon K (2004) A taste of Jordan algebras, Springer, New York.
       

              \bibitem{fgreensite:2022b} Greensite F (2022), A new proof of the Pythagorean Theorem and generalization of the usual algebra norm. arXiv:2209.14119
              
                    \bibitem{fgreensite:2022d} Greensite F (2022), Origin of the uncurling metric formalism from an Inverse Problems regularization method. arXiv:2209.14137

         \bibitem{birkhoff:1932} Birkhoff, GD (1932), A Set of Postulates for Plane Geometry (Based on Scale and Protractors), Annals of Mathematics, 33: 329–345.
         
               \bibitem{vogel:2002} Vogel C (2002) Computational Methods for Inverse Problems. Philadelphia: SIAM.
      
      \bibitem{engl:2000}  Engl HW,  Hanke M,  Neubauer A (2000) Regularization of Inverse Problems, Klewer, Boston.
      
      \bibitem{hansen:1992} Hansen PC (1992) Analysis of discrete ill-posed problems by means of the L-curve. SIAM Review, 34:563-580.
      
       
      \end{thebibliography}
  
 
   \end{document}  
   
 